% Traduction et production de textes divers liés à la citoyenneté — Allemand Tle A — Cours signé OnBuch+
% Programme officiel MINESEC. Compiler avec : tectonic AL30-traduction-et-production-de-textes-divers-lies-a-la-citoyenn.tex
\newcommand{\DOCMATIERE}{Allemand}
\newcommand{\DOCNIVEAU}{Tle A}
\newcommand{\DOCMODULE}{Chapitre 5 : Citoyenneté}
\newcommand{\DOCLECON}{AL30}
\newcommand{\DOCTITRE}{Traduction et production de textes divers liés à la citoyenneté}
% OnBuch+ — préambule « cours signés » ENRICHI (compilé avec tectonic / XeLaTeX)
% Système visuel « Pop » d'OnBuch+ : fond crème (jamais blanc), bordures encre
% épaisses, ombres « dures » décalées, titres Archivo Black en capitales.
% Version MATHÉMATIQUES : graphes (pgfplots/tikz), boîtes pédagogiques
% (définition, propriété, méthode, attention, le savais-tu, exercice résolu,
% exercice, TP, bilan), page de garde et signature OnBuch+ ancrée en bas.
%
% Macros attendues dans main.tex AVANT \input{preamble} :
%   \DOCMATIERE  \DOCNIVEAU  \DOCMODULE  \DOCLECON (numéro)  \DOCTITRE
\documentclass[11pt]{article}

\usepackage{fontspec}
\usepackage[ngerman,french]{babel} % français = langue principale ; l'allemand via \all{..} / textebox
\usepackage[a4paper,margin=2cm,top=3.4cm,bottom=2.8cm,headheight=1.4cm,footskip=1.3cm]{geometry}
\usepackage{amsmath,amssymb,amsfonts}
\usepackage{mathtools} % \xmapsto, \coloneqq... souvent utilisés par les modèles
\usepackage{cancel} % simplifications barrées (bilans de forces)
% \abs{z} (module, valeur absolue) et \norm{u} (norme), utilisés par les modèles
\providecommand{\abs}{}\let\abs\relax\DeclarePairedDelimiter{\abs}{\lvert}{\rvert}
\providecommand{\norm}{}\let\norm\relax\DeclarePairedDelimiter{\norm}{\lVert}{\rVert}
\usepackage[table]{xcolor} % [table] : fournit aussi \rowcolors (alternance de lignes)
\usepackage{pagecolor}
\usepackage{tikz}
\usetikzlibrary{arrows.meta,positioning,calc,shapes.geometric,shapes.misc,shapes.arrows,decorations.pathmorphing,decorations.pathreplacing,decorations.markings,patterns,shadows,mindmap,trees,fit,backgrounds,matrix,intersections,angles,quotes}
\usepackage{pgfplots}
\usepackage[version=4]{mhchem} % équations nucléaires \ce{^{235}_{92}U}
\usepackage[european,siunitx]{circuitikz} % schémas électriques
\pgfplotsset{compat=1.18}
\usepgfplotslibrary{fillbetween,groupplots} % aires entre courbes (calcul intégral)
\usepackage{enumitem}
\usepackage{fancyhdr}
% [most] charge toutes les bibliothèques tcolorbox (listings, minted, raster,
% documentation...) et gonfle inutilement la mémoire TeX sur un document long
% (~300 boîtes) : on ne charge que ce qui est réellement utilisé.
\usepackage[skins,breakable]{tcolorbox}
\usepackage{graphicx}
\usepackage{colortbl}
\usepackage{booktabs}
\usepackage{tabularx}
\usepackage{diagbox} % case d'en-tête diagonale des tableaux à double entrée
% Colonnes extensibles centrée / gauche / droite (C, L, R), souvent utilisées par les modèles
\newcolumntype{C}{>{\centering\arraybackslash}X}
\newcolumntype{L}{>{\raggedright\arraybackslash}X}
\newcolumntype{R}{>{\raggedleft\arraybackslash}X}
\usepackage{array}
\usepackage{multicol}
\usepackage{multirow}
\usepackage{siunitx}
% parse-numbers=false : accepte aussi \SI{2,5 \times 10^{-3}}{...}
\sisetup{locale=FR,output-decimal-marker={,},group-minimum-digits=4,parse-numbers=false}
\DeclareSIUnit\ohm{\ensuremath{\symup{\Omega}}} % U+2126 absent de la police
\DeclareSIUnit\division{div} % réglages d'oscilloscope (ms/div, V/div)
\DeclareSIUnit\tr{tr} % tours (vitesse de rotation, ex. \SI{1500}{\tr\per\minute})
\DeclareSIUnit\molar{M} % concentration molaire (ex. \SI{3}{\molar}, \SI{1}{\micro\molar})
\DeclareSIUnit\an{an} % année (le modèle confond parfois avec \year, un compteur TeX) — ex. \SI{5}{\kilo\meter\per\an}
\DeclareSIUnit\FCFA{FCFA} % franc CFA (ex. \SI{500}{\FCFA\per\kilo\gram})
\DeclareSIUnit\degreeBrix{\textdegree Brix} % teneur en sucre d'un sirop/jus (réfractométrie)
\DeclareSIUnit\month{mois} % le modèle utilise parfois \month, un compteur TeX, comme unité de durée
\usepackage{qrcode}
% \degree : commande fréquemment utilisée par les modèles pour un angle ou une
% température en dehors de \SI{}{} (non fournie par les paquets chargés).
\providecommand{\degree}{^{\circ}}
% \qed (fin de démonstration), utilisé par les modèles sans amsthm
\providecommand{\qed}{\hfill$\square$}
% \barre : double barre des valeurs interdites dans un tableau de variations (tkz-tab)
\providecommand{\barre}{\ensuremath{\|}}
% environnement proof (amsthm non chargé) utilisé par les modèles
\ifdefined\proof\else\newenvironment{proof}[1][Démonstration]{\par\noindent\textit{#1.}\ }{\qed\par}\fi
\usepackage{lastpage}
\usepackage{hyperref}
\hypersetup{hidelinks}

% ============================================================
% Palette « Pop » OnBuch+ — mêmes hex que la classe P (plus_ui.dart)
% ============================================================
\definecolor{poporange}{HTML}{F59321}
\definecolor{popdark}{HTML}{E07A0C}
\definecolor{popcream}{HTML}{FFF6E8}
\definecolor{popink}{HTML}{1C1714}
\definecolor{poppurple}{HTML}{7A5AE0}
\definecolor{popgreen}{HTML}{2E9E6A}
\definecolor{poppink}{HTML}{E0457B}
\definecolor{popblue}{HTML}{2F7DE1}
\definecolor{popgold}{HTML}{C98A12}
\definecolor{popmuted}{HTML}{8A7F76}
\definecolor{popline}{HTML}{EADFD2}
\definecolor{popgray}{HTML}{8A7F76}
\colorlet{popgrey}{popgray}
% Alias tolérants (le modèle nomme parfois les couleurs différemment)
\colorlet{popwhite}{white}
\colorlet{popblack}{popink}
\colorlet{popred}{poppink}
\colorlet{popyellow}{popgold}
% Teintes claires (fonds de tableaux/encadrés) — le modèle nomme parfois
% les couleurs avec un suffixe L (ex. popblueL) sans les avoir définies.
\colorlet{poporangeL}{poporange!20}
\colorlet{popdarkL}{popdark!20}
\colorlet{poppurpleL}{poppurple!20}
\colorlet{popgreenL}{popgreen!20}
\colorlet{poppinkL}{poppink!20}
\colorlet{popblueL}{popblue!20}
\colorlet{popgoldL}{popgold!20}
\colorlet{popredL}{popred!20}
\colorlet{popgrayL}{popgray!20}
% Teintes claires pour fonds de tableaux / zones
\colorlet{popblueL}{popblue!10!white}
\colorlet{poporangeL}{poporange!14!white}
\colorlet{popgreenL}{popgreen!12!white}
\colorlet{poppurpleL}{poppurple!12!white}
\colorlet{poppinkL}{poppink!12!white}
\colorlet{popgoldL}{popgold!14!white}

\pagecolor{popcream}

% ============================================================
% Typographie
% ============================================================
\defaultfontfeatures{Ligatures=TeX}
\setmainfont{PlusJakartaSans-Medium.ttf}[Path=../../../fonts/,BoldFont=PlusJakartaSans-Bold.ttf]
% Maths en sans-serif (Fira Math) avec chiffres et lettres droites de Plus
% Jakarta, pour que les formules se fondent dans le texte.
\usepackage[math-style=ISO,bold-style=ISO]{unicode-math}
\setmathfont{FiraMath-Regular.otf}[Path=../../../fonts/]
\setmathfont{PlusJakartaSans-Medium.ttf}[Path=../../../fonts/,range={up/num,up/{Latin,latin},\mathup}]
\usepackage{icomma} % virgule décimale sans espace en mode math
% Caractères mathématiques Unicode absents des polices de texte (Plus Jakarta,
% Archivo Black) : rendus par leur équivalent mathématique, en texte comme en
% formule — sinon ils s'affichent en carré vide (ex. « Arithmétique dans ℤ »).
\usepackage{newunicodechar}
\newunicodechar{ℝ}{\ensuremath{\mathbb{R}}}
\newunicodechar{ℤ}{\ensuremath{\mathbb{Z}}}
\newunicodechar{ℂ}{\ensuremath{\mathbb{C}}}
\newunicodechar{ℕ}{\ensuremath{\mathbb{N}}}
\newunicodechar{ℚ}{\ensuremath{\mathbb{Q}}}
\newunicodechar{𝔻}{\ensuremath{\mathbb{D}}}
\newunicodechar{α}{\ensuremath{\alpha}}
\newunicodechar{β}{\ensuremath{\beta}}
\newunicodechar{γ}{\ensuremath{\gamma}}
\newunicodechar{δ}{\ensuremath{\delta}}
\newunicodechar{ε}{\ensuremath{\varepsilon}}
\newunicodechar{θ}{\ensuremath{\theta}}
\newunicodechar{λ}{\ensuremath{\lambda}}
\newunicodechar{μ}{\ensuremath{\mu}}
\newunicodechar{σ}{\ensuremath{\sigma}}
\newunicodechar{τ}{\ensuremath{\tau}}
\newunicodechar{φ}{\ensuremath{\varphi}}
\newunicodechar{ψ}{\ensuremath{\psi}}
\newunicodechar{ω}{\ensuremath{\omega}}
\newunicodechar{Σ}{\ensuremath{\Sigma}}
% majuscules produites par \MakeUppercase dans les titres (α-aminés -> Α)
\newunicodechar{Α}{\ensuremath{\alpha}}
\newunicodechar{Β}{\ensuremath{\beta}}
\newunicodechar{Γ}{\ensuremath{\Gamma}}
\newunicodechar{Δ}{\ensuremath{\Delta}}
\newunicodechar{Ε}{\ensuremath{\varepsilon}}
\newunicodechar{Θ}{\ensuremath{\Theta}}
\newunicodechar{Λ}{\ensuremath{\Lambda}}
\newunicodechar{Μ}{\ensuremath{\mu}}
\newunicodechar{Τ}{\ensuremath{\tau}}
\newunicodechar{Φ}{\ensuremath{\Phi}}
\newunicodechar{Ψ}{\ensuremath{\Psi}}
\newunicodechar{Ω}{\ensuremath{\Omega}}
\newunicodechar{↦}{\ensuremath{\mapsto}}
\newunicodechar{∈}{\ensuremath{\in}}
\newunicodechar{∉}{\ensuremath{\notin}}
\newunicodechar{∧}{\ensuremath{\wedge}}
\newunicodechar{⇔}{\ensuremath{\Leftrightarrow}}
\newunicodechar{⟺}{\ensuremath{\Longleftrightarrow}}
\newunicodechar{⇒}{\ensuremath{\Rightarrow}}
\newunicodechar{⟹}{\ensuremath{\Longrightarrow}}
\newunicodechar{∘}{\ensuremath{\circ}}
\newunicodechar{∪}{\ensuremath{\cup}}
\newunicodechar{∩}{\ensuremath{\cap}}
\newunicodechar{≡}{\ensuremath{\equiv}}
\newunicodechar{⊂}{\ensuremath{\subset}}
\newunicodechar{⊥}{\ensuremath{\perp}}
\newunicodechar{∃}{\ensuremath{\exists}}
\newunicodechar{∀}{\ensuremath{\forall}}
\newunicodechar{∛}{\ensuremath{\sqrt[3]{\,}}}
\newunicodechar{∜}{\ensuremath{\sqrt[4]{\,}}}
\newunicodechar{⃗}{\ensuremath{{}^{\to}}}
\newunicodechar{ⁿ}{\ifmmode^{n}\else\textsuperscript{\ensuremath{n}}\fi}
\newunicodechar{ˣ}{\ifmmode^{x}\else\textsuperscript{\ensuremath{x}}\fi}
\newunicodechar{ᵇ}{\ifmmode^{b}\else\textsuperscript{\ensuremath{b}}\fi}
\newunicodechar{ʳ}{\ifmmode^{r}\else\textsuperscript{\ensuremath{r}}\fi}
\newunicodechar{ᵘ}{\ifmmode^{u}\else\textsuperscript{\ensuremath{u}}\fi}
\newunicodechar{ʸ}{\ifmmode^{y}\else\textsuperscript{\ensuremath{y}}\fi}
\newunicodechar{ᵃ}{\ifmmode^{a}\else\textsuperscript{\ensuremath{a}}\fi}
\newunicodechar{ᵉ}{\ifmmode^{e}\else\textsuperscript{\ensuremath{e}}\fi}
\newunicodechar{ᵏ}{\ifmmode^{k}\else\textsuperscript{\ensuremath{k}}\fi}
\newunicodechar{ᵖ}{\ifmmode^{p}\else\textsuperscript{\ensuremath{p}}\fi}
\newunicodechar{ᴺ}{\ifmmode^{N}\else\textsuperscript{\ensuremath{N}}\fi}
\newunicodechar{ᶜ}{\ifmmode^{c}\else\textsuperscript{\ensuremath{c}}\fi}
\newunicodechar{ʰ}{\ifmmode^{h}\else\textsuperscript{\ensuremath{h}}\fi}
\newunicodechar{ᵠ}{\ifmmode^{q}\else\textsuperscript{\ensuremath{q}}\fi}
\newunicodechar{ₙ}{\ifmmode_{n}\else\textsubscript{\ensuremath{n}}\fi}
\newunicodechar{ₐ}{\ifmmode_{a}\else\textsubscript{\ensuremath{a}}\fi}
\newunicodechar{ᵢ}{\ifmmode_{i}\else\textsubscript{\ensuremath{i}}\fi}
\newunicodechar{ᵣ}{\ifmmode_{r}\else\textsubscript{\ensuremath{r}}\fi}
\newunicodechar{ₖ}{\ifmmode_{k}\else\textsubscript{\ensuremath{k}}\fi}
\newunicodechar{ᵦ}{\ifmmode_{\beta}\else\textsubscript{\ensuremath{\beta}}\fi}
\newunicodechar{ₓ}{\ifmmode_{x}\else\textsubscript{\ensuremath{x}}\fi}
\newfontfamily\popdisplay{ArchivoBlack-Regular.ttf}[Path=../../../fonts/]
\newfontfamily\poplogo{SpaceGrotesk-Bold.ttf}[Path=../../../fonts/]
\newfontfamily\popsemi{PlusJakartaSans-SemiBold.ttf}[Path=../../../fonts/]
\newfontfamily\popxbold{PlusJakartaSans-ExtraBold.ttf}[Path=../../../fonts/]
\newfontfamily\popxboldls{PlusJakartaSans-ExtraBold.ttf}[Path=../../../fonts/,LetterSpace=3.0]

\setlist[enumerate]{leftmargin=1.7em,itemsep=5pt,topsep=4pt}
\setlist[itemize]{leftmargin=1.7em,itemsep=4pt,topsep=4pt,label={\color{poporange}\textbullet}}
\setlength{\parindent}{0pt}
\setlength{\parskip}{5pt}
\renewcommand{\arraystretch}{1.25}

% pgfplots : style Pop par défaut
\pgfplotsset{
  every axis/.append style={axis line style={popink,line width=1pt},tick style={popink},
    grid style={popline,line width=0.5pt},label style={font=\small},tick label style={font=\footnotesize}},
  popcurve/.style={poporange,line width=1.8pt,no markers,smooth},
}
% Même style disponible dans un \draw TikZ simple (hors pgfplots)
\tikzset{popcurve/.style={poporange,line width=1.8pt}}
\tikzset{popgrid/.style={popline,very thin}}
\colorlet{popcurve}{poporange} % le nom du style est parfois utilisé comme couleur

% ============================================================
% Pill / squiggle
% ============================================================
\newcommand{\poppill}[3]{%
  \tikz[baseline=-0.55ex]\node[draw=popink,line width=1.2pt,rounded corners=2.6pt,%
    fill=#2,inner xsep=6.5pt,inner ysep=3pt]{%
    \popxboldls\fontsize{7}{7}\selectfont\color{#3}\MakeUppercase{#1}};%
}
\newcommand{\popsquiggle}[1]{%
  \tikz[baseline=-0.4ex]{
    \draw[#1,line width=2.6pt,line cap=round]
      plot [smooth] coordinates {(0,0.09) (0.34,0.01) (0.68,0.21) (1.02,0.01) (1.36,0.10)};
  }%
}
% Mot-clé surligné (marqueur orange sous le texte)
\newcommand{\cle}[1]{{\popxbold\color{popdark}#1}}

% ============================================================
% En-tête / pied
% ============================================================
\pagestyle{fancy}
\fancyhf{}
\renewcommand{\headrulewidth}{0pt}
\renewcommand{\footrulewidth}{0pt}
\lhead{\raisebox{-0.15ex}{\poplogo\fontsize{13}{13}\selectfont\color{popink}OnBuch\color{poporange}+}}
\rhead{\poppill{\DOCMATIERE\ \textbullet\ \DOCNIVEAU\ \textbullet\ Leçon \DOCLECON}{white}{popink}}
\lfoot{\popxbold\fontsize{7.6}{7.6}\selectfont\color{popmuted}\MakeUppercase{Cours signé OnBuch+}\ \textbullet\ {\popsemi\DOCTITRE}}
\rfoot{\tikz[baseline=-0.6ex]\node[rounded corners=8.5pt,fill=popink,minimum height=17pt,minimum width=17pt,inner xsep=5pt,inner sep=0]{\popxbold\fontsize{8}{8}\selectfont\color{popcream}\thepage\,/\,\pageref*{LastPage}};}

% ============================================================
% Titres
% ============================================================
\newcommand{\chaptitre}[2]{%
  \begin{tcolorbox}[enhanced,colback=poporange,colframe=popink,boxrule=2.6pt,arc=11pt,
      drop shadow=popink,left=15pt,right=15pt,top=12pt,bottom=11pt,before skip=3pt,after skip=18pt]
    \poppill{#2}{popink}{popcream}\par\vspace{9pt}
    {\raggedright\hyphenpenalty=10000\exhyphenpenalty=10000\popdisplay\fontsize{22}{24}\selectfont\color{popcream}\MakeUppercase{#1}\par}
  \end{tcolorbox}
}
% Section numérotée : pastille orange avec le numéro + titre Archivo
\newcounter{popsec}
\newcommand{\coursec}[1]{%
  \stepcounter{popsec}\setcounter{popsub}{0}%
  \par\vspace{16pt}\noindent
  \begin{minipage}{\linewidth}
  \tikz[baseline=-0.7ex]\node[draw=popink,line width=1.6pt,fill=poporange,rounded corners=4pt,minimum size=22pt,inner sep=0]{\popdisplay\fontsize{12}{12}\selectfont\color{popcream}\thepopsec};%
  \hspace{8pt}{\popdisplay\fontsize{15}{17}\selectfont\color{popink}\MakeUppercase{#1}}\par
  \vspace{4pt}\noindent{\color{poporange}\rule{36pt}{3.6pt}}
  \end{minipage}\par\nopagebreak\vspace{8pt}
}
\newcounter{popsub}
\newcommand{\courssub}[1]{%
  \stepcounter{popsub}%
  \par\vspace{9pt}\noindent{\popxbold\fontsize{11.5}{13}\selectfont\color{popink}{\color{poporange}\thepopsec.\thepopsub}\hspace{6pt}#1}\par\nopagebreak\vspace{4pt}
}

% ============================================================
% Boîtes pédagogiques — même recette : fond blanc, bordure épaisse colorée,
% ombre dure encre, badge plein. Titre optionnel [#1].
% ============================================================
\tcbset{popbase/.style={breakable,enhanced,colback=white,boxrule=2.3pt,arc=9pt,
  shadow={3pt}{-3pt}{0pt}{popink},left=14pt,right=14pt,top=11pt,bottom=11pt,before skip=11pt,after skip=15pt,
  fontupper=\popsemi\fontsize{10.6}{14.5}\selectfont\color{popink},
  fontlower=\popsemi\fontsize{10.6}{14.5}\selectfont\color{popink}}}
% Coche dessinée (le glyphe ✓ n'existe pas dans les polices du document)
\newcommand{\popcheck}[1]{\tikz[baseline=-0.5ex]\draw[#1,line width=1.6pt,line cap=round,line join=round](-0.13,0.0)--(-0.04,-0.09)--(0.13,0.1);}
\providecommand{\coche}{\popcheck{popgreen}}
\providecommand{\resultat}[1]{\par\smallskip\noindent\fcolorbox{popgreen}{popgreenL}{\parbox{\dimexpr\linewidth-2\fboxsep-2\fboxrule\relax}{\textbf{Résultat.} #1}}\par\smallskip}
\newcommand{\popboxhead}[3]{\poppill{#1}{#2}{white}\ifx\relax#3\relax\else\hspace{6pt}{\popxbold\fontsize{10.6}{12}\selectfont\color{popink}#3}\fi\par\vspace{7pt}}

\newtcolorbox{aretenir}[1][]{popbase,colframe=popgreen,before upper={\popboxhead{À retenir}{popgreen}{#1}}}
% Texte en allemand (document de lecture, dialogue, extrait) : cadre
% d'étude. Un titre optionnel (source, auteur) s'affiche dans l'en-tête.
\newtcolorbox{textebox}[1][]{popbase,colframe=popblue,before upper={\popboxhead{Text}{popblue}{#1}\begin{otherlanguage}{ngerman}},after upper={\end{otherlanguage}}}
% Marques ✓ / ✗ (forme correcte / fautive) souvent utilisées par les modèles
\providecommand{\cross}{\textcolor{poppink}{\ensuremath{\times}}}
\providecommand{\xmark}{\cross}\providecommand{\crossmark}{\cross}\providecommand{\cmark}{\ensuremath{\checkmark}}
% Ligne de réponse / justification à compléter (inventée par les modèles)
\providecommand{\justification}[1]{\par\noindent\textit{Justification~:} \dotfill\par}
% \textipa (package tipa non chargé) : la police ne couvre pas l'API -> simple passage
\providecommand{\textipa}[1]{#1}
% Soulignement (ulem non chargé) : \uline, \ul
\providecommand{\uline}[1]{\underline{#1}}\providecommand{\ul}[1]{\underline{#1}}
% \glossaire{\item ...} inventée par les modèles : liste de glossaire
\providecommand{\glossaire}[1]{\begin{itemize}#1\end{itemize}}
% \alert (beamer) inventée par les modèles : texte mis en évidence
\providecommand{\alert}[1]{\textbf{\textcolor{poppink}{#1}}}
% variantes ASCII d'ordinaux écrites en macro
\providecommand{\iere}{\textsuperscript{re}}\providecommand{\ieme}{\textsuperscript{e}}\providecommand{\ere}{\textsuperscript{re}}\providecommand{\eme}{\textsuperscript{e}}\providecommand{\er}{\textsuperscript{er}}\providecommand{\ie}{\textsuperscript{e}}
% Ordinaux français écrits en macro par les modèles : 1\ière, 3\ième
\newcommand{\ière}{\textsuperscript{re}}\newcommand{\ième}{\textsuperscript{e}}
% \exemple (inventée par les modèles) : étiquette « Exemple : »
\providecommand{\exemple}{\textbf{Exemple~:}\ }
% \square utilisable aussi hors mode math (cases à cocher)
\AtBeginDocument{\let\squareMath\square\renewcommand{\square}{\ensuremath{\squareMath}}}
% Mot ou phrase en allemand à l'intérieur d'un paragraphe français
\newcommand{\all}[1]{\ifmmode\text{\textit{#1}}\else\begingroup\selectlanguage{ngerman}\itshape #1\endgroup\fi}
\let\esp\all % alias toléré
\newtcolorbox{exemplebox}[1][]{popbase,colframe=poporange,before upper={\popboxhead{Exemple}{poporange}{#1}}}
\newtcolorbox{definition}[1][]{popbase,colframe=popblue,before upper={\popboxhead{Définition}{popblue}{#1}}}
\newtcolorbox{propriete}[1][]{popbase,colframe=poppurple,before upper={\popboxhead{Propriété}{poppurple}{#1}}}
\newtcolorbox{methode}[1][]{popbase,colframe=popgold,before upper={\popboxhead{Méthode}{popgold}{#1}}}
\newtcolorbox{attention}[1][]{popbase,colframe=poppink,before upper={\popboxhead{Attention}{poppink}{#1}}}
\newtcolorbox{experience}[1][]{popbase,colframe=popblue,colback=popblueL,before upper={\popboxhead{Expérience}{popblue}{#1}}}
% « Le savais-tu ? » : carte violette pleine (contexte, culture, Cameroun)
\newtcolorbox{savaistu}[1][]{popbase,colback=poppurple,colframe=popink,
  fontupper=\popsemi\fontsize{10.4}{14.2}\selectfont\color{white},
  before upper={\poppill{Le savais-tu\,?}{white}{poppurple}\ifx\relax#1\relax\else\hspace{6pt}{\popxbold\color{white}#1}\fi\par\vspace{7pt}}}
% Exercice résolu : énoncé en haut, \tcblower puis la solution sur fond crème
\newcounter{popexo}\newcounter{popexr}
\newtcolorbox{exoresolu}[1][]{popbase,colframe=poporange,colbacklower=popcream,
  segmentation style={popink,line width=1.2pt,dashed},
  before upper={\stepcounter{popexr}\popboxhead{Exercice résolu \thepopexr}{poporange}{#1}},
  before lower={\poppill{Solution}{popink}{popcream}\par\vspace{6pt}}}
% Exercice (non résolu en place) — la correction est dans la partie Corrigés
\newtcolorbox{exercice}[1][]{popbase,colframe=popink,
  before upper={\stepcounter{popexo}\popboxhead{Exercice \thepopexo}{popink}{#1}}}
% Corrigé
\newtcolorbox{corrige}[1][]{popbase,colframe=popgreen,colback=popgreenL,
  before upper={\popboxhead{Corrigé}{popgreen}{#1}}}
% Objectifs / prérequis / bilan
\newtcolorbox{objectifs}{popbase,colframe=popink,colback=poporangeL,
  before upper={\popboxhead{Objectifs de la leçon}{popink}{}}}
\newtcolorbox{prerequis}{popbase,colframe=popink,colback=popblueL,
  before upper={\popboxhead{Prérequis}{popblue}{}}}
\newtcolorbox{bilan}{popbase,colframe=popink,colback=white,boxrule=2.8pt,
  before upper={\popboxhead{Fiche bilan}{poporange}{L'essentiel à connaître pour l'examen}}}
% Figure encadrée avec légende
\newtcolorbox{popfigure}[1][]{enhanced,breakable=false,colback=white,colframe=popline,boxrule=1.4pt,arc=8pt,
  left=10pt,right=10pt,top=10pt,bottom=8pt,before skip=10pt,after skip=12pt,halign=center,
  lower separated=false,halign lower=center,
  fontlower=\popsemi\fontsize{9}{11.5}\selectfont\color{popmuted},
  #1}
\newcommand{\legende}[1]{\par\vspace{6pt}{\popsemi\fontsize{9}{11.5}\selectfont\color{popmuted}#1}}

% ============================================================
% Signature OnBuch+
% ============================================================
\newcommand{\poplogoinline}[1]{{\poplogo\fontsize{#1}{#1}\selectfont\color{popink}OnBuch\color{poporange}+}}
\newcommand{\signatureonbuch}{%
  \begin{tcolorbox}[enhanced,colback=popink,colframe=popink,boxrule=2.2pt,arc=10pt,
      drop shadow=poporange,left=14pt,right=14pt,top=10pt,bottom=10pt,before skip=0pt,after skip=0pt]
    \tikz[baseline=-0.7ex]\node[circle,fill=popgreen,draw=popcream,line width=1.3pt,minimum size=22pt,inner sep=0]{\popcheck{white}};%
    \hspace{9pt}\begin{minipage}[c]{0.62\linewidth}
      {\popxbold\fontsize{12}{13}\selectfont\color{popcream}Signé\ \textbullet\ {\poplogo OnBuch\color{poporange}+}}\par\vspace{2pt}
      {\raggedright\popsemi\fontsize{8}{10.5}\selectfont\color{popline}\MakeUppercase{Équipe pédagogique OnBuch+}\\\MakeUppercase{Conforme au programme officiel MINESEC}\par}
    \end{minipage}\hfill
    {\poplogo\fontsize{22}{22}\selectfont\color{popcream}OnBuch\color{poporange}+}
  \end{tcolorbox}%
}

% ============================================================
% Page de garde
% #1 titre · #2 matière · #3 niveau · #4 module · #5 n° leçon · #6 durée ·
% #7 description / objectifs (texte ou itemize)
% ============================================================
\newcommand{\pagegarde}[7]{%
  \thispagestyle{empty}
  \newgeometry{a4paper,margin=1.7cm,top=1.6cm,bottom=1.4cm}%
  \begin{tikzpicture}[remember picture,overlay]
    \fill[poporange!18] ([xshift=-3.2cm,yshift=-3.6cm]current page.north east) circle (4.6cm);
    \fill[poppurple!14] ([xshift=2.4cm,yshift=7.6cm]current page.south west) circle (3.8cm);
  \end{tikzpicture}%
  {\poplogo\fontsize{30}{30}\selectfont\color{popink}OnBuch\color{poporange}+}\hfill
  \raisebox{0.6ex}{\poppill{Cours signé \textbullet\ Premium}{popink}{popcream}}\par
  \vspace{1pt}\popsquiggle{poppurple}\par
  \vspace{8pt}
  \begin{tcolorbox}[enhanced,colback=poporange,colframe=popink,boxrule=2.8pt,arc=13pt,
      drop shadow=popink,left=18pt,right=18pt,top=12pt,bottom=13pt,after skip=10pt]
    \poppill{#2 \textbullet\ #3}{popink}{popcream}\hspace{5pt}\poppill{#4}{white}{popink}\par\vspace{8pt}
    {\popdisplay\fontsize{14}{14}\selectfont\color{popink}LEÇON #5}\par\vspace{4pt}
    {\raggedright\hyphenpenalty=10000\exhyphenpenalty=10000\popdisplay\fontsize{25}{27}\selectfont\color{popcream}\MakeUppercase{#1}\par}
    \vspace{8pt}
    {\popxbold\fontsize{9.5}{9.5}\selectfont\color{popink}\MakeUppercase{Durée conseillée : #6}}
  \end{tcolorbox}
  \begin{tcolorbox}[enhanced,colback=white,colframe=popink,boxrule=2.2pt,arc=10pt,
      drop shadow=popink,left=16pt,right=16pt,top=10pt,bottom=10pt,after skip=8pt]
    \poppill{Dans cette leçon}{poppurple}{white}\par\vspace{5pt}
    \popsemi\fontsize{9.3}{12.6}\selectfont\color{popink}#7
  \end{tcolorbox}
  \noindent\begin{minipage}[t]{0.487\linewidth}
    \begin{tcolorbox}[enhanced,colback=popgreen,colframe=popink,boxrule=2.2pt,arc=10pt,
        drop shadow=popink,left=11pt,right=11pt,top=10pt,bottom=10pt,height=3.1cm,valign=center]
      \tcbox[colback=white,colframe=popink,boxrule=1.3pt,arc=5pt,boxsep=0pt,nobeforeafter,
        left=3pt,right=3pt,top=3pt,bottom=3pt,valign=center]{\qrcode[height=1.9cm]{https://play.google.com/store/apps/details?id=com.luvvix.onbuch&hl=fr}}%
      \hspace{8pt}\begin{minipage}[c]{0.5\linewidth}
        {\popdisplay\fontsize{11}{12.5}\selectfont\color{white}\MakeUppercase{Télécharge OnBuch}}\par\vspace{4pt}
        {\raggedright\popsemi\fontsize{8.4}{11}\selectfont\color{white}Gratuit \textbullet\ Google Play\\Tuteur IA, annales, concours, résultats\par}
      \end{minipage}
    \end{tcolorbox}
  \end{minipage}\hfill
  \begin{minipage}[t]{0.487\linewidth}
    \begin{tcolorbox}[enhanced,colback=poppurple,colframe=popink,boxrule=2.2pt,arc=10pt,
        drop shadow=popink,left=11pt,right=11pt,top=10pt,bottom=10pt,height=3.1cm,valign=center]
      \tikz[baseline=-0.7ex]\node[circle,fill=white,draw=popink,line width=1.3pt,minimum size=1.6cm,inner sep=0]{\popdisplay\fontsize{24}{24}\selectfont\color{poppurple}+};%
      \hspace{8pt}\begin{minipage}[c]{0.6\linewidth}
        {\popdisplay\fontsize{11}{12.5}\selectfont\color{white}\MakeUppercase{Passe à OnBuch+}}\par\vspace{4pt}
        {\raggedright\popsemi\fontsize{8.4}{11}\selectfont\color{white}Tous les cours signés, Tuteur IA illimité, fiches PDF — onglet OnBuch+ de l'app\par}
      \end{minipage}
    \end{tcolorbox}
  \end{minipage}
  \vspace{8pt}
  \popcontenu
  \vfill
  \signatureonbuch
  \restoregeometry
  \newpage
}

% Rangée « Ce cours contient » (page de garde)
\newcommand{\poptile}[3]{% #1 couleur #2 gros texte #3 légende
  \begin{tcolorbox}[enhanced,colback=#1,colframe=popink,boxrule=2pt,arc=9pt,shadow={2.5pt}{-2.5pt}{0pt}{popink},
      left=6pt,right=6pt,top=9pt,bottom=9pt,halign=center,valign=center,height=1.75cm,width=\linewidth,nobeforeafter]
    {\popdisplay\fontsize{12.5}{13}\selectfont\color{white}\MakeUppercase{#2}}\par\vspace{3pt}
    {\centering\popsemi\fontsize{7.8}{9.5}\selectfont\color{white}#3\par}
  \end{tcolorbox}}
\newcommand{\popcontenu}{%
  \par\noindent{\popxboldls\fontsize{7.5}{7.5}\selectfont\color{popmuted}CE COURS SIGNÉ CONTIENT}\par\vspace{7pt}
  \noindent\begin{minipage}[t]{0.235\linewidth}\poptile{popblue}{Cours}{complet, illustré, pas à pas}\end{minipage}\hfill
  \begin{minipage}[t]{0.235\linewidth}\poptile{poporange}{Méthodes}{et exercices résolus}\end{minipage}\hfill
  \begin{minipage}[t]{0.235\linewidth}\poptile{poppink}{Exercices}{type examen + corrigés}\end{minipage}\hfill
  \begin{minipage}[t]{0.235\linewidth}\poptile{popgreen}{Fiche bilan}{l'essentiel à retenir}\end{minipage}\par}

% Sommaire
\newtcolorbox{sommaire}{enhanced,colback=white,colframe=popink,boxrule=2.2pt,arc=10pt,
  drop shadow=popink,left=18pt,right=18pt,top=13pt,bottom=13pt,before skip=6pt,after skip=20pt,
  before upper={\poppill{Sommaire}{poppurple}{white}\par\vspace{9pt}\popsemi\fontsize{10.6}{14.5}\selectfont\color{popink}}}

% Signature de fin de cours — poussée tout en bas de la dernière page
\newcommand{\findecours}{%
  \par\vfill\nopagebreak
  \begin{center}\popsquiggle{poporange}\end{center}\vspace{4pt}
  \signatureonbuch
}
\AtBeginDocument{\def\llbracket{\mathopen{[\mkern-3.5mu[}}\def\rrbracket{\mathclose{]\mkern-3.5mu]}}} % intervalles d'entiers (glyphe ⟦ absent de la police)
% Glyphes absents de Fira Math : redéfinis à partir de glyphes disponibles
\AtBeginDocument{%
  \def\setminus{\mathbin{\backslash}}\let\smallsetminus\setminus
  \def\vdots{\mathord{\vbox{\baselineskip3.2pt\lineskiplimit0pt\kern4pt\hbox{.}\hbox{.}\hbox{.}}}}%
  \def\ddots{\mathinner{\mkern1mu\raise6.5pt\hbox{.}\mkern2mu\raise3.6pt\hbox{.}\mkern2mu\raise0.7pt\hbox{.}\mkern1mu}}%
  \def\triangle{\mathord{\scalebox{0.9}{$\symup{\Delta}$}}}%
  \def\nabla{\mathord{\raisebox{\depth}{\scalebox{1}[-1]{$\symup{\Delta}$}}}}%
  \def\checkmark{\popcheck{popdark}}%
  \def\star{\mathbin{\ast}}\def\bigstar{\mathord{\ast}}%
  \def\diamond{\mathbin{\scalebox{0.6}{$\Diamond$}}}%
  \def\bigcup{\mathop{\mathchoice{\raisebox{-0.3em}{\scalebox{1.7}{$\cup$}}}{\scalebox{1.25}{$\cup$}}{\cup}{\cup}}\displaylimits}%
  \def\bigcap{\mathop{\mathchoice{\raisebox{-0.3em}{\scalebox{1.7}{$\cap$}}}{\scalebox{1.25}{$\cap$}}{\cap}{\cap}}\displaylimits}%
  \def\lesseqqgtr{\mathrel{\lessgtr}}\def\bigoplus{\mathop{\oplus}\displaylimits}%
}
\providecommand{\ssi}{si et seulement si} % abréviation fréquente
\AtBeginDocument{\providecommand{\wideparen}{\overparen}} % arc de cercle

%% FIGFIX-BEGIN v3 — correctifs globaux des figures (docs/onbuch-generation/tools/figfix)
\usepackage{adjustbox}\usepackage{etoolbox}
% toute figure TikZ/pgfplots plus large que la ligne est réduite à la largeur disponible
\BeforeBeginEnvironment{tikzpicture}{\begin{adjustbox}{max width=\linewidth}}
\AfterEndEnvironment{tikzpicture}{\end{adjustbox}}
% légendes pgfplots lisibles par-dessus les courbes
\pgfplotsset{every axis/.append style={legend style={fill=white,fill opacity=0.92,text opacity=1,draw=black!15,inner sep=3pt}}}
% étiquettes d'arêtes (syntaxe edge["..."]) avec fond blanc
\tikzset{every edge quotes/.append style={fill=white,inner sep=1.2pt}}
% bulles de cartes mentales : texte à la ligne dans la bulle, bulle un peu plus grande
\tikzset{every concept/.append style={text width=2.0cm,align=center,minimum size=2.3cm,inner sep=2pt}}
%% FIGFIX-END
\begin{document}

\pagegarde{Traduction et production de textes divers liés à la citoyenneté}{Allemand}{Tle A}{Chapitre 5 : Citoyenneté}{AL30}{2 h}{Cette leçon mixte entraîne les élèves de Terminale A à produire et à traduire des textes variés (discours, lettre ouverte, interview, article) sur le thème de la citoyenneté. Elle combine lecture-commentaire d'un texte support, lexique thématique, révision générale de la grammaire du programme et entraînement au thème et à la version.
\vspace{4pt}
\begin{multicols}{2}\small
\begin{itemize}
\item À la fin de cette leçon, je sais lire et commenter un texte allemand sur la citoyenneté en dégageant le thème, la thèse et les arguments.
\item À la fin de cette leçon, je sais utiliser le lexique de la citoyenneté (droits, devoirs, élections, engagement) dans des phrases correctes.
\item À la fin de cette leçon, je sais rédiger un discours, une lettre ouverte, une interview ou un article sur la citoyenneté en respectant les conventions du genre.
\item À la fin de cette leçon, je sais traduire de courtes phrases et un texte court du français vers l'allemand (thème) en respectant l'ordre des mots et les cas.
\item À la fin de cette leçon, je sais traduire un texte court de l'allemand vers le français (version) sans calquer les structures allemandes.
\item À la fin de cette leçon, je sais réinvestir les structures grammaticales du programme (déclinaisons, subordonnées, passif, subjonctif) dans mes productions et traductions.
\end{itemize}\end{multicols}}

\begin{sommaire}
\begin{multicols}{2}
\begin{enumerate}
\item Texte support : lecture et compréhension
\item Lexique thématique : la citoyenneté
\item Méthode du commentaire de texte et commentaire modèle
\item Révision grammaticale ciblée pour la production et la traduction
\item Production de textes : discours, lettre ouverte, interview, article
\item Traduction : thème et version sur la citoyenneté
\item Activité d'intégration
\item Méthodes et exercices résolus
\item Exercices
\item Corrigés des exercices
\item Fiche bilan
\end{enumerate}
\end{multicols}
\end{sommaire}

\begin{objectifs}
\begin{itemize}
\item À la fin de cette leçon, je sais lire et commenter un texte allemand sur la citoyenneté en dégageant le thème, la thèse et les arguments.
\item À la fin de cette leçon, je sais utiliser le lexique de la citoyenneté (droits, devoirs, élections, engagement) dans des phrases correctes.
\item À la fin de cette leçon, je sais rédiger un discours, une lettre ouverte, une interview ou un article sur la citoyenneté en respectant les conventions du genre.
\item À la fin de cette leçon, je sais traduire de courtes phrases et un texte court du français vers l'allemand (thème) en respectant l'ordre des mots et les cas.
\item À la fin de cette leçon, je sais traduire un texte court de l'allemand vers le français (version) sans calquer les structures allemandes.
\item À la fin de cette leçon, je sais réinvestir les structures grammaticales du programme (déclinaisons, subordonnées, passif, subjonctif) dans mes productions et traductions.
\item À la fin de cette leçon, je sais relire et corriger ma production : orthographe, majuscules, place du verbe, cas et accords.
\item À la fin de cette leçon, je sais éviter les pièges de traduction fréquents entre le français et l'allemand (faux amis, prépositions, temps).
\end{itemize}
\end{objectifs}

\begin{prerequis}
\begin{itemize}
\item Lexique de la citoyenneté vu en Première (die Wahl, der Bürger, die Rechte und Pflichten)
\item Ordre des mots en proposition principale et subordonnée (verbe en 2e position / verbe final)
\item Déclinaisons des articles et des adjectifs (Nominativ, Akkusativ, Dativ, Genitiv)
\item Passif (Vorgangspassiv) et subordonnées avec weil, dass, obwohl, damit
\item Konjunktiv II pour exprimer l'opinion et le souhait (man sollte, ich würde)
\end{itemize}
\end{prerequis}

\newpage

\coursec{Texte support : lecture et compréhension}

\courssub{Situation d'accroche et problématique}

Aux dernières élections scolaires de ton lycée, tu as voté pour le délégué de classe : c'est déjà un acte de citoyenneté ! En Allemagne, les jeunes de 16 ans peuvent voter aux élections européennes, et le débat sur la participation citoyenne des jeunes anime régulièrement les médias. Imagine qu'un journaliste allemand te demande ton avis sur l'engagement des jeunes Camerounais : sauras-tu rédiger une lettre ouverte ou un article en allemand ? Et sauras-tu traduire son article pour tes camarades francophones ?

Cette leçon te donne les outils pour lire, écrire et traduire des textes sur la citoyenneté, compétence clé de l'épreuve du Bac. Tout commence par la lecture : avant de produire ou de traduire, il faut savoir \cle{comprendre} un texte allemand authentique. Notre problématique est donc la suivante : \textbf{comment lire, comprendre et reformuler un texte allemand sur la citoyenneté ?}

\begin{methode}[Lire un texte en 3 étapes]
Avant toute chose, retiens la démarche professionnelle de lecture d'un texte allemand :
\begin{enumerate}
\item \textbf{Survol} (\all{Überfliegen}) : lis rapidement le titre, la première et la dernière phrase. Objectif : identifier le \cle{thème} général et le type de texte, sans dictionnaire.
\item \textbf{Lecture détaillée} (\all{genaues Lesen}) : relis phrase par phrase, en t'aidant du glossaire. Souligne les mots-clés, repère les arguments et les exemples.
\item \textbf{Reformulation} (\all{Wiedergabe}) : résume les idées principales avec tes propres mots, d'abord en français, puis en allemand. C'est le premier pas vers la traduction (la \cle{version}).
\end{enumerate}
\end{methode}

\begin{attention}[Piège n°1 du lecteur francophone]
Ne traduis \textbf{jamais} mot à mot pendant la lecture ! Cherche d'abord le \cle{sens global}. En allemand, le verbe conjugué est en 2\up{e} position et l'infinitif ou le participe se trouve souvent \textbf{à la fin} de la phrase : si tu traduis mot à mot, tu perds le fil. Exemple : \all{In Kamerun engagieren sich viele Schüler in ihren Schulen.} — le sens se débloque quand tu as repéré le verbe réfléchi \all{sich engagieren} (s'engager) et son pronom \all{sich}. Lis la phrase entière avant de conclure.
\end{attention}

\courssub{Lecture du texte : « Jugendliche und die Politik »}

\textbf{Consigne de lecture.} Lis d'abord le texte en silence en appliquant l'étape 1 (survol). Puis ton professeur le lit à voix haute ; écoute la prononciation : \all{Jugendliche} se prononce « you-guend-li-ré », \all{Staatsbürgerschaft} « chtâts-bur-gueur-chaft », \all{Wahlbeteiligung} « vâl-bé-taï-li-goung ». Enfin, relis avec le glossaire.

\begin{textebox}[Jugendliche und die Politik — Zeitungsartikel]
In Deutschland dürfen Jugendliche ab 16 Jahren bei Europawahlen wählen. Viele Politiker meinen, dass junge Menschen sich früh für die Demokratie interessieren sollten. Andere finden, dass sie zu jung sind, um wichtige Entscheidungen zu treffen.

Die Wahlbeteiligung der jungen Wähler ist oft niedrig. „Viele Jugendliche denken, Politik ist nur etwas für Erwachsene“, sagt eine Lehrerin aus Berlin. „Aber das stimmt nicht: Politik bestimmt unser Leben — die Schule, die Ausbildung, sogar den Preis des Handys!“

Für das Wahlrecht ab 16 sprechen gute Argumente: Jugendliche, die früh wählen, bleiben oft ihr ganzes Leben aktive Wähler. Außerdem müssen junge Menschen heute wichtige Entscheidungen treffen, zum Beispiel bei der Berufswahl. Gegner des Wahlrechts ab 16 sagen dagegen: Sechzehnjährige haben zu wenig Erfahrung und lassen sich leicht von den Medien beeinflussen.

In Kamerun engagieren sich viele Schüler in ihren Schulen und Vierteln: Sie wählen ihre Klassensprecher, organisieren Sportfeste und helfen bei sozialen Projekten, zum Beispiel beim Müllsammeln oder beim Bau von Brunnen. Andere junge Leute leisten einen Freiwilligendienst in Krankenhäusern oder bei Hilfsorganisationen.

Das zeigt: Staatsbürgerschaft beginnt nicht erst mit 18 Jahren, sondern im Alltag. Wer in der Schule Verantwortung übernimmt, lernt die Demokratie praktisch. Vielleicht ist das der beste Weg zu einer starken Jugend — in Deutschland, in Österreich, in der Schweiz und in Kamerun.
\end{textebox}

\begin{definition}[Glossaire du texte]
\begin{itemize}
\item \all{die Staatsbürgerschaft} (pl. \all{-en}) : la nationalité, la citoyenneté
\item \all{der Staatsbürger} (\all{-}) : le citoyen
\item \all{das Wahlrecht} (\all{-e}) : le droit de vote
\item \all{die Europawahl} (\all{-en}) : l'élection européenne
\item \all{die Wahlbeteiligung} (\all{-en}) : la participation électorale
\item \all{sich engagieren} (verbe réfléchi) : s'engager — \all{Er engagiert sich in der Schule.} « Il s'engage dans son école. »
\item \all{der Freiwilligendienst} (\all{-e}) : le service volontaire, le bénévolat organisé — \all{einen Freiwilligendienst leisten} « effectuer un service volontaire »
\item \all{eine Entscheidung treffen} : prendre une décision — \all{eine wichtige Entscheidung treffen} « prendre une décision importante »
\item \all{der Klassensprecher} (\all{-}) : le délégué de classe
\item \all{die Verantwortung übernehmen} : assumer une responsabilité
\item \all{niedrig} (adj.) : bas, faible — contraire : \all{hoch} « élevé »
\end{itemize}
\end{definition}

\begin{savaistu}[Voter à 16 ans : une réalité européenne]
En Allemagne, le vote aux élections européennes est possible dès 16 ans. En Autriche, on peut voter à 16 ans à \textbf{toutes} les élections, y compris les élections nationales : c'est un cas unique en Europe. En Suisse, le débat est vif : certains cantons (comme Glaris) autorisent le vote à 16 ans pour les affaires cantonales. Et au Cameroun ? L'âge légal du vote est fixé à 20 ans, mais l'engagement citoyen des jeunes passe par d'autres voies : délégués de classe, clubs, associations de quartier, projets communautaires.
\end{savaistu}

\courssub{Repérage des informations générales}

Après le survol, complète la fiche d'identité du texte. C'est un réflexe exigé au Bac : tout commentaire de texte commence par la présentation du document.

\begin{center}
\begin{tabularx}{\linewidth}{|l|X|}
\hline
\rowcolor{popblueL} \textbf{Critère} & \textbf{Réponse} \\
\hline
Type de texte (\all{Textsorte}) & Article de presse (\all{Zeitungsartikel}) : il informe et argumente. \\
\hline
Source probable & Un journal ou un magazine pour jeunes, en Allemagne. \\
\hline
Thème (\all{Thema}) & L'engagement citoyen des jeunes : le droit de vote à 16 ans et l'engagement au quotidien. \\
\hline
Destinataires (\all{Leser}) & Les jeunes lecteurs, les élèves, mais aussi les adultes intéressés par la politique. \\
\hline
Intention de l'auteur & Informer, présenter les arguments pour et contre, et montrer que la citoyenneté commence dans la vie quotidienne. \\
\hline
\end{tabularx}
\end{center}

\textbf{Questions de compréhension globale} (à l'oral, en allemand) :
\begin{enumerate}
\item \all{Worum geht es im Text?} — De quoi parle le texte ? Réponse attendue : \all{Es geht um Jugendliche und die Politik, besonders um das Wahlrecht ab 16 und um das Engagement junger Menschen.} « Il s'agit des jeunes et de la politique, en particulier du droit de vote à 16 ans et de l'engagement des jeunes. »
\item \all{Wer spricht?} — Qui parle ? Un journaliste ; il cite aussi une professeure de Berlin.
\item \all{An wen richtet sich der Text?} — À qui s'adresse le texte ? Aux jeunes lecteurs et au public intéressé par la jeunesse et la démocratie.
\end{enumerate}

\courssub{Compréhension détaillée : arguments et exemples}

Le cœur du texte est un débat : faut-il accorder le droit de vote dès 16 ans ? Repérons les arguments avec précision.

\textbf{Arguments POUR (\all{Pro}) :}
\begin{itemize}
\item \all{Jugendliche, die früh wählen, bleiben oft ihr ganzes Leben aktive Wähler.} « Les jeunes qui votent tôt restent souvent des électeurs actifs toute leur vie. »
\item \all{Junge Menschen müssen heute wichtige Entscheidungen treffen, zum Beispiel bei der Berufswahl.} « Les jeunes doivent aujourd'hui prendre des décisions importantes, par exemple pour le choix du métier. »
\item Argument implicite de la professeure : la politique concerne la vie des jeunes (école, formation, prix du téléphone portable).
\end{itemize}

\textbf{Arguments CONTRE (\all{Contra}) :}
\begin{itemize}
\item \all{Sechzehnjährige haben zu wenig Erfahrung.} « Les jeunes de seize ans ont trop peu d'expérience. »
\item \all{Sie lassen sich leicht von den Medien beeinflussen.} « Ils se laissent facilement influencer par les médias. »
\end{itemize}

\textbf{Exemples d'engagement cités :} au Cameroun, les élèves élisent leurs délégués (\all{Klassensprecher}), organisent des fêtes sportives, ramassent les ordures (\all{Müllsammeln}), construisent des puits (\all{Brunnen}) ou effectuent un service volontaire (\all{einen Freiwilligendienst leisten}) dans les hôpitaux.

\begin{exemplebox}[La structure argumentative du texte]
Observe le connecteur d'opposition : \all{Gegner des Wahlrechts ab 16 sagen \textbf{dagegen}: ...} « Les adversaires du droit de vote à 16 ans objectent en revanche... ». Les mots \all{außerdem} (de plus), \all{zum Beispiel} (par exemple), \all{dagegen} (en revanche) sont les balises qui t'orientent dans le texte. Repère-les toujours en premier : ils annoncent la fonction de chaque phrase. Ajoute à ta liste : \all{aber} (mais), \all{deshalb} (c'est pourquoi), \all{das zeigt} (cela montre), qui introduit la conclusion.
\end{exemplebox}

\begin{popfigure}
\begin{tikzpicture}[
  box/.style={draw=popblue, fill=popblueL, rounded corners=2pt, align=center, font=\small, text width=3.1cm, minimum height=1cm},
  arr/.style={-{Stealth[length=3mm]}, thick, popdark}]
\node[box, align=center] (e) at (0,0){\textbf{Einleitung}\\ Vote à 16 ans\\ en Allemagne};
\node[box, fill=popgreenL, draw=popgreen, align=center] (p) at (4.6,0){\textbf{Argumente Pro}\\ électeurs actifs,\\ décisions importantes};
\node[box, fill=poporangeL, draw=poporange, align=center] (c) at (9.2,0){\textbf{Argumente Contra}\\ peu d'expérience,\\ influence des médias};
\node[box, fill=poppurpleL, draw=poppurple, align=center] (s) at (13.8,0){\textbf{Schluss}\\ la citoyenneté\\ commence au quotidien};
\draw[arr] (e) -- (p);
\draw[arr] (p) -- (c);
\draw[arr] (c) -- (s);
\end{tikzpicture}
\legende{Structure du texte : introduction, arguments pour, arguments contre, conclusion.}
\end{popfigure}

\courssub{Carte mentale du vocabulaire : la citoyenneté}

Pour réinvestir ce texte dans tes productions (discours, lettre ouverte, article), organise le lexique en cinq familles autour du mot central « citoyenneté » — en allemand \all{die Staatsbürgerschaft}.

\begin{popfigure}
\begin{tikzpicture}[
  mindmap,
  grow cyclic,
  every node/.style={concept, align=center, font=\small},
  concept color=popblue,
  level 1/.append style={level distance=3.4cm, sibling angle=72},
  level 2/.append style={level distance=2.2cm, sibling angle=45, font=\scriptsize}]
\node[concept color=popblue, text=white, align=center]{\all{Staatsbürgerschaft}\\ citoyenneté}
  child[concept color=popgreen] { node[align=center]{\all{Rechte}\\ droits}
    child { node {\all{das Wahlrecht}} }
    child { node {\all{die Meinungsfreiheit}} } }
  child[concept color=poporange] { node[align=center]{\all{Pflichten}\\ devoirs}
    child { node {\all{die Verantwortung}} }
    child { node {\all{Gesetze respektieren}} } }
  child[concept color=poppurple] { node[align=center]{\all{Wahlen}\\ élections}
    child { node {\all{die Wahlbeteiligung}} }
    child { node {\all{der Wähler}} } }
  child[concept color=poppink] { node {\all{Engagement}}
    child { node {\all{der Freiwilligendienst}} }
    child { node {\all{sich engagieren}} } }
  child[concept color=popgold] { node[align=center]{\all{Medien}\\ médias}
    child { node {\all{beeinflussen}} }
    child { node {\all{die Information}} } };
\end{tikzpicture}
\legende{Carte mentale lexicale : cinq familles de mots autour de la citoyenneté.}
\end{popfigure}

\begin{center}
\begin{tabularx}{\linewidth}{|X|X|X|}
\hline
\rowcolor{popblueL} \textbf{Français} & \textbf{Deutsch} & \textbf{Remarque} \\
\hline
la citoyenneté, la nationalité & \all{die Staatsbürgerschaft, -en} & mot composé : \all{Staat} + \all{Bürger} + \all{-schaft} \\
\hline
le droit de vote & \all{das Wahlrecht, -e} & neutre, comme tous les mots en \all{-recht} \\
\hline
s'engager & \all{sich engagieren} & verbe réfléchi : \all{ich engagiere mich} \\
\hline
la participation électorale & \all{die Wahlbeteiligung, -en} & féminin (suffixe \all{-ung}) \\
\hline
prendre une décision & \all{eine Entscheidung treffen} & expression figée : \all{treffen} = ici « prendre » \\
\hline
assumer une responsabilité & \all{Verantwortung übernehmen} & \all{übernehmen} : verbe fort, participe \all{übernommen} \\
\hline
le service volontaire & \all{der Freiwilligendienst, -e} & du volontariat au bénévolat organisé \\
\hline
influencer & \all{beeinflussen} & verbe faible régulier \\
\hline
la liberté d'expression & \all{die Meinungsfreiheit} & féminin (suffixe \all{-heit}) \\
\hline
respecter les lois & \all{die Gesetze respektieren} & \all{das Gesetz, -e} : la loi \\
\hline
\end{tabularx}
\end{center}

\begin{attention}[Faux amis et pièges de traduction]
\begin{itemize}
\item \all{die Politik} signifie « la politique » (l'activité) mais aussi « la politique menée » (les mesures) ; ne traduis jamais par « le politique ».
\item \all{aktiv} = « actif » ; attention, « les électeurs actuels » ne se dit pas avec \all{aktuell} dans ce sens : on dit \all{die heutigen Wähler}. \all{aktuell} signifie « d'actualité, actuel » au sens de « récent » : \all{aktuelle Nachrichten} « les nouvelles actuelles ».
\item \all{sich engagieren} exige le pronom réfléchi : on dit \all{Ich engagiere \textbf{mich} in meinem Viertel.} et jamais \all{ich engagiere in meinem Viertel}.
\item Majuscules obligatoires : \all{die Jugendlichen} (les jeunes), \all{die Erwachsenen} (les adultes) sont des adjectifs substantivés : ils prennent la majuscule comme tous les noms allemands, et se déclinent comme des adjectifs (\all{die jungen Wähler}, mais \all{die Jugendlichen}).
\item \all{wählen} (élire, voter) ne prend pas de préposition : \all{Sie wählen ihre Klassensprecher.} « Elles élisent leurs délégués de classe. » Ne dis pas \all{wählen für}.
\end{itemize}
\end{attention}

\courssub{Reformulation en français : premier pas vers la version}

Dernière étape de la méthode : reformuler les idées principales en français, avec tes propres mots. Cet exercice prépare directement l'épreuve de \cle{version} (allemand → français).

\textbf{Modèle de reformulation (idées principales du texte) :}
\begin{enumerate}
\item En Allemagne, les jeunes peuvent voter aux élections européennes dès 16 ans, mais la société est divisée sur ce droit de vote précoce.
\item Les partisans estiment que voter tôt forme des électeurs actifs pour toute la vie ; les adversaires jugent les jeunes trop inexpérimentés et trop influençables par les médias.
\item Au Cameroun, les jeunes s'engagent autrement : délégués de classe, projets sociaux, service volontaire.
\item Conclusion de l'auteur : la citoyenneté n'attend pas la majorité, elle s'apprend dans la vie quotidienne.
\end{enumerate}

\begin{exoresolu}[Reformuler une phrase difficile]
\textbf{Énoncé.} Reformule en français clair la phrase suivante, sans traduction mot à mot : \all{Wer in der Schule Verantwortung übernimmt, lernt die Demokratie praktisch.}
\tcblower
\textbf{Solution détaillée.}
\begin{enumerate}
\item \textbf{Repérage de la structure} : \all{Wer ... , ...} est une subordonnée relative générale = « Celui qui... ». Le verbe de la subordonnée (\all{übernimmt}) est en fin de proposition ; le verbe principal (\all{lernt}) vient juste après la virgule (règle de la 2\up{e} position : la subordonnée entière occupe la 1\up{re} place).
\item \textbf{Lexique} : \all{Verantwortung übernehmen} = « assumer une responsabilité » ; \all{praktisch} = ici « de manière concrète, en pratique » (et non « pratique » au sens d'« utile »).
\item \textbf{Reformulation} : « Celui qui assume des responsabilités à l'école apprend la démocratie de façon concrète. » Autre formulation possible : « Prendre des responsabilités à l'école, c'est apprendre la démocratie en la vivant. »
\end{enumerate}
\textbf{Erreur à éviter} : la traduction littérale « Qui à l'école responsabilité prend, apprend la démocratie pratiquement » est incompréhensible en français. On reformule toujours dans un français naturel.
\end{exoresolu}

\begin{exercice}[À toi de jouer]
\begin{enumerate}
\item Relis le texte et retrouve les trois exemples d'engagement des jeunes Camerounais cités par l'auteur.
\item Traduis en français : \all{Die Wahlbeteiligung der jungen Wähler ist oft niedrig.}
\item Réponds en allemand en une phrase complète : \all{Findest du das Wahlrecht ab 16 gut? Warum (nicht)?}
\item Reformule en français la conclusion du texte (dernier paragraphe) en deux phrases maximum.
\item Relève dans le texte deux connecteurs d'addition et un connecteur d'opposition, et traduis-les.
\end{enumerate}
\end{exercice}

\begin{corrige}[Exercice « À toi de jouer »]
\begin{enumerate}
\item Les élèves élisent leurs délégués de classe, organisent des fêtes sportives et participent à des projets sociaux (ramassage des ordures, construction de puits) ; d'autres effectuent un service volontaire dans les hôpitaux ou les organisations d'aide.
\item « La participation électorale des jeunes électeurs est souvent faible. »
\item Réponses possibles : \all{Ich finde das Wahlrecht ab 16 gut, weil junge Menschen dann früh Demokratie lernen.} / \all{Ich finde es nicht gut, weil Sechzehnjährige oft zu wenig Erfahrung haben.} (Attention : après \all{weil}, le verbe conjugué va à la fin.)
\item « La citoyenneté ne commence pas à 18 ans mais dans la vie quotidienne : assumer des responsabilités à l'école, c'est apprendre concrètement la démocratie. »
\item Addition : \all{außerdem} « de plus », \all{und} « et » ; opposition : \all{dagegen} « en revanche » (on accepte aussi \all{aber} « mais »).
\end{enumerate}
\end{corrige}

\begin{aretenir}[L'essentiel de cette section]
\begin{itemize}
\item Un texte allemand se lit en trois étapes : survol, lecture détaillée avec glossaire, reformulation.
\item On identifie toujours d'abord le type de texte, la source, le thème et le destinataire.
\item Un texte argumentatif s'organise autour d'arguments \all{Pro} et \all{Contra}, balisés par des connecteurs (\all{außerdem}, \all{dagegen}, \all{zum Beispiel}, \all{das zeigt}).
\item Le vocabulaire de la citoyenneté (\all{die Staatsbürgerschaft}, \all{das Wahlrecht}, \all{sich engagieren}, \all{die Wahlbeteiligung}, \all{der Freiwilligendienst}) est la base de toutes les productions et traductions de ce chapitre.
\item La reformulation en français est le premier pas vers la version : on restitue le sens, jamais les mots un à un.
\end{itemize}
\end{aretenir}

\coursec{Lexique thématique : la citoyenneté}

Dans la section précédente, nous avons lu et commenté un texte sur la participation citoyenne des jeunes. Vous avez remarqué que ce texte contenait de nombreux mots nouveaux : \all{das Wahlrecht}, \all{sich engagieren}, \all{die Pflicht}... Pour produire ou traduire un texte sur la citoyenneté — discours, lettre ouverte, interview ou article —, il faut d'abord maîtriser ce vocabulaire avec précision : le \cle{genre} de chaque nom, son \cle{pluriel}, et la construction exacte de chaque verbe. C'est l'objectif de cette section : construire, champ par champ, votre boîte à outils lexicale.

\begin{methode}[Comment apprendre un mot allemand]
\begin{enumerate}
\item Apprenez toujours un nom avec son \textbf{article} et son \textbf{pluriel} : \all{das Recht, die Rechte} (et non « Recht » tout seul).
\item Apprenez un verbe avec sa \textbf{construction} : \all{sich engagieren für + Akkusativ} (s'engager pour quelque chose).
\item Notez la prononciation en lettres françaises si elle est difficile : \all{die Gleichberechtigung} se prononce « glaïkh-be-rékh-ti-goung ».
\item Réemployez le mot dans une phrase complète, jamais dans une liste isolée.
\end{enumerate}
\end{methode}

\courssub{Champ lexical 1 : les droits}

\begin{definition}[Les droits — die Rechte]
\begin{center}
\begin{tabularx}{\linewidth}{|l|X|}
\hline
\rowcolor{popblueL} \textbf{Deutsch} & \textbf{Français} \\
\hline
das Recht (-e) & le droit \\
\hline
das Wahlrecht & le droit de vote \\
\hline
die Meinungsfreiheit & la liberté d'expression \\
\hline
die Gleichberechtigung & l'égalité des droits \\
\hline
das Demonstrationsrecht & le droit de manifestation \\
\hline
die Menschenrechte (Pl.) & les droits de l'homme \\
\hline
die Pressefreiheit & la liberté de la presse \\
\hline
\end{tabularx}
\end{center}
\end{definition}

Observez les exemples suivants, tirés de notre texte support :

\all{Jeder Bürger hat Rechte und Pflichten.} « Chaque citoyen a des droits et des devoirs. »

\all{In Deutschland gilt das Wahlrecht ab 18 Jahren.} « En Allemagne, le droit de vote s'applique à partir de 18 ans. »

\all{Die Meinungsfreiheit ist ein wichtiges Recht in einer Demokratie.} « La liberté d'expression est un droit important dans une démocratie. »

Remarquez la composition des mots allemands : \all{Wahl} (élection) + \all{Recht} (droit) = \all{Wahlrecht} ; \all{Meinung} (opinion) + \all{Freiheit} (liberté) = \all{Meinungsfreiheit}. C'est le principe des \cle{mots composés} : le dernier élément donne le genre et le sens principal.

\begin{propriete}[Les noms abstraits en -heit, -keit, -ung sont féminins]
Tous les noms allemands terminés en \textbf{-heit}, \textbf{-keit} et \textbf{-ung} sont féminins, sans exception :
\begin{itemize}
\item \all{die Freiheit} (la liberté), \all{die Gleichheit} (l'égalité), \all{die Menschheit} (l'humanité) ;
\item \all{die Möglichkeit} (la possibilité), \all{die Gerechtigkeit} (la justice) ;
\item \all{die Regierung} (le gouvernement), \all{die Meinung} (l'opinion), \all{die Gleichberechtigung} (l'égalité des droits).
\end{itemize}
Leur pluriel se forme en \textbf{-en} : \all{die Freiheit, die Freiheiten} ; \all{die Regierung, die Regierungen}.
\end{propriete}

\begin{attention}[Faux ami : das Recht]
\all{das Recht} signifie « le droit » au sens juridique (\all{das Wahlrecht} = le droit de vote), mais aussi « la justice d'une décision ». Attention : « avoir raison » se dit \all{Recht haben} : \all{Du hast Recht!} « Tu as raison ! ». En revanche, « la droite » (direction) se dit \all{die Rechte} ou \all{rechts}. Ne traduisez jamais « tourner à droite » par \all{nach Recht} : il faut \all{nach rechts}.
\end{attention}

\courssub{Champ lexical 2 : les devoirs}

\begin{definition}[Les devoirs — die Pflichten]
\begin{center}
\begin{tabularx}{\linewidth}{|l|X|}
\hline
\rowcolor{popblueL} \textbf{Deutsch} & \textbf{Français} \\
\hline
die Pflicht (-en) & le devoir \\
\hline
Steuern zahlen & payer des impôts \\
\hline
das Gesetz (-e) respektieren & respecter la loi \\
\hline
wählen gehen & aller voter \\
\hline
die Schulpflicht & l'obligation scolaire \\
\hline
der Wehrdienst & le service militaire \\
\hline
die Steuer (-n) & l'impôt, la taxe \\
\hline
\end{tabularx}
\end{center}
\end{definition}

\all{Steuern zahlen ist eine Pflicht für alle Bürger.} « Payer des impôts est un devoir pour tous les citoyens. »

\all{Wer das Gesetz nicht respektiert, wird bestraft.} « Celui qui ne respecte pas la loi est puni. »

\all{In Deutschland gibt es die Schulpflicht: alle Kinder müssen zur Schule gehen.} « En Allemagne, il y a l'obligation scolaire : tous les enfants doivent aller à l'école. »

\begin{exemplebox}[Droits et devoirs : la phrase clé]
\all{Jeder Bürger hat Rechte und Pflichten.} « Chaque citoyen a des droits et des devoirs. »

Cette phrase est un excellent point de départ pour toute production écrite sur la citoyenneté. Notez l'ordre des mots : sujet (\all{Jeder Bürger}) — verbe conjugué en 2\up{e} position (\all{hat}) — compléments.
\end{exemplebox}

\begin{attention}[« Devoir » : die Pflicht ou müssen ?]
Le nom « le devoir » se traduit par \all{die Pflicht}. Mais le verbe français « devoir » se traduit par le modal \all{müssen} : « Les citoyens doivent respecter la loi » = \all{Die Bürger müssen das Gesetz respektieren.} N'écrivez jamais \all{*Die Bürger pflichten das Gesetz}.
\end{attention}

\courssub{Champ lexical 3 : les élections et la vie politique}

\begin{definition}[Élections et politique]
\begin{center}
\begin{tabularx}{\linewidth}{|l|X|}
\hline
\rowcolor{popblueL} \textbf{Deutsch} & \textbf{Français} \\
\hline
die Wahl (-en) & l'élection ; le choix \\
\hline
wählen (wählt, wählte, hat gewählt) & choisir, élire, voter \\
\hline
der Wähler (-) & l'électeur \\
\hline
der Kandidat (-en) / die Kandidatin (-nen) & le candidat / la candidate \\
\hline
die Partei (-en) & le parti (politique) \\
\hline
die Regierung (-en) & le gouvernement \\
\hline
das Parlament (-e) & le parlement \\
\hline
die Stimme (-n) & la voix (électorale) \\
\hline
der Wahltag & le jour du scrutin \\
\hline
\end{tabularx}
\end{center}
\end{definition}

\all{Am Wahltag gehen Millionen Wähler zur Wahl.} « Le jour du scrutin, des millions d'électeurs vont voter. »

\all{Die Partei hat einen neuen Kandidaten gewählt.} « Le parti a élu un nouveau candidat. »

\all{Das Parlament wählt die Regierung.} « Le parlement élit le gouvernement. »

\begin{attention}[wählen / die Wahl / der Wähler : ne pas confondre]
Ces trois mots de la même famille ont des fonctions différentes :
\begin{itemize}
\item \all{wählen} est un \textbf{verbe} : « choisir » ou « élire ». \all{Ich wähle diesen Kandidaten.} « Je choisis/vote pour ce candidat. »
\item \all{die Wahl} est un \textbf{nom} : « l'élection » ou « le choix ». \all{Die Wahl findet im Oktober statt.} « L'élection a lieu en octobre. »
\item \all{der Wähler} désigne la \textbf{personne} : « l'électeur ». \all{Die Wähler entscheiden.} « Les électeurs décident. »
\end{itemize}
Erreur fréquente des francophones : écrire \all{*der Wahl} pour « l'électeur ». Retenez : \textbf{-er} = la personne qui fait l'action (comme \all{der Lehrer}, \all{der Arbeiter}).
\end{attention}

\begin{savaistu}[D'où vient le mot « Bürger » ?]
\all{der Bundesbürger} désigne officiellement le citoyen allemand (littéralement « citoyen de la Fédération »). Le mot \all{Bürger} vient de \all{die Burg} (le château fort) : au Moyen Âge, le « Bürger » était l'habitant de la ville fortifiée, qui jouissait de droits particuliers, contrairement au paysan. Le mot français « bourgeois » a exactement la même origine ! Aujourd'hui, \all{der Bürger} signifie simplement « le citoyen », et l'adjectif \all{bürgerlich} signifie « civil » : \all{das bürgerliche Engagement} = l'engagement citoyen.
\end{savaistu}

\courssub{Champ lexical 4 : l'engagement citoyen}

\begin{definition}[L'engagement — das Engagement]
\begin{center}
\begin{tabularx}{\linewidth}{|l|X|}
\hline
\rowcolor{popblueL} \textbf{Deutsch} & \textbf{Français} \\
\hline
sich engagieren (für + Akk.) & s'engager (pour) \\
\hline
der Freiwillige (-n) / die Freiwillige (-n) & le/la bénévole \\
\hline
das soziale Projekt (-e) & le projet social \\
\hline
helfen (hilft, half, hat geholfen) + Dat. & aider \\
\hline
die Solidarität & la solidarité \\
\hline
die Verantwortung & la responsabilité \\
\hline
die Freiwilligenarbeit & le bénévolat \\
\hline
die Jugendorganisation (-en) & l'organisation de jeunesse \\
\hline
\end{tabularx}
\end{center}
\end{definition}

\all{Die Jugendlichen engagieren sich in sozialen Projekten.} « Les jeunes s'engagent dans des projets sociaux. »

\all{Viele Freiwillige helfen den alten Menschen in der Nachbarschaft.} « Beaucoup de bénévoles aident les personnes âgées du quartier. »

\all{Solidarität bedeutet, Verantwortung für andere zu übernehmen.} « La solidarité signifie prendre ses responsabilités envers les autres. »

\begin{propriete}[Deux constructions à mémoriser]
\begin{enumerate}
\item \all{sich engagieren für + Akkusativ} : \all{Er engagiert sich für die Umwelt.} « Il s'engage pour l'environnement. » On trouve aussi \all{sich engagieren in + Dativ} (s'engager dans un projet, une association).
\item \all{helfen + Dativ} : le verbe \all{helfen} se construit toujours avec le datif, contrairement au français « aider quelqu'un ». \all{Ich helfe dem alten Mann.} « J'aide le vieil homme. » C'est un verbe fort : \all{helfen — hilft — half — hat geholfen}.
\end{enumerate}
\end{propriete}

\courssub{Champ lexical 5 : les problèmes et les débats}

\begin{definition}[Les problèmes de société]
\begin{center}
\begin{tabularx}{\linewidth}{|l|X|}
\hline
\rowcolor{popblueL} \textbf{Deutsch} & \textbf{Français} \\
\hline
die Korruption & la corruption \\
\hline
die Ungerechtigkeit (-en) & l'injustice \\
\hline
die Armut & la pauvreté \\
\hline
der Extremismus & l'extrémisme \\
\hline
die Desinformation & la désinformation \\
\hline
die Diskriminierung & la discrimination \\
\hline
der Protest (-e) & la protestation, la contestation \\
\hline
die Debatte (-n) & le débat \\
\hline
kämpfen gegen + Akk. & lutter contre \\
\hline
\end{tabularx}
\end{center}
\end{definition}

\all{Die Bürger kämpfen gegen die Korruption.} « Les citoyens luttent contre la corruption. »

\all{Die Desinformation in den sozialen Netzwerken ist eine Gefahr für die Demokratie.} « La désinformation sur les réseaux sociaux est un danger pour la démocratie. »

\all{Armut und Ungerechtigkeit sind große Probleme in vielen Ländern.} « La pauvreté et l'injustice sont de grands problèmes dans de nombreux pays. »

Notez que \all{die Armut} n'a pas de pluriel (nom abstrait), et que \all{kämpfen gegen} exige l'accusatif : \all{gegen die Korruption}.

\courssub{Expressions utiles pour l'argumentation}

Pour rédiger un discours, une lettre ouverte ou un article, vous devez exprimer votre opinion et structurer votre argumentation. Voici les connecteurs indispensables :

\begin{center}
\begin{tabularx}{\linewidth}{|X|X|X|}
\hline
\rowcolor{popblueL} \textbf{Français} & \textbf{Deutsch} & \textbf{Remarque} \\
\hline
à mon avis & meiner Meinung nach & suivi du verbe en 2\up{e} position \\
\hline
je suis d'avis que... & ich bin der Ansicht, dass... & subordonnée : verbe à la fin \\
\hline
d'une part... d'autre part... & einerseits... andererseits... & paire inséparable \\
\hline
il est important que... & es ist wichtig, dass... & subordonnée avec \all{dass} \\
\hline
à mon sens / selon moi & meines Erachtens & registre soutenu (discours) \\
\hline
je suis convaincu que... & ich bin überzeugt, dass... & verbe à la fin \\
\hline
en conclusion & abschließend / zum Schluss & pour clore un discours \\
\hline
\end{tabularx}
\end{center}

\all{Meiner Meinung nach ist die Meinungsfreiheit das wichtigste Recht.} « À mon avis, la liberté d'expression est le droit le plus important. »

\all{Ich bin der Ansicht, dass alle Jugendlichen wählen gehen sollten.} « Je suis d'avis que tous les jeunes devraient aller voter. »

\all{Einerseits haben die Bürger viele Rechte, andererseits müssen sie auch ihre Pflichten erfüllen.} « D'une part, les citoyens ont beaucoup de droits ; d'autre part, ils doivent aussi remplir leurs devoirs. »

\begin{attention}[Ordre des mots avec les expressions d'opinion]
Après \all{meiner Meinung nach} ou \all{einerseits}, le verbe conjugué reste en \textbf{deuxième position} : \all{Meiner Meinung nach \textbf{ist} die Wahl wichtig.} (et non \all{*Meiner Meinung nach die Wahl ist wichtig}). En revanche, après \all{dass}, le verbe conjugué va \textbf{à la fin} : \all{Ich bin der Ansicht, dass die Wahl wichtig \textbf{ist}.}
\end{attention}

\courssub{Carte mentale : le portrait du bon citoyen}

\begin{popfigure}
\begin{tikzpicture}[
  mindmap,
  grow cyclic,
  every node/.style={concept},
  root concept/.append style={concept color=popblue, font=\large\bfseries},
  level 1/.append style={level distance=3.6cm, sibling angle=90},
  level 1 concept/.append style={font=\small},
  concept color=popblue
]
\node[root concept]{der gute Bürger}
  child[concept color=popgreen] { node[align=center]{Rechte kennen\\ \footnotesize das Wahlrecht, die Meinungsfreiheit} }
  child[concept color=poporange] { node[align=center]{Pflichten erfüllen\\ \footnotesize Steuern zahlen, das Gesetz respektieren} }
  child[concept color=poppurple] { node[align=center]{wählen\\ \footnotesize wählen gehen, die Stimme abgeben} }
  child[concept color=poppink] { node[align=center]{sich engagieren\\ \footnotesize helfen, Solidarität zeigen} };
\end{tikzpicture}
\legende{Carte mentale lexicale : les quatre dimensions du « bon citoyen » (der gute Bürger).}
\end{popfigure}

\courssub{Texte d'application : un jeune engagé à Yaoundé}

\begin{textebox}[Ein engagierter Jugendlicher in Jaunde]
Marius ist 19 Jahre alt und wohnt in Jaunde, der Hauptstadt Kameruns. Er ist ein guter Bürger: Er kennt seine Rechte, aber er erfüllt auch seine Pflichten. „Meiner Meinung nach hat jeder Bürger Rechte und Pflichten", sagt er. „Einerseits haben wir das Recht auf Meinungsfreiheit, andererseits müssen wir das Gesetz respektieren."

Marius engagiert sich in einem sozialen Projekt: Jeden Samstag hilft er Kindern aus armen Familien bei den Hausaufgaben. „Die Freiwilligenarbeit ist wichtig für mich", erklärt er. „Ich bin der Ansicht, dass Solidarität unsere Gesellschaft stärker macht." Mit anderen Freiwilligen kämpft er auch gegen die Desinformation in den sozialen Netzwerken: Sie organisieren Workshops in Schulen und erklären den Schülern, wie man falsche Nachrichten erkennt.

Nächstes Jahr finden Wahlen statt. Marius wird zum ersten Mal wählen gehen. „Es ist wichtig, dass alle jungen Leute ihre Stimme abgeben", sagt er. „Die Wähler entscheiden über die Zukunft des Landes. Wer nicht wählt, darf sich später nicht beschweren!" Marius träumt davon, später Journalist zu werden, um über die Probleme seiner Gesellschaft zu berichten: die Armut, die Ungerechtigkeit, aber auch die vielen positiven Initiativen der Jugend.
\end{textebox}

\textbf{Glossaire :} \all{erfüllen} = remplir, accomplir ; \all{die Hausaufgaben (Pl.)} = les devoirs (scolaires) ; \all{stärker machen} = rendre plus fort ; \all{der Workshop (-s)} = l'atelier ; \all{erkennen (erkannte, hat erkannt)} = reconnaître ; \all{die Stimme abgeben} = voter (litt. « donner sa voix ») ; \all{sich beschweren} = se plaindre ; \all{berichten über + Akk.} = faire un reportage sur.

\textbf{Questions de compréhension :}
\begin{enumerate}
\item \all{Was macht Marius jeden Samstag?} (Que fait Marius chaque samedi ?)
\item \all{Gegen was kämpfen Marius und die anderen Freiwilligen?} (Contre quoi luttent-ils ?)
\item \all{Warum will Marius wählen gehen?} (Pourquoi veut-il aller voter ?)
\end{enumerate}

\begin{exoresolu}[Traduction guidée : du français vers l'allemand]
\textbf{Énoncé.} Traduisez en allemand : « À mon avis, chaque citoyen doit respecter la loi et aider les autres. D'une part, nous avons des droits ; d'autre part, nous avons aussi des devoirs. »
\tcblower
\textbf{Solution détaillée.}
\begin{enumerate}
\item « À mon avis » = \all{Meiner Meinung nach} ; le verbe suit immédiatement (2\up{e} position).
\item « chaque citoyen doit respecter la loi » : modal \all{müssen} en 2\up{e} position, infinitif \all{respektieren} à la fin : \all{jeder Bürger muss das Gesetz respektieren}.
\item « aider les autres » : \all{helfen + Dativ} : \all{den anderen helfen} (datif pluriel avec \all{-n} !).
\item « d'une part... d'autre part » = \all{einerseits... andererseits...}.
\end{enumerate}
\textbf{Traduction complète :} \all{Meiner Meinung nach muss jeder Bürger das Gesetz respektieren und den anderen helfen. Einerseits haben wir Rechte, andererseits haben wir auch Pflichten.}
\end{exoresolu}

\begin{exercice}[À vous de jouer]
\begin{enumerate}
\item Complétez avec le bon article : \all{... Wahl, ... Kandidat, ... Solidarität, ... Gesetz, ... Meinungsfreiheit, ... Parlament}.
\item Traduisez en allemand : « Les jeunes s'engagent pour la solidarité et luttent contre la pauvreté. »
\item Traduisez en français : \all{Ich bin der Ansicht, dass die Gleichberechtigung ein wichtiges Recht ist.}
\end{enumerate}
\end{exercice}

\begin{corrige}[Exercice : À vous de jouer]
\begin{enumerate}
\item \all{die Wahl, der Kandidat, die Solidarität, das Gesetz, die Meinungsfreiheit, das Parlament}.
\item \all{Die Jugendlichen engagieren sich für die Solidarität und kämpfen gegen die Armut.} (Attention : \all{gegen + Akkusativ}.)
\item « Je suis d'avis que l'égalité des droits est un droit important. » (Verbe \all{ist} à la fin dans la subordonnée avec \all{dass}.)
\end{enumerate}
\end{corrige}

\begin{aretenir}[L'essentiel du lexique de la citoyenneté]
\begin{itemize}
\item Droits : \all{das Recht, das Wahlrecht, die Meinungsfreiheit, die Gleichberechtigung} ; devoirs : \all{die Pflicht, Steuern zahlen, das Gesetz respektieren, wählen gehen}.
\item Vie politique : \all{die Wahl} (l'élection), \all{wählen} (voter), \all{der Wähler} (l'électeur) — trois mots à ne pas confondre.
\item Engagement : \all{sich engagieren für + Akk.}, \all{helfen + Dat.}, \all{der Freiwillige}, \all{die Solidarität}.
\item Argumentation : \all{meiner Meinung nach} + verbe en 2\up{e} position ; \all{ich bin der Ansicht, dass...} + verbe à la fin ; \all{einerseits... andererseits...}.
\item Tous les noms en \textbf{-heit}, \textbf{-keit}, \textbf{-ung} sont féminins.
\end{itemize}
\end{aretenir}

\coursec{Méthode du commentaire de texte et commentaire modèle}

Dans la section précédente, nous avons constitué un riche lexique thématique de la citoyenneté : \all{das Wahlrecht}, \all{die Staatsbürgerschaft}, \all{die Pflicht}, \all{die Verantwortung}, \all{die Demokratie}... Ces mots ne doivent pas rester inertes dans votre cahier de vocabulaire. Nous allons maintenant apprendre à les \emph{mobiliser} dans une tâche d'examen classique au Baccalauréat : le \cle{commentaire de texte} (\all{die Textanalyse} ou \all{der Kommentar}). Commenter un texte, ce n'est ni le résumer ni le recopier : c'est le présenter, poser la question qu'il soulève, analyser ses arguments, puis donner un avis personnel justifié. Cette démarche en quatre étapes est la colonne vertébrale de toutes vos productions écrites de Terminale.

\courssub{Repères pour identifier le type de texte (Étape 1)}

Avant d'entrer dans la méthode, rappelons les caractéristiques des quatre types de textes du programme que vous devrez identifier lors de la présentation du document.

\begin{definition}[Caractéristiques des types de textes au programme]
\begin{itemize}
\item \textbf{Article de presse (\all{Zeitungsartikel})} : titre accrocheur, auteur (souvent signé), date, lieu (\all{Berlin, -}), structure en paragraphes logiques (introduction, développement, conclusion), ton informatif ou argumentatif, pas d'adresse directe au lecteur.
\item \textbf{Discours (\all{Rede})} : allocution initiale (\all{Meine Damen und Herren,}), marqueurs d'oralité (\all{wir}, \all{ich}, questions rhétoriques), structure rhétorique (anaphores, parallélismes, climax), appel final à l'action (\all{Appell}), public présent.
\item \textbf{Lettre ouverte (\all{Offener Brief})} : en-tête (expéditeur, destinataire public/institution), formule d'appel formelle (\all{Sehr geehrte Frau Bundeskanzlerin,}), argumentation personnelle et engagée, ton solennel, signature, vocation à être publiée.
\item \textbf{Interview (\all{Interview})} : alternance stricte questions (journaliste \all{Frage}) / réponses (invité \all{Antwort}), marques de l'oral spontané (\all{also}, \all{ja}, \all{eigentlich}), tutoiement ou vouvoiement selon le contexte, sujet d'actualité brûlant.
\end{itemize}
\end{definition}

\courssub{Observer d'abord : un extrait de commentaire modèle}

Avant toute règle, lisons attentivement l'extrait suivant, tiré d'un commentaire rédigé sur un article consacré au droit de vote à 16 ans. Repérez les expressions qui présentent le texte, celles qui formulent la question, celles qui analysent et celle qui donne un avis.

\begin{textebox}[Extrait du commentaire modèle]
\all{Der Text zeigt zwei Meinungen über das Wahlrecht ab 16. Einerseits sind die Jugendlichen reif und interessiert, andererseits fehlt ihnen die Erfahrung. Ich finde, dass man früh lernen muss, Verantwortung zu tragen.}
\end{textebox}

\begin{center}
\begin{tabular}{|l|l|}
\hline
\rowcolor{popblueL} \textbf{Mot allemand} & \textbf{Français} \\
\hline
\all{die Meinung, die Meinungen} & l'opinion, l'avis \\
\hline
\all{das Wahlrecht} & le droit de vote \\
\hline
\all{reif} & mûr, mature \\
\hline
\all{fehlen (+ Dativ)} & manquer (à quelqu'un) \\
\hline
\all{die Erfahrung, die Erfahrungen} & l'expérience \\
\hline
\all{Verantwortung tragen} & assumer une responsabilité \\
\hline
\end{tabular}
\end{center}

Observons la mécanique : la première phrase \emph{présente} le contenu (\all{Der Text zeigt...}), la deuxième \emph{analyse} les deux positions grâce au couple de connecteurs \all{einerseits... andererseits}, et la dernière \emph{juge} (\all{Ich finde, dass...}). Toute la méthode du commentaire tient dans ces trois gestes, précédés d'une présentation formelle du document.

\courssub{Étape 1 : présenter le texte}

La première étape consiste à identifier le document : qui l'a écrit, où et quand il a paru, de quel type de texte il s'agit (article, discours, lettre ouverte, interview) et quel est son thème. En allemand, on utilise des formules figées qu'il faut connaître par cœur.

\begin{definition}[Formules de présentation du texte]
\begin{itemize}
\item \all{Der Text handelt von...} = Le texte traite de... (suivi du \textbf{datif} : \all{Der Text handelt vo\textbf{m} Wahlrecht ab 16.})
\item \all{Der Artikel erschien in der Zeitung...} = L'article a paru dans le journal...
\item \all{Der Autor / Die Autorin ist...} = L'auteur / L'auteure est...
\item \all{Es geht um...} = Il s'agit de... (suivi de l'\textbf{accusatif} : \all{Es geht um die Jugend und die Politik.})
\item \all{Das Thema des Textes ist...} = Le thème du texte est...
\item \all{Es handelt sich um einen Artikel / eine Rede / einen offenen Brief / ein Interview.} = Il s'agit d'un article / discours / lettre ouverte / interview.
\end{itemize}
\end{definition}

\begin{exemplebox}[Présenter le document]
\all{Der Text ist ein Artikel aus einer deutschen Zeitung. Er handelt von der Frage, ob Jugendliche ab 16 wählen dürfen.}\\
« Le texte est un article tiré d'un journal allemand. Il traite de la question de savoir si les jeunes peuvent voter à partir de 16 ans. »
\end{exemplebox}

\begin{attention}[handeln von : deux pièges]
Premier piège : \all{handeln} seul signifie « agir » ; c'est \all{handeln \textbf{von}} qui signifie « traiter de ». Deuxième piège : la préposition \all{von} exige le \textbf{datif} : on écrit \all{Der Text handelt vo\textbf{m} Thema} (von dem), jamais \all{von das Thema}. Les francophones oublient souvent cette déclinaison car le français « de » ne marque aucun cas.
\end{attention}

\courssub{Étape 2 : dégager la problématique}

Un bon commentaire ne décrit pas seulement le texte : il formule la \cle{problématique}, c'est-à-dire la question centrale à laquelle le texte répond. En allemand, la problématique s'exprime le plus souvent par une \textbf{interrogative indirecte}, introduite par \all{ob} (si) ou par un mot interrogatif (\all{wann}, \all{warum}, \all{wie}), avec le \textbf{verbe conjugué rejeté à la fin}.

\begin{propriete}[L'interrogative indirecte : forme et emploi]
\textbf{Forme} : verbe introducteur + \all{ob} / mot interrogatif + sujet + ... + verbe conjugué \textbf{en dernière position}.\\
\textbf{Emploi} : pour rapporter une question sans la poser directement.\\
\all{Die Frage ist: Sollen Jugendliche ab 16 wählen dürfen?} (question directe : le verbe \all{sollen} est en deuxième position)\\
$\rightarrow$ \all{Die Frage ist, ob Jugendliche ab 16 wählen dürfen \textbf{sollten}.} (question indirecte : \all{sollten} passe à la fin ; \all{dürfen} reste à l'infinitif devant lui)\\
« La question est : les jeunes devraient-ils pouvoir voter dès 16 ans ? »
\end{propriete}

\begin{exemplebox}[Formuler la problématique]
\all{Der Text stellt die Frage, ob das Wahlalter gesenkt werden soll.}\\
« Le texte pose la question de savoir si l'âge du vote doit être abaissé. »\\
\all{Man fragt sich, warum so wenige junge Leute wählen gehen.}\\
« On se demande pourquoi si peu de jeunes vont voter. »
\end{exemplebox}

\courssub{Étape 3 : analyser la structure et les arguments}

C'est le cœur du commentaire. Il faut montrer \emph{comment} le texte est construit : quels arguments sont avancés, dans quel ordre, pour quelle thèse. Pour un texte argumentatif, on distingue classiquement les arguments \cle{pour} (\all{dafür / pro}) et les arguments \cle{contre} (\all{dagegen / contra}). L'analyse s'organise grâce aux \textbf{connecteurs d'ordre} et aux \textbf{verbes de déclaration}.

\begin{center}
\begin{tabularx}{\linewidth}{|X|X|X|}
\hline
\rowcolor{popblueL} \textbf{Fonction} & \textbf{Deutsch} & \textbf{Français} \\
\hline
Ordre : premièrement & \all{zuerst} & d'abord, premièrement \\
\hline
Ordre : ensuite & \all{dann / danach} & ensuite \\
\hline
Ordre : de plus & \all{außerdem} & de plus, en outre \\
\hline
Ordre : enfin & \all{schließlich / zuletzt} & enfin, finalement \\
\hline
Opposition & \all{einerseits... andererseits} & d'une part... d'autre part \\
\hline
Opposition & \all{aber / jedoch} & mais / cependant \\
\hline
Concession & \all{zwar... aber} & certes... mais \\
\hline
Conclusion & \all{deshalb / daher} & c'est pourquoi \\
\hline
\end{tabularx}
\end{center}

\begin{propriete}[Les verbes de déclaration + dass]
Pour rapporter les arguments de l'auteur, on utilise des verbes de déclaration suivis d'une subordonnée en \all{dass} (verbe à la fin) :\\
\all{behaupten} (affirmer), \all{meinen} (estimer), \all{erklären} (expliquer), \all{betonen} (souligner), \all{kritisieren} (critiquer), \all{zugeben} (admettre).\\
\all{Der Autor behauptet, dass die Jugendlichen reif genug sind.}\\
« L'auteur affirme que les jeunes sont assez mûrs. »
\end{propriete}

\begin{attention}[Connecteurs et place du verbe]
\all{Zuerst}, \all{dann}, \all{außerdem} et \all{schließlich} occupent la \textbf{première place} de la phrase : le verbe conjugué vient immédiatement après, \emph{avant} le sujet (inversion). On écrit \all{Außerdem \textbf{sind} die Jugendlichen politisch interessiert.} et jamais \all{Außerdem die Jugendlichen sind...}. C'est l'erreur la plus fréquente des francophones, car en français « de plus, les jeunes sont... » garde l'ordre sujet-verbe.
\end{attention}

\courssub{Étape 4 : donner son avis argumenté}

Le commentaire se termine obligatoirement par votre \cle{avis personnel}, introduit par une formule d'opinion et \textbf{justifié} par au moins un argument (souvent introduit par \all{weil} ou \all{denn}).

\begin{definition}[Formules d'opinion]
\begin{itemize}
\item \all{Meiner Meinung nach...} = À mon avis... (suivi de l'inversion : \all{Meiner Meinung nach \textbf{ist} das Wählen ein Recht.})
\item \all{Ich bin der Meinung, dass...} = Je suis de l'avis que... (+ verbe à la fin)
\item \all{Ich finde, dass... / Ich denke, dass...} = Je trouve que... / Je pense que...
\item \all{Ich bin dafür / dagegen, dass...} = Je suis pour / contre le fait que...
\end{itemize}
\end{definition}

\begin{exemplebox}[Un avis justifié]
\all{Meiner Meinung nach ist das Wählen ein Recht und eine Pflicht.}\\
« À mon avis, voter est un droit et un devoir. »\\
\all{Ich bin der Meinung, dass Jugendliche ab 16 wählen dürfen sollten, weil sie früh Verantwortung lernen müssen.}\\
« Je suis d'avis que les jeunes devraient pouvoir voter dès 16 ans, parce qu'ils doivent apprendre tôt la responsabilité. »
\end{exemplebox}

\courssub{L'organigramme des quatre étapes}

\begin{popfigure}
\begin{center}
\begin{tikzpicture}[
  node distance=0.55cm,
  etape/.style={rectangle, rounded corners=4pt, draw=popblue, fill=popblueL, thick, text width=6.2cm, align=center, minimum height=1cm, font=\small},
  fleche/.style={-{Stealth[length=3mm]}, thick, popdark}
]
\node[etape, align=center] (e1){\textbf{1. Présentation}\\ \all{Der Text handelt von...}};
\node[etape, below=of e1, align=center] (e2){\textbf{2. Problématique}\\ \all{Die Frage ist, ob...}};
\node[etape, below=of e2, align=center] (e3){\textbf{3. Analyse pro/contra}\\ \all{zuerst, dann, außerdem, schließlich}};
\node[etape, below=of e3, align=center] (e4){\textbf{4. Avis personnel}\\ \all{Meiner Meinung nach...}};
\draw[fleche] (e1) -- (e2);
\draw[fleche] (e2) -- (e3);
\draw[fleche] (e3) -- (e4);
\end{tikzpicture}
\end{center}
\legende{Les quatre étapes obligatoires du commentaire de texte : présenter, problématiser, analyser, juger.}
\end{popfigure}

\begin{methode}[Comment réussir un commentaire au Bac]
\begin{enumerate}
\item \textbf{Lire} le texte deux fois et souligner la thèse, les arguments pour et les arguments contre.
\item \textbf{Présenter} le document en deux phrases : type, source, thème (\all{Der Text handelt von...}).
\item \textbf{Formuler} la problématique sous forme de question indirecte (\all{Die Frage ist, ob...}).
\item \textbf{Analyser} les arguments en suivant l'ordre du texte, avec les connecteurs \all{zuerst}, \all{dann}, \all{außerdem}, \all{schließlich} et les verbes de déclaration.
\item \textbf{Conclure} par un avis personnel justifié (\all{Meiner Meinung nach... weil...}).
\item \textbf{Relire} : vérifier la place du verbe, les majuscules des noms et les cas après les prépositions.
\end{enumerate}
\end{methode}

\begin{aretenir}[La règle d'or]
Le commentaire suit toujours le plan : \textbf{présenter} $\rightarrow$ \textbf{problématiser} $\rightarrow$ \textbf{analyser} $\rightarrow$ \textbf{juger}. Ne jamais résumer sans analyser : un résumé n'est pas un commentaire.
\end{aretenir}

\courssub{Commentaire modèle rédigé}

Voici un commentaire complet (environ dix lignes) sur l'article étudié en classe, consacré au droit de vote à 16 ans. Lisez-le en repérant les quatre étapes.

\begin{textebox}[Commentaire modèle : Wahlrecht ab 16]
\all{Der vorliegende Text ist ein Zeitungsartikel über das Wahlrecht ab 16. Der Autor stellt die Frage, ob Jugendliche in diesem Alter schon wählen dürfen sollten. Zuerst nennt er die Argumente der Befürworter: Die Jugendlichen sind heute gut informiert und politisch interessiert. Dann erklärt er, dass viele Entscheidungen der Politik ihre Zukunft betreffen. Außerdem behauptet er, dass das Wählen die junge Generation stärker an die Demokratie bindet. Andererseits erwähnt er die Kritiker: Sie meinen, dass den Sechzehnjährigen die Erfahrung fehlt. Schließlich betont der Autor, dass man in anderen Ländern schon mit 16 wählen darf. Meiner Meinung nach ist das Wählen ein Recht und eine Pflicht. Ich bin der Meinung, dass Jugendliche früh lernen müssen, Verantwortung zu tragen, weil sie die Bürger von morgen sind.}
\end{textebox}

\begin{center}
\begin{tabular}{|l|l|}
\hline
\rowcolor{popblueL} \textbf{Mot allemand} & \textbf{Français} \\
\hline
\all{vorliegend} & présent, en question \\
\hline
\all{der Befürworter, die Befürworter} & le partisan \\
\hline
\all{betreffen (betrifft, betraf, hat betroffen)} & concerner \\
\hline
\all{binden (bindet, band, hat gebunden)} & lier, attacher \\
\hline
\all{der Kritiker, die Kritiker} & le critique, l'opposant \\
\hline
\all{der Bürger, die Bürger} & le citoyen \\
\hline
\end{tabular}
\end{center}

Questions de vérification : 1) Quelle formule introduit la problématique ? 2) Combien d'arguments « pour » l'auteur cite-t-il ? 3) Quel connecteur introduit la position adverse ? 4) Comment l'élève justifie-t-il son avis final ?

\courssub{Entraînement guidé}

\begin{exoresolu}[Formuler la problématique]
Énoncé : transformez la question directe en problématique de commentaire.\\
\all{Dürfen Jugendliche in Kamerun schon mit 16 Jahren wählen?}
\tcblower
Solution détaillée : on introduit la question par \all{Die Frage ist, ob...} et l'on rejette le verbe conjugué \all{dürfen} à la fin de la subordonnée. L'ordre des autres éléments ne change pas.\\
\all{Die Frage ist, ob Jugendliche in Kamerun schon mit 16 Jahren wählen \textbf{dürfen}.}\\
« La question est de savoir si les jeunes au Cameroun peuvent voter dès l'âge de 16 ans. »
\end{exoresolu}

\begin{exercice}[Entraînement oral : une phrase chacun]
À partir du plan ci-dessous, chaque élève formule à voix haute \textbf{une phrase} du commentaire sur le thème « Faut-il un cours d'éducation civique obligatoire au lycée ? ».
\begin{enumerate}
\item Présentation : \all{Der Text handelt von...}
\item Problématique : \all{Die Frage ist, ob...}
\item Argument pour : \all{Zuerst behauptet der Autor, dass...}
\item Argument contre : \all{Andererseits meinen die Kritiker, dass...}
\item Avis personnel : \all{Meiner Meinung nach..., weil...}
\end{enumerate}
\end{exercice}

\begin{attention}[Le piège du Bac : la paraphrase]
Au Baccalauréat, le commentaire exige un \textbf{avis personnel justifié}. Une simple paraphrase du texte — répéter ce que dit l'auteur avec d'autres mots — fait perdre des points, même si l'allemand est correct. Le correcteur attend la trace de votre jugement : \all{Ich bin der Meinung, dass...}, \all{Meiner Meinung nach...}, suivi d'un argument personnel (\all{weil...}). Un commentaire sans avis est un commentaire inachevé.
\end{attention}

\begin{savaistu}[Le débat sur le vote à 16 ans]
En Allemagne, les jeunes de 16 ans peuvent voter aux élections européennes depuis 2024, et dans plusieurs Länder (comme le Brandebourg ou Brême) aux élections régionales. En Autriche, le vote à 16 ans existe pour toutes les élections depuis 2007 : c'est un pionnier en Europe. Ce débat passionne aussi les jeunes Camerounais, où la majorité électorale est fixée à 20 ans — un excellent sujet de comparaison pour votre avis personnel au Bac.
\end{savaistu}

\coursec{Révision grammaticale ciblée pour la production et la traduction}

Après avoir appris à commenter un texte, il faut maintenant sécuriser les outils grammaticaux indispensables pour \cle{produire} et \cle{traduire} sur le thème de la citoyenneté. Les examinateurs sanctionnent presque toujours les mêmes fautes : ordre des mots, déclinaisons, majuscules, particules séparables. Cette section reprend donc, de manière ciblée, les structures que vous devez maîtriser parfaitement en Terminale.

\courssub{Observation : lisons d'abord un court texte}

\begin{textebox}[Ein offener Brief an die Jugend — extrait]
Liebe Jugendliche,\\
heute engagieren sich viele junge Menschen in Vereinen, aber zu wenige gehen wählen. Ich meine, dass alle Bürger an den Wahlen teilnehmen sollten, weil die Demokratie nur funktioniert, wenn alle mitmachen. Es wäre gut, wenn jede Schule einen Wahlabend organisieren würde. Man sollte auch mehr junge Kandidaten wählen, damit die Politik jünger wird. Obwohl viele Gesetze respektiert werden, bleibt viel zu tun. Werdet aktiv! Eure Stimme zählt.
\end{textebox}

\textbf{Glossaire :} \all{der Verein, -e} (l'association) ; \all{die Wahl, -en} (l'élection) ; \all{teilnehmen an + Dat.} (participer à) ; \all{der Kandidat, -en} (le candidat) ; \all{die Stimme, -n} (la voix, le vote) ; \all{zählen} (compter) ; \all{mitmachen} (participer, s'impliquer) ; \all{organisieren} (organiser) ; \all{aktiv} (actif/active).

\textbf{Observez :} où se trouve le verbe conjugué dans \all{heute engagieren sich viele junge Menschen} ? Et dans \all{dass alle Bürger an den Wahlen teilnehmen sollten} ? C'est toute la différence entre la phrase principale et la subordonnée allemandes. Notez aussi \all{teilnehmen} et \all{mitmachen} : ce sont des \cle{verbes à particule séparable}.

\courssub{L'ordre des mots : la règle d'or de la phrase allemande}

\begin{propriete}[Le verbe conjugué en deuxième position]
Dans toute phrase déclarative allemande, le \cle{verbe conjugué} occupe la \textbf{deuxième position}, quoi qu'il y ait en première position (sujet ou complément). C'est la différence fondamentale avec le français.
\begin{itemize}
\item Sujet en tête : \all{Viele Jugendliche engagieren sich heute.} « Beaucoup de jeunes s'engagent aujourd'hui. »
\item Complément en tête : \all{Heute engagieren sich viele Jugendliche.} « Aujourd'hui, beaucoup de jeunes s'engagent. » — le sujet passe \textbf{après} le verbe : c'est l'\cle{inversion}.
\end{itemize}
\end{propriete}

\begin{attention}[Erreur fréquente des francophones]
Après un complément placé en tête, le français garde l'ordre sujet–verbe ; l'allemand, jamais.\\
Faux : *\all{Heute viele Jugendliche engagieren sich.}\\
Correct : \all{Heute engagieren sich viele Jugendliche.}\\
Autres exemples : \all{In Kamerun wählen die Bürger alle sieben Jahre.} « Au Cameroun, les citoyens votent tous les sept ans. » / \all{Nächstes Jahr werde ich zum ersten Mal wählen.} « L'année prochaine, je voterai pour la première fois. »
\end{attention}

\begin{propriete}[Le verbe à la fin dans les subordonnées]
Dans une subordonnée introduite par \all{dass} (que), \all{weil} (parce que), \all{obwohl} (bien que), \all{wenn} (quand, si), \all{damit} (pour que, afin que), le \cle{verbe conjugué va à la fin} de la subordonnée. S'il y a un auxiliaire ou un modal, il se place après l'infinitif ou le participe.
\begin{itemize}
\item \all{Ich meine, dass alle Bürger wählen gehen sollten.} « Je pense que tous les citoyens devraient aller voter. »
\item \all{Weil die Demokratie wichtig ist, müssen alle wählen.} « Parce que la démocratie est importante, tous doivent voter. »
\item \all{Obwohl das Gesetz streng ist, respektieren es nicht alle.} « Bien que la loi soit stricte, tous ne la respectent pas. »
\item \all{Man sollte mehr junge Kandidaten wählen, damit die Politik sich erneuert.} « On devrait élire davantage de jeunes candidats, afin que la politique se renouvelle. »
\end{itemize}
\textbf{Attention :} quand la subordonnée est en tête, elle compte comme la position 1 ; le verbe de la principale vient juste après la virgule : \all{Weil die Demokratie wichtig ist, \textbf{müssen} alle wählen.}
\end{propriete}

\begin{propriete}[Verbes à particule séparable : \cle{trennbare Verben}]
Beaucoup de verbes courants sont composés d'un radical et d'une particule (préposition ou adverbe) : \all{teilnehmen, mitmachen, zuhören, aufgeben, anfangen, ausgehen, einkaufen}.
\begin{itemize}
\item \textbf{Phrase principale :} la particule se détache et va \textbf{à la fin de la phrase} (après les compléments). Le radical conjugué reste en position 2.
\all{Ich nehme an der Versammlung teil.} « Je participe à l'assemblée. »
\item \textbf{Subordonnée :} la particule \textbf{ne se sépare pas}, le verbe entier (particule + radical conjugué) va à la fin.
\all{..., dass ich an der Versammlung \textbf{teilnehme}.}
\item \textbf{Infinitif (avec zu / modal) :} \all{zu} s'intercale entre la particule et le radical.
\all{Ich habe vor, an der Versammlung \textbf{teilzunehmen}.} / \all{Ich muss an der Versammlung \textbf{teilnehmen}.}
\item \textbf{Participe II :} le \all{ge-} s'intercale : \all{teilgenommen}.
\end{itemize}
\end{propriete}

\begin{popfigure}
\begin{tikzpicture}[
box/.style={draw=popblue, fill=popblueL, rounded corners, align=center, minimum height=0.9cm, font=\small},
sub/.style={draw=poporange, fill=poporangeL, rounded corners, align=center, minimum height=0.9cm, font=\small},
arr/.style={-{Stealth[length=2.5mm]}, thick, popdark}]
\node[box, align=center] (p1){Position 1\\ Sujet \emph{ou} complément\\ \all{Heute}};
\node[box, right=0.5cm of p1, align=center] (p2){Position 2\\ \textbf{Verbe conjugué}\\ \all{engagieren}};
\node[box, right=0.5cm of p2, align=center] (p3){Milieu\\ sujet, compléments\\ \all{sich viele Jugendliche}};
\node[box, right=0.5cm of p3, align=center] (p4){Fin\\ participe / infinitif /\\ \textbf{particule séparable}};
\draw[arr] (p1) -- (p2);
\draw[arr] (p2) -- (p3);
\draw[arr] (p3) -- (p4);
\node[sub, below=1.0cm of p1, align=center] (s1){Subordonnée\\ \all{dass, weil, wenn...}\\ \all{dass alle Bürger}};
\node[sub, below=1.0cm of p3, align=center] (s2){milieu de la\\ subordonnée\\ \all{an den Wahlen}};
\node[sub, below=1.0cm of p4, align=center] (s3){\textbf{Verbe conjugué}\\ \textbf{à la fin}\\ \all{teilnehmen sollten}};
\draw[arr] (s1) -- (s2);
\draw[arr] (s2) -- (s3);
\end{tikzpicture}
\legende{La structure de la phrase allemande : verbe en position 2 dans la principale (particule à la fin), verbe entier à la fin dans la subordonnée.}
\end{popfigure}

\courssub{Les quatre cas : déclinaisons et fonctions}

Pour produire et traduire sans faute, il faut reconnaître la \cle{fonction} de chaque groupe nominal et choisir le bon cas.

\begin{definition}[Les quatre cas et leurs fonctions]
\begin{itemize}
\item \textbf{Nominativ} : le sujet. \all{Der Bürger wählt.} « Le citoyen vote. »
\item \textbf{Akkusativ} : le COD. \all{Er respektiert das Gesetz.} « Il respecte la loi. »
\item \textbf{Dativ} : le COI (et après certaines prépositions : \all{mit, von, zu, bei, nach, aus, seit, gegenüber}). \all{Wir helfen dem Nachbarn.} « Nous aidons le voisin. »
\item \textbf{Genitiv} : la possession, la complétion du nom. \all{Die Rechte des Bürgers.} « Les droits du citoyen. »
\end{itemize}
\end{definition}

\begin{center}
\begin{tabular}{|l|l|l|l|l|}
\hline
\rowcolor{popblueL} Cas & der (masc.) & die (fém.) & das (neutre) & die (pluriel) \\
\hline
Nominativ & der / ein & die / eine & das / ein & die / -- \\
\hline
Akkusativ & den / einen & die / eine & das / ein & die / -- \\
\hline
Dativ & dem / einem & der / einer & dem / einem & den + n / -- + n \\
\hline
Genitiv & des + (e)s / eines + (e)s & der / einer & des + (e)s / eines + (e)s & der / -- \\
\hline
\end{tabular}
\end{center}

\begin{propriete}[La déclinaison de l'adjectif]
\begin{itemize}
\item \textbf{Déclinaison faible} (après article défini \all{der/die/das} ou \all{dieser/jeder/mancher}) : l'adjectif prend \textbf{-e} au nominatif/accusatif singulier (masc., fém., neutre) et au nominatif/accusatif pluriel \textbf{-en} ; partout ailleurs (datif, génitif, accusatif masc. sg.) il prend \textbf{-en}.
\all{der junge Kandidat} (Nom), \all{den jungen Kandidaten} (Akk), \all{dem jungen Kandidaten} (Dat), \all{des jungen Kandidaten} (Gen). Pluriel : \all{die jungen Kandidaten} (Nom/Akk), \all{den jungen Kandidaten} (Dat).
\item \textbf{Déclinaison forte} (sans article) : l'adjectif porte la marque du cas comme l'article défini. \all{freie Wahlen} (Nom/Akk Pl), \all{mit freiem Mut} (Dat Sg Neutre), \all{wegen freier Meinungsäußerung} (Gen Sg Fém).
\item \textbf{Déclinaison mixte} (après article indéfini \all{ein/kein/mein} ou possessifs) : l'adjectif prend la marque du genre \textbf{seulement là où l'article ne la porte pas} (Nom masc. \all{-er}, Nom/Acc neutre \all{-es}, Gen fém./pl. \all{-er}) ; ailleurs \textbf{-en}.
\all{ein junger Kandidat}, \all{ein freies Land}, \all{meiner freien Zeit} (Dat fém.), \all{keine jungen Kandidaten} (Akk Pl).
\end{itemize}
\end{propriete}

\begin{exemplebox}[Exemples traduits sur le thème de la citoyenneté]
\begin{itemize}
\item \all{Der junge Wähler gibt seine Stimme ab.} « Le jeune électeur dépose son vote. » (Nominativ)
\item \all{Die Partei unterstützt den jungen Kandidaten.} « Le parti soutient le jeune candidat. » (Akkusativ)
\item \all{Wir danken den freiwilligen Helfern.} « Nous remercions les bénévoles. » (Dativ pluriel, \textbf{-n} au nom et \textbf{-en} à l'adjectif)
\item \all{Die Meinung des Volkes ist wichtig.} « L'opinion du peuple est importante. » (Genitiv, \all{des Volkes})
\end{itemize}
\end{exemplebox}

\begin{attention}[Majuscules : le réflexe allemand]
En allemand, \textbf{tous} les noms communs prennent la majuscule, contrairement au français : \all{der Bürger} (le citoyen), \all{die Wahl} (l'élection), \all{das Recht} (le droit), \all{die Pflicht} (le devoir), \all{die Freiheit} (la liberté), \all{der Wahlabend} (la soirée électorale). En thème, chaque nom français traduit doit déclencher le réflexe « majuscule ». C'est la faute la plus fréquente — et la plus facile à éliminer à la relecture.
\end{attention}

\courssub{Le passif}

\begin{propriete}[Formation et emploi du passif]
Le passif se forme avec \all{werden} + \cle{participe II} : \all{werden} porte la conjugaison, le participe reste à la fin.
\begin{itemize}
\item Présent : \all{Das Gesetz \textbf{wird} respektiert.} « La loi est respectée. »
\item Prétérit : \all{Viele Projekte \textbf{wurden} von den Schülern organisiert.} « De nombreux projets ont été organisés par les élèves. »
\item Parfait : \all{Die Wahl \textbf{ist} verschoben \textbf{worden}.} « L'élection a été reportée. » (attention : \all{worden}, pas \all{geworden})
\item Avec modal : \all{Die Rechte müssen geschützt werden.} « Les droits doivent être protégés. »
\end{itemize}
L'agent est introduit par \all{von} + Dativ : \all{von den Schülern} « par les élèves ». Le passif est très fréquent dans les textes civiques et journalistiques allemands ; en version, pensez souvent à le traduire par « on » : \all{Es wird viel diskutiert.} « On discute beaucoup. »
\end{propriete}

\courssub{Le Konjunktiv II : exprimer l'avis et le conseil}

\begin{propriete}[Formes et emplois du Konjunktiv II]
Le Konjunktiv II exprime le souhait, le conseil, l'hypothèse, la politesse. Formes à connaître par cœur :
\begin{itemize}
\item \all{sein} : \all{wäre} — \all{Es wäre gut, wenn alle wählen gingen.} « Il serait bien que tous aillent voter. »
\item \all{haben} : \all{hätte} — \all{Wir hätten mehr Zeit.} « Nous aurions plus de temps. »
\item Modaux : \all{sollte} (devrait), \all{müsste} (faudrait), \all{könnte} (pourrait), \all{würde} + infinitif (conditionnel général pour tous les autres verbes).
\item \all{Man sollte die Gesetze respektieren.} « On devrait respecter les lois. »
\item \all{Man sollte mehr junge Kandidaten wählen.} « On devrait élire davantage de jeunes candidats. »
\end{itemize}
Dans la subordonnée en \all{wenn}, le verbe conjugué reste à la fin : \all{..., wenn alle wählen \textbf{gingen}.}
\end{propriete}

\begin{savaistu}[Le Konjunktiv II dans les discours politiques]
Le Konjunktiv II est extrêmement fréquent dans les discours et les interviews politiques en Allemagne, en Autriche et en Suisse : il permet de formuler des propositions sans les imposer. \all{Wir sollten mehr in die Bildung investieren.} « Nous devrions investir davantage dans l'éducation. » ; \all{Man müsste den Jugendlichen mehr zuhören.} « Il faudrait écouter davantage les jeunes. » Dans vos propres discours et lettres ouvertes, utilisez-le : c'est la marque d'un allemand soigné et nuancé.
\end{savaistu}

\courssub{Prépositions et cas : pièges pour la traduction}

\begin{center}
\begin{tabularx}{\linewidth}{|X|X|X|}
\hline
\rowcolor{popblueL} Préposition & Cas & Exemple (Citoyenneté) \\
\hline
\all{an, in, auf, über, unter, vor, hinter, neben, zwischen} (Wechselpräpositionen) & \textbf{Datif} (lieu \emph{où ?}) / \textbf{Akkusativ} (lieu \emph{où vers ?}) & \all{teilnehmen an + Dat.} (participer à) / \all{in die Politik gehen} (aller en politique = direction) \\
\hline
\all{mit, von, zu, bei, nach, aus, seit, gegenüber} & \textbf{Dativ} & \all{mit den Bürgern}, \all{von der Regierung}, \all{zu den Wahlen}, \all{gegenüber dem Parlament} \\
\hline
\all{durch, für, gegen, ohne, um} & \textbf{Akkusativ} & \all{für die Rechte}, \all{gegen das Gesetz}, \all{ohne Wahl} \\
\hline
\all{trotz, statt, während, wegen} & \textbf{Génitif} (langue soutenue) / souvent Dativ (langue courante) & \all{trotz des Verbots}, \all{wegen der Wahl} \\
\hline
\end{tabularx}
\end{center}

\courssub{Comparaison français–allemand : les pièges de la traduction}

\begin{center}
\begin{tabularx}{\linewidth}{|X|X|X|}
\hline
\rowcolor{popblueL} Français & Deutsch & Remarque \\
\hline
Il est important que tous votent. (subjonctif) & \all{Es ist wichtig, dass alle wählen.} & Pas de subjonctif en allemand : \textbf{indicatif} après \all{dass}, verbe à la fin. \\
\hline
Parce que la démocratie est importante, tous doivent voter. & \all{Weil die Demokratie wichtig ist, müssen alle wählen.} & Subordonnée en tête = position 1 ; inversion dans la principale. \\
\hline
On devrait respecter les lois. & \all{Man sollte die Gesetze respektieren.} & « On » = \all{man} ; « devrait » = \all{sollte} (Konj. II). \\
\hline
La loi est respectée par tous. & \all{Das Gesetz wird von allen respektiert.} & Passif : \all{werden} + participe II ; « par » = \all{von} + Dativ. \\
\hline
les droits de l'homme & \all{die Menschenrechte} & Majuscules aux noms ; mot composé fréquent en allemand. \\
\hline
Il participe à la réunion. & \all{Er nimmt an der Versammlung teil.} & Verbe séparable : particule \all{teil} à la fin (principale). \\
\hline
... qu'il participe à la réunion. & \all{..., dass er an der Versammlung teilnimmt.} & Subordonnée : verbe entier \all{teilnimmt} à la fin. \\
\hline
\end{tabularx}
\end{center}

\begin{attention}[Trois réflexes de relecture pour le thème]
\begin{enumerate}
\item Chaque \textbf{nom} a-t-il sa majuscule et son article décliné au bon cas ?
\item Le \textbf{verbe conjugué} est-il en position 2 (principale) ou à la fin (subordonnée) ? La \textbf{particule séparable} est-elle à sa place (fin de principale / collée au verbe en subordonnée) ?
\item Après « il est important que / il faut que », ai-je bien mis l'\textbf{indicatif} allemand et non un pseudo-subjonctif ? \all{Es ist wichtig, dass jeder Bürger seine Pflicht kennt.} « Il est important que chaque citoyen connaisse son devoir. »
\end{enumerate}
\end{attention}

\begin{exoresolu}[Thème guidé : phrases sur la citoyenneté]
Traduisez en allemand : 1. Aujourd'hui, beaucoup de jeunes s'engagent dans des associations. 2. Je pense que tous les citoyens devraient voter. 3. Bien que la loi soit stricte, elle n'est pas toujours respectée. 4. Il serait bon que chaque école organise un débat.
\tcblower
\textbf{Solution détaillée :}
\begin{enumerate}
\item Complément en tête $\Rightarrow$ inversion ; verbe séparable \all{sich engagieren} : \all{Heute engagieren sich viele Jugendliche in Vereinen.}
\item Subordonnée en \all{dass} $\Rightarrow$ verbe à la fin ; « devraient » = \all{sollten} : \all{Ich meine, dass alle Bürger wählen sollten.}
\item \all{obwohl} $\Rightarrow$ verbe à la fin ; subordonnée en tête $\Rightarrow$ inversion dans la principale ; passif au présent : \all{Obwohl das Gesetz streng ist, wird es nicht immer respektiert.}
\item Konjunktiv II : \all{wäre} ; subordonnée en \all{wenn} avec verbe final : \all{Es wäre gut, wenn jede Schule eine Debatte organisieren würde.}
\end{enumerate}
Points de vigilance : majuscules (\all{Jugendliche, Bürger, Gesetz, Schule, Debatte}), place du verbe, déclinaison de \all{jede Schule} (nominatif féminin), particule séparable (ex. 1).
\end{exoresolu}

\begin{exercice}[À vous de jouer]
Traduisez en allemand :
\begin{enumerate}
\item Au Cameroun, les citoyens doivent respecter les lois.
\item Il est important que les jeunes participent aux élections.
\item On devrait écouter davantage les jeunes, parce qu'ils ont de bonnes idées.
\item De nombreux projets ont été organisés par les élèves.
\end{enumerate}
Puis traduisez en français :
\all{Es wäre gut, wenn mehr Frauen in die Politik gingen, damit die Gesellschaft gerechter wird.}
\end{exercice}

\begin{corrige}[Exercice : À vous de jouer]
\begin{enumerate}
\item \all{In Kamerun müssen die Bürger die Gesetze respektieren.} (complément en tête, inversion ; modal \all{müssen} en position 2, infinitif final).
\item \all{Es ist wichtig, dass die Jugendlichen an den Wahlen teilnehmen.} (indicatif après \all{dass}, verbe à la fin ; \all{teilnehmen an} + Dativ pluriel \all{den Wahlen}).
\item \all{Man sollte den Jugendlichen mehr zuhören, weil sie gute Ideen haben.} (\all{zuhören} + Dativ $\rightarrow$ \all{den Jugendlichen} ; \all{weil} + verbe final \all{haben}).
\item \all{Viele Projekte wurden von den Schülern organisiert.} (passif au prétérit : \all{wurden} + participe II ; agent \all{von} + Dativ pluriel \all{den Schülern}).
\end{enumerate}
Version : « Il serait bon que davantage de femmes s'engagent en politique / aillent en politique, afin que la société devienne plus juste. » (\all{in die Politik gehen} = direction $\rightarrow$ Akkusatif ; \all{gingen} = Konj. II Präteritum de \all{gehen} ; \all{damit} + indicatif présent \all{wird} pour exprimer le but visé).
\end{corrige}

\begin{aretenir}[L'essentiel de la révision]
\begin{itemize}
\item Verbe conjugué en \textbf{position 2} dans la principale ; \textbf{inversion} après un complément en tête ; verbe \textbf{à la fin} après \all{dass, weil, obwohl, wenn, damit}.
\item \textbf{Verbes à particule séparable} : particule à la fin de la principale (\all{nimmt ... teil}), collée au verbe à la fin de la subordonnée (\all{... teilnimmt}).
\item Quatre cas, quatre fonctions : Nominativ (sujet), Akkusativ (COD), Dativ (COI / prépositions \all{mit, von, an+Dat...}), Genitiv (possession / prépositions \all{trotz, wegen...}) ; déclinaison article / nom / adjectif (faible, forte, mixte).
\item Passif : \all{werden} + participe II ; agent avec \all{von} + Dativ ; parfait \all{ist ... worden}.
\item Konjunktiv II pour le conseil et l'avis : \all{man sollte, es wäre gut, man müsste, würde + infinitif}.
\item Pas de subjonctif après « il est important que » ; majuscule à \textbf{tous} les noms.
\end{itemize}
\end{aretenir}

\coursec{Production de textes : discours, lettre ouverte, interview, article}

Dans la section précédente, nous avons révisé les structures grammaticales essentielles : la place du verbe, les cas, les connecteurs et les temps. Cette grammaire n'a de sens que si elle sert un objectif de communication. C'est précisément l'objet de cette section : apprendre à \cle{produire} des textes authentiques sur le thème de la citoyenneté, en respectant les conventions de quatre genres fondamentaux que l'on retrouve à l'examen comme dans la vie réelle : le discours, la lettre ouverte, l'interview et l'article de presse. Un élève de Terminale doit pouvoir écrire au maire de Yaoundé, interviewer un responsable associatif ou rédiger un article pour le journal du lycée — en allemand correct et adapté au destinataire.

\courssub{Avant d'écrire : les trois questions fondamentales}

\begin{methode}[Les trois questions avant toute production]
Avant d'écrire la moindre phrase, pose-toi toujours trois questions, dans cet ordre :
\begin{enumerate}
\item \textbf{À qui j'écris ?} (Wem schreibe ich ?) — le destinataire détermine le registre : \all{du} pour un camarade, \all{Sie} pour une autorité.
\item \textbf{Pourquoi j'écris ?} (Warum schreibe ich ?) — informer, convaincre, protester, proposer : le but détermine les arguments.
\item \textbf{Quel genre de texte ?} (Welche Textsorte ?) — le genre impose la structure et les formules fixes : on ne commence pas un discours comme une lettre.
\end{enumerate}
Si tu réponds clairement à ces trois questions, la moitié du travail est faite : le reste est une question de grammaire et de vocabulaire, que tu maîtrises déjà.
\end{methode}

\begin{center}
\begin{tabularx}{\linewidth}{|X|X|X|X|}
\hline
\rowcolor{popblueL} \textbf{Genre} & \textbf{Destinataire} & \textbf{Registre} & \textbf{Formules types} \\
\hline
die Rede (discours) & un public présent & soutenu, direct & \all{Liebe Mitbürger!} / \all{Ich danke Ihnen.} \\
\hline
der offene Brief (lettre ouverte) & une autorité + le public & formel & \all{Sehr geehrter Herr Bürgermeister,} / \all{Mit freundlichen Grüßen} \\
\hline
das Interview (interview) & un invité, via un média & formel ou courant & \all{Warum engagieren Sie sich...?} \\
\hline
der Artikel (article) & les lecteurs d'un journal & neutre, informatif & titre + chapeau + citations \\
\hline
\end{tabularx}
\end{center}

\courssub{Le discours (die Rede)}

\textbf{Observation.} Lis d'abord cet extrait de discours prononcé par une élève lors de la journée de la citoyenneté au lycée :

\begin{textebox}[Extrait de discours — die Rede]
„Liebe Mitschülerinnen und Mitschüler, sehr geehrte Lehrerinnen und Lehrer!

Heute spreche ich über ein Thema, das uns alle betrifft: unsere Verantwortung als junge Bürger. Viele sagen, die Jugend sei nicht interessiert an Politik. Ich sage: Das stimmt nicht! Wir kümmern uns um unsere Schule, wir helfen in unseren Vierteln, wir wählen die Delegierten unserer Klassen.

Aber wir können mehr tun. Deshalb schlage ich vor: Jede Klasse organisiert einmal im Monat eine Reinigungsaktion im Quartier. Das kostet nur zwei Stunden, aber es verändert alles.

Liebe Freunde, die Zukunft unseres Landes beginnt heute, und sie beginnt mit uns. Arbeiten wir zusammen! Ich danke Ihnen allen für Ihre Aufmerksamkeit."
\end{textebox}

\textbf{Glossaire :} \all{die Verantwortung} (la responsabilité) ; \all{der Bürger, -} (le citoyen) ; \all{sich kümmern um + Akk.} (s'occuper de) ; \all{die Reinigungsaktion, -en} (l'opération de nettoyage) ; \all{die Aufmerksamkeit} (l'attention) ; \all{der Delegierte, -n} (le délégué).

\begin{propriete}[Structure d'un discours (Aufbau einer Rede)]
Un discours suit toujours cinq étapes :
\begin{enumerate}
\item \textbf{Formule d'appel} (Anrede) : \all{Liebe Mitschüler!} / \all{Sehr geehrte Damen und Herren!} / \all{Liebe Mitbürger!}
\item \textbf{Introduction du thème} : \all{Heute spreche ich über...} (Aujourd'hui je parle de...)
\item \textbf{Arguments} avec connecteurs : \all{zuerst} (d'abord), \all{außerdem} (de plus), \all{deshalb} (c'est pourquoi).
\item \textbf{Appel à l'action} (Appell) : impératif ou \all{wir}-form : \all{Arbeiten wir zusammen!} (Travaillons ensemble !)
\item \textbf{Clôture} : \all{Ich danke Ihnen für Ihre Aufmerksamkeit.} (Je vous remercie de votre attention.)
\end{enumerate}
\end{propriete}

\begin{exemplebox}[Phrases utiles pour le discours]
\all{Liebe Mitbürger, wir müssen zusammenarbeiten!} = Chers concitoyens, nous devons travailler ensemble !

\all{Es ist Zeit, etwas zu ändern.} = Il est temps de changer quelque chose.

\all{Jeder von uns kann einen Beitrag leisten.} = Chacun de nous peut apporter sa contribution.
\end{exemplebox}

\courssub{La lettre ouverte (der offene Brief)}

La lettre ouverte est une lettre adressée officiellement à une personnalité (maire, ministre, directeur) mais destinée à être lue par le public, souvent publiée dans un journal. Elle combine donc la forme de la lettre formelle et la force argumentative du discours.

\begin{textebox}[Début de lettre ouverte modèle]
„Sehr geehrter Herr Bürgermeister,

als Schüler dieses Gymnasiums schreibe ich Ihnen, weil die Sauberkeit unseres Viertels uns alle betrifft. Jeden Morgen sehen wir Müll auf den Straßen, und Mülleimer fehlen oft oder sind überfüllt. Wir schlagen vor, dass die Stadt mehr Mülleimer aufstelle und dass die Schulen an einer großen Reinigungsaktion teilnehmen. Wir sind sicher: Wenn die Bürger und die Stadt zusammenarbeiten, wird unser Viertel sauberer und schöner.

Mit freundlichen Grüßen"
\end{textebox}

\textbf{Glossaire :} \all{die Sauberkeit} (la propreté) ; \all{das Viertel, -} (le quartier) ; \all{der Müll} (les ordures) ; \all{der Mülleimer, -} (la poubelle) ; \all{aufstellen} (installer) ; \all{teilnehmen an + Dat.} (participer à) ; \all{überfüllt} (plein à craquer, débordant).

\begin{propriete}[Structure et formules de la lettre formelle]
\begin{itemize}
\item \textbf{Expéditeur} (Absender) en haut à gauche, \textbf{destinataire} (Empfänger) en dessous, lieu et date à droite : \all{Yaoundé, den 15. März 2025}.
\item \textbf{Objet} (Betreff) : une phrase nominale courte, sans verbe conjugué : \all{Betreff: Sauberkeit in unserem Viertel}.
\item \textbf{Formule d'appel} (Anrede) : \all{Sehr geehrter Herr Bürgermeister,} / \all{Sehr geehrte Frau Direktorin,} — suivie d'une \textbf{virgule}, et la phrase suivante commence par une \textbf{minuscule}.
\item \textbf{Corps argumenté} : présentation du problème, arguments, propositions concrètes (\all{Wir schlagen vor, dass...} + \textbf{Konjunktiv I}).
\item \textbf{Formule de politesse} (Grußformel), sans virgule à la fin, puis la signature.
\end{itemize}
\end{propriete}

\begin{propriete}[Les formules de clôture (Grußformeln)]
\begin{center}
\begin{tabular}{|l|l|l|}
\hline
\rowcolor{popblueL} \textbf{Formule} & \textbf{Registre} & \textbf{Équivalent français} \\
\hline
\all{Hochachtungsvoll} & très formel (autorité supérieure) & Veuillez agréer... \\
\hline
\all{Mit freundlichen Grüßen} & formel standard & Cordialement / Salutations distinguées \\
\hline
\all{Freundliche Grüße} & semi-formel & Bien cordialement \\
\hline
\all{Viele Grüße} & moins formel (connaissance) & Amitiés \\
\hline
\all{Liebe Grüße} & amical & Bien à toi \\
\hline
\end{tabular}
\end{center}
\cle{Attention} : la Grußformel n'est \textbf{jamais} suivie d'une virgule en allemand, contrairement au français.
\end{propriete}

\begin{attention}[Le « Sie » majuscule : une faute éliminatoire]
Dans les lettres et interviews officielles, le vouvoiement s'écrit avec une \textbf{majuscule} : \all{Sie}, \all{Ihnen}, \all{Ihr}, \all{Ihre}. Écrire \all{ich danke sie} est une faute grave. Compare :
\begin{itemize}
\item \all{Wir bitten Sie um eine Antwort.} (Nous vous demandons une réponse.) — vouvoiement.
\item \all{Wir bitten sie um eine Antwort.} (Nous leur demandons une réponse.) — 3\textsuperscript{e} personne !
\end{itemize}
La majuscule distingue « vous » de « ils/elles ». De même : \all{Ihr Brief} (votre lettre) ≠ \all{ihr Brief} (leur/son lettre).
\end{attention}

\begin{popfigure}
\begin{center}
\begin{tikzpicture}[node distance=0.35cm,
box/.style={draw=popblue, fill=popblueL, rounded corners=3pt, text width=6.2cm, align=center, font=\small, inner sep=4pt},
arr/.style={-{Stealth[length=2.5mm]}, thick, popdark}]
\node[box] (abs) {\textbf{Absender} : Name, Adresse des Schreibers};
\node[box, below=of abs] (emp) {\textbf{Empfänger} : Herr Bürgermeister, Rathaus};
\node[box, below=of emp] (bet) {\textbf{Betreff} : Sauberkeit im Viertel};
\node[box, below=of bet] (anr) {\textbf{Anrede} : Sehr geehrter Herr Bürgermeister,};
\node[box, below=of anr] (txt) {\textbf{Text} : Problem — Argumente — Vorschläge};
\node[box, below=of txt] (gru) {\textbf{Grußformel} : Mit freundlichen Grüßen + Unterschrift};
\draw[arr] (abs) -- (emp);
\draw[arr] (emp) -- (bet);
\draw[arr] (bet) -- (anr);
\draw[arr] (anr) -- (txt);
\draw[arr] (txt) -- (gru);
\end{tikzpicture}
\end{center}
\legende{Schéma de la lettre ouverte : de l'expéditeur (Absender) à la formule de politesse (Grußformel).}
\end{popfigure}

\courssub{L'interview (das Interview)}

L'interview est un dialogue écrit entre un journaliste et un invité. Sa réussite repose sur la qualité des questions.

\begin{propriete}[Les questions ouvertes : les W-Fragen]
Pour obtenir des réponses développées, utilise les \textbf{questions ouvertes} introduites par un mot interrogatif (W-Fragen), avec le verbe en deuxième position :
\begin{center}
\begin{tabular}{|l|l|l|}
\hline
\rowcolor{popblueL} \textbf{Mot} & \textbf{Sens} & \textbf{Exemple} \\
\hline
\all{Was} & que / quoi & \all{Was macht Ihre Organisation?} \\
\hline
\all{Warum} & pourquoi & \all{Warum engagieren Sie sich für die Jugend?} \\
\hline
\all{Wie} & comment & \all{Wie kann man Bürger werden?} \\
\hline
\all{Wann} & quand & \all{Wann haben Sie angefangen?} \\
\hline
\all{Wie lange / Wie oft} & depuis combien / combien de fois & \all{Wie oft treffen Sie sich?} \\
\hline
\end{tabular}
\end{center}
Évite les questions fermées (réponse oui/non) ; si tu en poses une, ajoute une \textbf{relance} : \all{Können Sie das genauer erklären?} (Pouvez-vous expliquer cela plus précisément ?)
\end{propriete}

\begin{exemplebox}[Modèle d'échange d'interview]
\textbf{Journalist:} \all{Warum engagieren Sie sich für die Jugend?} = Pourquoi vous engagez-vous pour la jeunesse ?

\textbf{Invitée:} \all{Ich engagiere mich für die Jugend, weil junge Menschen die Zukunft unseres Landes sind.} = Je m'engage pour la jeunesse parce que les jeunes sont l'avenir de notre pays.

\textbf{Journalist:} \all{Was raten Sie den Schülern?} = Que conseillez-vous aux élèves ?

\textbf{Invitée:} \all{Ich rate ihnen, in der Schule fleißig zu sein und sich im Quartier zu engagieren.} = Je leur conseille d'être assidus à l'école et de s'engager dans leur quartier.
\end{exemplebox}

\begin{attention}[Ordre des mots dans les questions]
Après le mot interrogatif, le verbe vient \textbf{immédiatement}, puis le sujet : \all{Warum engagieren Sie sich?} et jamais « Warum Sie engagieren sich? ». De plus, \all{sich engagieren} est réfléchi : n'oublie pas le \all{sich} !
\end{attention}

\courssub{L'article de presse (der Artikel)}

L'article informe. Il obéit à une architecture précise, héritée de la presse germanophone :

\begin{propriete}[Structure d'un article (Aufbau eines Artikels)]
\begin{enumerate}
\item \textbf{Titre accrocheur} (Überschrift) : court, souvent nominal : \all{Schüler putzen das Viertel} (Des élèves nettoient le quartier).
\item \textbf{Chapeau} (Vorspann) : un ou deux phrases qui répondent aux questions \textbf{Wer ? Was ? Wo ? Wann ?} (qui ? quoi ? où ? quand ?).
\item \textbf{Paragraphes courts} : une idée par paragraphe, du plus important au moins important (pyramide inversée).
\item \textbf{Citations de témoins} entre guillemets allemands : \all{„Wir sind stolz auf unsere Arbeit", sagt ein Schüler.} (« Nous sommes fiers de notre travail », dit un élève.) Note que le verbe de parole suit la citation.
\end{enumerate}
\end{propriete}

\begin{exemplebox}[Chapeau d'article modèle]
\all{Yaoundé. Am Samstag haben 50 Schüler des Gymnasiums ihr Viertel gereinigt. Die Aktion war ein großer Erfolg.} = Yaoundé. Samedi, 50 élèves du lycée ont nettoyé leur quartier. L'opération a été un grand succès.

On y trouve : \textbf{où} (Yaoundé), \textbf{quand} (am Samstag), \textbf{qui} (50 Schüler), \textbf{quoi} (haben gereinigt).
\end{exemplebox}

\courssub{Production guidée : lettre ouverte au maire}

\begin{exoresolu}[Rédiger une lettre ouverte (10 à 12 lignes)]
\textbf{Énoncé.} En groupe, rédigez une lettre ouverte au maire de votre ville sur la propreté du quartier. Plan imposé : 1) formule d'appel ; 2) présentation du problème ; 3) deux arguments ; 4) deux propositions ; 5) formule de clôture.
\tcblower
\textbf{Solution détaillée (modèle de rédaction).}

\all{Sehr geehrter Herr Bürgermeister,}

\all{als Schüler des Gymnasiums schreiben wir Ihnen, weil die Sauberkeit in unserem Viertel ein großes Problem ist. Jeden Tag liegt Müll auf den Straßen, und viele Bürger werfen Plastik einfach weg. Das ist gefährlich für die Gesundheit, denn der Müll bringt Krankheiten. Außerdem gibt unser Viertel ein schlechtes Bild der Stadt. Deshalb schlagen wir vor, dass die Stadt mehr Mülleimer aufstelle und dass jede Schule einmal im Monat eine Reinigungsaktion organisiere. Wir Schüler sind bereit zu helfen. Wir hoffen auf Ihre Unterstützung und danken Ihnen im Voraus.}

\all{Mit freundlichen Grüßen}

\all{Die Schüler der Klasse Tle A}

\textbf{Commentaire de méthode :} la lettre respecte le plan (appel → problème → arguments avec \all{denn} et \all{außerdem} → propositions avec \all{Wir schlagen vor, dass...} + \textbf{Konjunktiv I} (\all{aufstelle}, \all{organisiere}) → clôture). Le vouvoiement \all{Sie/Ihnen/Ihre} porte toujours une majuscule ; les verbes sont en deuxième position ; la subordonnée avec \all{dass} rejette le verbe à la fin.
\end{exoresolu}

\begin{experience}[Atelier d'écriture en groupe]
Par groupes de quatre : 1) le groupe choisit un problème de citoyenneté (propreté, circulation, sécurité) ; 2) chaque membre rédige une partie de la lettre ouverte selon le plan ; 3) le groupe assemble, relit avec la grille ci-dessous ; 4) un porte-parole lit la lettre à la classe, qui joue le rôle du conseil municipal et pose des questions.
\end{experience}

\courssub{La grille de relecture}

\begin{methode}[Relire sa production en cinq passages]
Relis ton texte cinq fois, en ne cherchant qu'\textbf{un seul} type d'erreur à chaque passage :
\begin{enumerate}
\item \textbf{Orthographe} : Umlaute (ä, ö, ü), ß, double consonnes.
\item \textbf{Majuscules} : tous les noms, et \all{Sie/Ihnen/Ihr} du vouvoiement.
\item \textbf{Place du verbe} : 2\textsuperscript{e} position en phrase déclarative, fin de phrase après \all{weil, dass, wenn}, participe en fin au Perfekt.
\item \textbf{Cas} : article et adjectif corrects après préposition (\all{für die Jugend}, \all{mit dem Bürgermeister}).
\item \textbf{Connecteurs} : au moins trois connecteurs variés (\all{zuerst, außerdem, deshalb, trotzdem}).
\end{enumerate}
Cette méthode systématique fait gagner plusieurs points à l'examen : la plupart des fautes sont des erreurs d'inattention, pas d'ignorance.
\end{methode}

\begin{aretenir}[L'essentiel de la section]
\begin{itemize}
\item Chaque genre a ses formules obligatoires : \all{Liebe Mitbürger!} (discours), \all{Sehr geehrter Herr...} + \all{Mit freundlichen Grüßen} (lettre), W-Fragen (interview), titre + chapeau (article).
\item Le vouvoiement exige la majuscule : \all{Sie, Ihnen, Ihr}.
\item Après la formule d'appel d'une lettre : virgule, puis minuscule ; la Grußformel ne prend pas de virgule.
\item Dans une lettre formelle, après \all{vorschlagen, dass} / \all{fordern, dass} : \textbf{Konjunktiv I} (verbe à la fin, forme du subjonctif I : \all{aufstelle}, \all{organisiere}).
\item Avant d'écrire : À qui ? Pourquoi ? Quel genre ? — Après avoir écrit : la grille de relecture en cinq passages.
\end{itemize}
\end{aretenir}

\begin{exercice}[À toi de jouer]
1. Rédige le début d'un discours (4 lignes) pour la journée de l'arbre à ton lycée : formule d'appel, présentation du thème, un argument.

2. Écris trois questions d'interview (avec \all{Was}, \all{Warum}, \all{Wie}) pour un jeune volontaire qui nettoie les marchés de Douala.

3. Rédige le titre et le chapeau d'un article sur une élection de délégués de classe à ton lycée (réponds à : qui ? quoi ? où ? quand ?).
\end{exercice}

\begin{corrige}[Exercice : À toi de jouer]
1. Exemple : \all{Liebe Mitschüler, liebe Lehrer! Heute spreche ich über den Tag des Baumes. Bäume sind wichtig, denn sie geben uns Schatten und saubere Luft.} — Appel, thème (\all{ich spreche über}), argument introduit par \all{denn} (verbe en 2\textsuperscript{e} position après \all{denn}, conjonction de coordination).

2. Exemples : \all{Was machen Sie genau auf dem Markt?} / \all{Warum ist Ihnen die Sauberkeit so wichtig?} / \all{Wie können andere Jugendliche Ihnen helfen?} — Verbe en 2\textsuperscript{e} position, \all{Sie/Ihnen} avec majuscule. \all{wichtig} + Dat. (\all{Ihnen}), \all{helfen} + Dat. (\all{Ihnen}).

3. Exemple : Titre : \all{Tle A wählt ihre Delegierten}. Chapeau : \all{Yaoundé. Am Montag haben die Schüler der Klasse Tle A ihre neuen Klassendelegierten gewählt.} — Où (Yaoundé), quand (am Montag), qui (die Schüler), quoi (haben gewählt).
\end{corrige}

\coursec{Traduction : thème et version sur la citoyenneté}

Après avoir appris à \emph{produire} des textes en allemand (discours, lettre ouverte, interview, article), nous franchissons la dernière étape de cette leçon : la \cle{traduction}. Traduire, c'est produire un texte à partir d'un autre texte : les compétences que vous venez de mobiliser (ordre des mots, cas, temps, vocabulaire de la citoyenneté) sont exactement celles dont vous aurez besoin ici. Nous distinguerons le \cle{thème} (français $\rightarrow$ allemand) et la \cle{version} (allemand $\rightarrow$ français), les deux exercices classiques du baccalauréat.

\courssub{Observation : que se passe-t-il quand on traduit ?}

Observons d'abord trois phrases traduites dans les deux sens et repérons les transformations.

\begin{exemplebox}[Trois phrases, deux sens de traduction]
\textbf{Thème (français $\rightarrow$ allemand) :}\par
\all{Die Bürger müssen die Gesetze respektieren, weil die Demokratie wichtig ist.}\par
« Les citoyens doivent respecter les lois parce que la démocratie est importante. »\par\medskip
\textbf{Version (allemand $\rightarrow$ français) :}\par
\all{Viele Jugendliche bekommen die Chance, sich freiwillig zu engagieren.}\par
« Beaucoup de jeunes reçoivent la chance de s'engager bénévolement. »\par\medskip
\textbf{Thème :}\par
\all{In unserer Schule wählen die Schüler jedes Jahr ihre Vertreter.}\par
« Dans notre école, les élèves élisent leurs délégués chaque année. »
\end{exemplebox}

Qu'observe-t-on ? Trois types de transformations systématiques :

\begin{enumerate}
\item \textbf{L'ordre des mots change.} Dans la subordonnée introduite par \all{weil}, le verbe conjugué allemand \all{ist} va en \emph{fin} de phrase, alors que le français garde l'ordre normal : « parce que la démocratie \emph{est} importante ».
\item \textbf{Les cas remplacent les prépositions françaises.} « Respecter \emph{les} lois » devient \all{die Gesetze respektieren} (accusatif, sans préposition) ; « s'engager » devient \all{sich engagieren} (verbe pronominal avec \all{sich} à l'accusatif).
\item \textbf{Les temps ne coïncident pas toujours.} L'allemand n'a pas de passé composé distinct du présent dans certains emplois oraux, et le Präteritum allemand peut correspondre à un passé composé ou à un imparfait français selon le contexte.
\end{enumerate}

\begin{aretenir}[Le principe fondamental]
Traduire, ce n'est \textbf{jamais} remplacer mot à mot. C'est \emph{reformuler le sens} d'une langue dans le système grammatical de l'autre. Le mot français n'a pas toujours d'équivalent exact : c'est la \emph{phrase entière} qu'il faut reconstruire.
\end{aretenir}

\courssub{La méthode du thème (français $\rightarrow$ allemand)}

Le thème est l'exercice le plus exigeant : c'est vous qui construisez la phrase allemande. Voici la démarche en quatre étapes.

\begin{methode}[Traduire du français vers l'allemand en 4 étapes]
\begin{enumerate}
\item \textbf{Identifier le verbe} de chaque proposition et son infinitif allemand (attention aux verbes à particule : \all{sich engagieren}, \all{teilnehmen}).
\item \textbf{Placer le verbe correctement} : en position 2 dans la principale (\all{Die Bürger respektieren...}), en fin dans la subordonnée (\all{..., weil die Demokratie wichtig ist}), l'auxiliaire en position 2 et le participe en fin au Perfekt (\all{Die Schüler haben ihre Vertreter gewählt}).
\item \textbf{Accorder les cas} : déterminer la fonction de chaque nom (sujet = nominatif, COD = accusatif, COI = datif) et vérifier les prépositions à cas fixe (\all{für} + Akk., \all{mit} + Dat., \all{wegen} + Gen.).
\item \textbf{Vérifier les majuscules} (tous les noms !), l'orthographe (Umlaute, \all{ß}) et relire la phrase entière à voix haute.
\end{enumerate}
\end{methode}

\begin{popfigure}
\begin{center}
\begin{tikzpicture}[
  node distance=0.55cm and 1.1cm,
  box/.style={draw=popblue, fill=popblueL, rounded corners=3pt, align=center, text width=4.6cm, minimum height=1cm, font=\small},
  arr/.style={-{Stealth[length=3mm]}, thick, poporange}
]
\node[box, align=center] (lire){\textbf{1. LIRE}\\ Comprendre le sens global de la phrase française};
\node[box, below=of lire, align=center] (verbe){\textbf{2. IDENTIFIER LE VERBE}\\ Infinitif allemand ? Verbe à particule ? Subordonnée ?};
\node[box, below=of verbe, align=center] (phrase){\textbf{3. CONSTRUIRE LA PHRASE}\\ Verbe en position 2 ou en fin ; sujet et compléments aux bons cas};
\node[box, below=of phrase, align=center] (verif){\textbf{4. VÉRIFIER}\\ Cas, ordre des mots, majuscules, Umlaute, orthographe};
\draw[arr] (lire) -- (verbe);
\draw[arr] (verbe) -- (phrase);
\draw[arr] (phrase) -- (verif);
\draw[arr, dashed, poppurple] (verif.east) .. controls +(1.6,0) and +(1.6,0) .. (verbe.east) node[midway, right, font=\small\itshape, text=poppurple, align=center]{erreur ?\\ on recommence};
\end{tikzpicture}
\end{center}
\legende{Organigramme de la méthode de traduction (thème)}
\end{popfigure}

\begin{exoresolu}[Thème guidé : une phrase de citoyenneté]
\textbf{Énoncé.} Traduisez en allemand : « Les citoyens doivent respecter les lois parce que la démocratie est importante. »
\tcblower
\textbf{Solution détaillée.}\par
\textbf{Étape 1 — Les verbes.} « Doivent respecter » = verbe modal \all{müssen} + infinitif \all{respektieren} ; « est » = \all{sein}. La subordonnée est introduite par « parce que » = \all{weil}.\par
\textbf{Étape 2 — La place des verbes.} Principale : \all{müssen} en position 2, \all{respektieren} en fin. Subordonnée avec \all{weil} : \all{ist} en fin absolue.\par
\textbf{Étape 3 — Les cas.} « Les citoyens » = sujet $\rightarrow$ \all{die Bürger} (nominatif pluriel). « Les lois » = COD de \all{respektieren} $\rightarrow$ \all{die Gesetze} (accusatif pluriel, identique au nominatif au pluriel). « La démocratie » = sujet de la subordonnée $\rightarrow$ \all{die Demokratie}.\par
\textbf{Étape 4 — Vérification.} Majuscules : \all{Bürger}, \all{Gesetze}, \all{Demokratie}. Umlaut : \all{Bürger}.\par
\textbf{Traduction :} \all{Die Bürger müssen die Gesetze respektieren, weil die Demokratie wichtig ist.}
\end{exoresolu}

\begin{attention}[Le piège de la subordonnée]
Après \all{weil}, \all{dass}, \all{obwohl}, \all{wenn}, le verbe conjugué allemand va \textbf{en fin de phrase}, alors que le français garde l'ordre normal. Erreur typique du francophone : \all{*weil die Demokratie ist wichtig}. Correct : \all{..., weil die Demokratie wichtig IST.} = « ...parce que la démocratie EST importante. » Pensez à « pousser » le verbe au bout de la subordonnée avant d'écrire la suite.
\end{attention}

\courssub{La méthode de la version (allemand $\rightarrow$ français)}

En version, le danger est inverse : il ne faut pas \emph{calquer} la syntaxe allemande. Le français attendu doit être naturel, idiomatique, élégant.

\begin{methode}[Traduire de l'allemand vers le français en 4 étapes]
\begin{enumerate}
\item \textbf{Lire tout le texte} avant de traduire la moindre phrase : repérer le thème, le ton, l'époque (temps des verbes).
\item \textbf{Comprendre le sens global} de chaque phrase : découper en segments de sens (sujet / verbe / compléments), en repérant le verbe conjugué (position 2) et les infinitifs ou participes renvoyés en fin.
\item \textbf{Restituer en français naturel} : remettre le verbe à sa place française, remplacer les tournures allemandes par des tournures françaises (\all{Es gibt viele Möglichkeiten} = « Il existe de nombreuses possibilités », non « \emph{Il donne beaucoup de possibilités} »).
\item \textbf{Relire uniquement le français} : la phrase doit pouvoir être prononcée par un francophone qui ne connaît pas l'allemand, sans qu'on devine la langue d'origine.
\end{enumerate}
\end{methode}

\begin{propriete}[Règle d'or de la version]
En version, on traduit \textbf{le sens, jamais les mots isolés}. Un mot allemand peut avoir plusieurs équivalents français selon le contexte : \all{sich engagieren} = « s'engager », « s'investir », « faire du bénévolat » ; c'est le contexte qui décide.
\end{propriete}

\begin{savaistu}[La version au baccalauréat]
Au Bac, la version est notée sur \textbf{deux critères} : la \emph{fidélité au sens} du texte allemand ET la \emph{qualité du français}. Une traduction trop littérale — le fameux « calque » — est pénalisée, même si tous les mots sont « justes ». Écrire « \emph{Les jeunes reçoivent la possibilité, eux-mêmes à engager} » pour \all{Jugendliche bekommen die Möglichkeit, sich zu engagieren} fait perdre des points : il fallait écrire « Les jeunes ont la possibilité de s'engager. »
\end{savaistu}

\courssub{Les pièges fréquents : faux amis, prépositions, temps}

\begin{attention}[Faux amis n° 1 : bekommen]
\all{bekommen} signifie « \textbf{recevoir} », jamais « devenir » ! « Devenir » se dit \all{werden}. Ainsi \all{Ich bekomme 18 Jahre} est un faux sens (on dirait : « je reçois 18 ans ») ; correct : \all{Ich werde 18 Jahre alt.} = « Je vais avoir 18 ans. » Et \all{Viele Jugendliche bekommen die Chance} = « Beaucoup de jeunes \emph{reçoivent / ont} la chance ».
\end{attention}

\begin{center}
\begin{tabularx}{\linewidth}{|l|l|X|X|}
\hline
\rowcolor{popblueL} \textbf{Mot allemand} & \textbf{Vrai sens} & \textbf{Faux ami français} & \textbf{Exemple traduit} \\
\hline
\all{bekommen} & recevoir & $\neq$ « devenir » (= \all{werden}) & \all{Sie bekommt einen Brief.} = Elle reçoit une lettre. \\
\hline
\all{aktuell} & actuel, en ce moment & $\neq$ « actuel » au sens de « réel » & \all{das aktuelle Thema} = le sujet d'actualité \\
\hline
\all{die Mappe, -n} & la pochette, le classeur & $\neq$ « la mappe » au sens de carte & \all{die Landkarte} = la carte géographique \\
\hline
\all{der Chef, -s} & le patron, le chef d'entreprise & $\neq$ « le chef » (cuisinier = \all{der Koch}) & \all{Mein Chef engagiert sich.} = Mon patron s'engage. \\
\hline
\all{die Wahl, -en} & le choix ; l'élection & $\neq$ seulement « élection » & \all{die freie Wahl} = le libre choix \\
\hline
\end{tabularx}
\end{center}

\begin{propriete}[Prépositions à cas fixe à revoir pour la citoyenneté]
\begin{itemize}
\item \all{für} + \textbf{accusatif} : \all{sich für die Umwelt engagieren} = « s'engager pour l'environnement » ;
\item \all{mit} + \textbf{datif} : \all{mit den Nachbarn sprechen} = « parler avec les voisins » ;
\item \all{wegen} + \textbf{génitif} : \all{wegen des Gesetzes} = « à cause de la loi » (au lycée, le datif \all{wegen dem Gesetz} est toléré à l'oral, mais préférez le génitif à l'écrit) ;
\item \all{an} + datif : \all{an einer Wahl teilnehmen} = « participer à une élection ».
\end{itemize}
\end{propriete}

\begin{propriete}[Präteritum allemand et passé français]
Le Präteritum allemand correspond, selon le contexte, au \textbf{passé composé} ou à l'\textbf{imparfait} français : \all{Letztes Jahr nahmen viele Schüler an der Wahl teil.} = « L'année dernière, beaucoup d'élèves \emph{ont participé} à l'élection. » En version, choisissez le temps français selon le sens (action ponctuelle = passé composé ; habitude ou description = imparfait), pas selon la forme allemande.
\end{propriete}

\courssub{Texte de version : le bénévolat des jeunes en Allemagne}

\begin{textebox}[Jugendliche engagieren sich]
\all{In Deutschland engagieren sich viele Jugendliche freiwillig. Sie arbeiten zum Beispiel in Sportvereinen, bei der Feuerwehr oder in Umweltgruppen. Viele Schüler bekommen in der Schule die Chance, an sozialen Projekten teilzunehmen: Sie helfen älteren Menschen, sammeln Geld für Kinder oder pflanzen Bäume in ihrer Stadt. „Ich habe letztes Jahr mit meiner Klasse einen Spielplatz gebaut“, erzählt Lena, 17 Jahre alt, aus Hamburg. „Das war schwer, aber wir haben viel gelernt.“ Politiker sagen oft, dass das Ehrenamt für die Demokratie wichtig ist, weil die Bürger dabei Verantwortung lernen. Auch in Kamerun engagieren sich junge Leute: Sie organisieren Fußballturniere oder reinigen die Straßen ihres Viertels. Engagement verbindet die Menschen — überall auf der Welt.}
\end{textebox}

\textbf{Glossaire :} \all{freiwillig} = bénévolement ; \all{der Sportverein, -e} = le club de sport ; \all{die Feuerwehr} = les pompiers ; \all{die Umweltgruppe, -n} = le groupe écologiste ; \all{teilnehmen an} (+ dat.) = participer à ; \all{das Geld sammeln} = collecter de l'argent ; \all{der Spielplatz, -plätze} = l'aire de jeux ; \all{das Ehrenamt, -er} = le bénévolat ; \all{die Verantwortung} = la responsabilité ; \all{das Fußballturnier, -e} = le tournoi de football ; \all{das Viertel, -} = le quartier ; \all{verbinden} = relier, unir.

\textbf{Questions de compréhension :}
\begin{enumerate}
\item \all{Wo engagieren sich die deutschen Jugendlichen? Nennen Sie drei Beispiele.}
\item \all{Was hat Lena letztes Jahr gemacht?}
\item \all{Warum ist das Ehrenamt laut den Politikern wichtig für die Demokratie?}
\item Quel parallèle le texte établit-il avec le Cameroun ?
\end{enumerate}

\begin{exoresolu}[Version guidée : deux phrases du texte]
\textbf{Énoncé.} Traduisez en français : \all{Viele Schüler bekommen in der Schule die Chance, an sozialen Projekten teilzunehmen. Politiker sagen oft, dass das Ehrenamt für die Demokratie wichtig ist, weil die Bürger dabei Verantwortung lernen.}
\tcblower
\textbf{Solution détaillée.}\par
\textbf{Phrase 1.} Segments de sens : \all{Viele Schüler} (sujet) / \all{bekommen} (verbe : « reçoivent », faux ami !) / \all{die Chance} (COD) / \all{an sozialen Projekten teilzunehmen} (infinitive : « participer à des projets sociaux »). Traduction naturelle : « Beaucoup d'élèves \emph{ont} la chance, à l'école, de participer à des projets sociaux. » (On préfère « avoir la chance » au calque « recevoir la chance ».)\par
\textbf{Phrase 2.} Subordonnée \all{dass} : verbe \all{ist} en fin ; subordonnée \all{weil} : verbe \all{lernen} en fin. Traduction : « Les hommes politiques répètent souvent que le bénévolat est important pour la démocratie, parce que les citoyens y apprennent la responsabilité. » Notez : \all{dabei} = « en le faisant, dans cette activité » $\rightarrow$ « y » en français naturel.
\end{exoresolu}

\courssub{Entraînement au thème}

\begin{exercice}[Thème : cinq phrases sur la citoyenneté]
Traduisez en allemand :
\begin{enumerate}
\item Les citoyens doivent respecter les lois parce que la démocratie est importante.
\item Beaucoup de jeunes s'engagent pour l'environnement.
\item L'année dernière, les élèves ont élu leurs délégués.
\item Bien que le travail soit difficile, les bénévoles aident les personnes âgées.
\item Dans notre quartier, nous participons à un projet social avec nos voisins.
\end{enumerate}
\end{exercice}

\begin{corrige}[Exercice : thème]
\begin{enumerate}
\item \all{Die Bürger müssen die Gesetze respektieren, weil die Demokratie wichtig ist.} (Modal \all{müssen} en position 2, infinitif en fin ; verbe \all{ist} en fin de subordonnée \all{weil}.)
\item \all{Viele Jugendliche engagieren sich für die Umwelt.} (Verbe pronominal \all{sich engagieren} ; \all{für} + accusatif.)
\item \all{Letztes Jahr haben die Schüler ihre Vertreter gewählt.} (Perfekt : auxiliaire \all{haben} en position 2, participe \all{gewählt} en fin ; « l'année dernière » en tête fait passer le sujet après le verbe.)
\item \all{Obwohl die Arbeit schwer ist, helfen die Freiwilligen den älteren Menschen.} (\all{obwohl} : verbe en fin ; \all{helfen} + datif : \all{den älteren Menschen} ; quand la subordonnée ouvre la phrase, le verbe de la principale suit immédiatement : \all{..., helfen...}.)
\item \all{In unserem Viertel nehmen wir mit unseren Nachbarn an einem sozialen Projekt teil.} (Verbe à particule \all{teilnehmen} : \all{teil} en fin ; \all{an} + datif ; \all{mit} + datif pluriel \all{unseren Nachbarn}.)
\end{enumerate}
\end{corrige}

\begin{exercice}[Version : le texte sur le bénévolat]
Traduisez en français le texte \all{Jugendliche engagieren sich} dans son intégralité. Relisez ensuite \emph{uniquement} votre français : aucune phrase ne doit « sonner allemand ».
\end{exercice}

\begin{corrige}[Exercice : version (proposition de traduction)]
« En Allemagne, de nombreux jeunes s'engagent bénévolement. Ils travaillent par exemple dans des clubs de sport, chez les pompiers ou dans des groupes écologistes. À l'école, beaucoup d'élèves ont la chance de participer à des projets sociaux : ils aident les personnes âgées, collectent de l'argent pour les enfants ou plantent des arbres dans leur ville. „L'année dernière, j'ai construit une aire de jeux avec ma classe“, raconte Lena, 17 ans, de Hambourg. „C'était difficile, mais nous avons beaucoup appris.“ Les hommes politiques répètent souvent que le bénévolat est important pour la démocratie, parce que les citoyens y apprennent la responsabilité. Au Cameroun aussi, des jeunes s'engagent : ils organisent des tournois de football ou nettoient les rues de leur quartier. L'engagement relie les hommes — partout dans le monde. »
\end{corrige}

\begin{aretenir}[Bilan de la section]
\begin{itemize}
\item \textbf{Thème} : verbe d'abord (place !), cas ensuite, majuscules et Umlaute à la relecture.
\item \textbf{Version} : le sens, jamais le mot à mot ; le français final doit être naturel.
\item \textbf{Pièges} : \all{bekommen} = recevoir ; verbe en fin après \all{weil}, \all{dass}, \all{obwohl} ; prépositions à cas fixe ; Präteritum $\neq$ un seul temps français.
\end{itemize}
Vous disposez désormais de toutes les clés pour réussir l'épreuve de traduction sur le thème de la citoyenneté. \all{Viel Erfolg!}
\end{aretenir}

\newpage

\coursec{Activité d'intégration}

\coursec{Leçon AL30 : Traduction et production de textes divers liés à la citoyenneté}

\courssub{Objectifs de l'activité}

À l'issue de cette activité d'intégration, tu seras capable de :
\begin{itemize}
\item rédiger en allemand un texte court et authentique (article ou lettre ouverte) sur un acte de \cle{citoyenneté} réalisé dans ton lycée ;
\item traduire en français un texte allemand reçu d'un autre groupe (travail de \cle{version}) ;
\item réinvestir le lexique du chapitre (\all{die Staatsbürgerschaft, die Wahl, die Solidarität, die Freiwilligenarbeit}…) et les structures grammaticales révisées (ordre des mots, cas, temps du passé, voix passive, discours indirect simple) ;
\item relire et corriger un texte à l'aide d'une grille de correction (majuscules des noms, place du verbe, cas, orthographe avec \all{ä, ö, ü, ß}) ;
\item coopérer en groupe dans une tâche de médiation linguistique proche de la réalité professionnelle.
\end{itemize}

\courssub{Situation}

Ton lycée de Yaoundé est partenaire, depuis trois ans, du \all{Gymnasium am Stadtpark} de Hambourg. Chaque trimestre, les deux écoles échangent un journal scolaire bilingue intitulé \all{„Brücke der Jugend“} („ Le pont de la jeunesse “). Le club d'allemand de ton lycée est chargé du prochain numéro, dont le dossier central porte sur la citoyenneté au quotidien. La professeure coordinatrice allemande a envoyé le courriel suivant :

\begin{textebox}[Courriel de la partenaire allemande]
Liebe Freunde in Yaoundé,\\
für die nächste Ausgabe unserer Schülerzeitung „Brücke der Jugend“ brauchen wir eure Beiträge zum Thema „Junge Bürger handeln“. Schreibt uns einen kurzen Artikel oder einen offenen Brief über eine konkrete Aktion der Zivilcourage oder der Solidarität in eurer Schule: eine Wahl, eine Putzaktion im Viertel, eine Hilfsaktion für Bedürftige. Wir möchten wissen: Was habt ihr gemacht? Wer hat teilgenommen? Warum war diese Aktion wichtig für eure Gemeinschaft?\\
Wir freuen uns auf eure Texte!\\
Herzliche Grüße\\
Frau Berger, Koordinatorin
\end{textebox}

\begin{definition}[Glossaire du courriel]
\all{die Ausgabe, -n} : le numéro (d'une revue) ; \all{der Beitrag, -e} : l'article, la contribution ; \all{die Zivilcourage} (n. f., sans pluriel) : le courage civique ; \all{die Putzaktion, -en} : l'opération de nettoyage ; \all{die Hilfsaktion, -en} : l'opération d'aide ; \all{bedürftig} (adj.) : nécessiteux ; \all{teilnehmen (nahm teil, hat teilgenommen)} : participer ; \all{die Gemeinschaft, -en} : la communauté.
\end{definition}

Le proviseur a accepté que les meilleurs textes soient affichés au tableau du club d'allemand et envoyés à Hambourg. Ta classe se met donc au travail : par groupes de quatre, vous formez la rédaction camerounaise du journal.

\courssub{Consignes}

\begin{enumerate}
\item \textbf{Préparation collective (10 min).} Dans ton groupe, choisissez un acte de citoyenneté réellement réalisé (ou réaliste) dans votre lycée : élection des délégués de classe, nettoyage du quartier avec les jeunes du voisinage, collecte de vivres pour un orphelinat, campagne contre les déchets plastiques. Notez en allemand cinq mots-clés du lexique de la citoyenneté que vous voulez absolument employer.
\item \textbf{Production écrite (25 min).} Rédigez en allemand un article ou une lettre ouverte de 10 à 12 lignes. Votre texte doit contenir : un titre, une présentation de l'action (qui ? quoi ? où ? quand ?), au moins une phrase au \cle{Perfekt} et une au \cle{Präteritum}, une phrase à la voix passive, une citation au discours indirect (\all{Er sagte, dass…}) et une conclusion sur l'importance de l'action pour la communauté.
\item \textbf{Échange et version (20 min).} Échangez votre texte avec celui d'un autre groupe. Traduisez en français le texte reçu, en phrases complètes et naturelles : c'est un travail de \cle{version}, pas de mot à mot.
\item \textbf{Relecture croisée (15 min).} À l'aide de la grille de correction ci-dessous, corrigez le texte allemand que vous venez de traduire : majuscules des noms, place du verbe (V2 en proposition principale, verbe en fin de subordonnée), cas après les prépositions, particules séparables, orthographe (\all{ä, ö, ü, ß}). Renvoyez le texte corrigé à ses auteurs avec vos remarques.
\item \textbf{Finalisation (10 min).} Chaque groupe intègre les corrections, réécrit la version définitive et la remet pour l'affichage. La classe vote pour les deux meilleurs textes qui partiront à Hambourg.
\end{enumerate}

\begin{methode}[Grille de correction pour la relecture croisée]
\begin{enumerate}
\item \textbf{Majuscules} : tous les noms communs allemands portent-ils une majuscule (\all{die Wahl, der Müll, die Schüler}) ?
\item \textbf{Place du verbe} : le verbe conjugué est-il en deuxième position dans les phrases déclaratives ? En fin de proposition après \all{dass, weil, als} ?
\item \textbf{Particules séparables} : \all{teilnehmen} donne-t-il bien \all{Die Schüler nahmen teil} ?
\item \textbf{Cas} : \all{mit dem Viertel}, \all{für die Bedürftigen}, \all{nach der Aktion} — les déterminants sont-ils corrects ?
\item \textbf{Temps} : le récit au passé utilise-t-il le Perfekt (récit) et le Präteritum (\all{war, hatte, gab}) de façon cohérente ?
\item \textbf{Orthographe} : Umlaute et \all{ß} sont-ils présents et bien placés ?
\end{enumerate}
\end{methode}

\courssub{Ressources à mobiliser}

\begin{itemize}
\item \textbf{Lexique de la citoyenneté} : \all{der Bürger, - / die Bürgerin, -nen} (le/la citoyen·ne) ; \all{die Wahl, -en} (l'élection) ; \all{wählen} (élire, voter) ; \all{der Delegierte, -n} (le délégué) ; \all{die Solidarität} (la solidarité) ; \all{die Freiwilligenarbeit} (le bénévolat) ; \all{sich engagieren für + Akk.} (s'engager pour) ; \all{die Verantwortung} (la responsabilité) ; \all{der Müll} (les déchets) ; \all{spenden} (faire un don) ; \all{gemeinsam} (ensemble).
\item \textbf{Ordre des mots} : verbe conjugué en position 2 ; si un complément ouvre la phrase, le sujet passe après le verbe : \all{Am Samstag haben wir das Viertel gereinigt.}
\item \textbf{Passé} : Perfekt pour le récit (\all{wir haben gesammelt, wir sind gegangen}) ; Präteritum pour \all{sein} et \all{haben} (\all{es war, wir hatten}) ainsi que pour les verbes modaux.
\item \textbf{Voix passive (Vorgangspassiv)} : \all{Es wurden zwanzig Säcke Müll gesammelt.} („ Vingt sacs de déchets ont été ramassés. “) — auxiliaire \all{werden} + participe passé.
\item \textbf{Discours indirect} : \all{Der Delegierte sagte, dass alle Schüler mitgemacht hätten.} (Konjunktiv II \all{hätten} + participe passé \all{mitgemacht} en fin de subordonnée).
\item \textbf{Subordonnées} : verbe conjugué en fin après \all{weil, dass, als, damit}.
\end{itemize}

\begin{attention}[Pièges fréquents dans cette tâche]
Ne traduis jamais « nettoyer le quartier » par \all{putzen das Viertel} seul : préfère \all{das Viertel reinigen} ou \all{sauber machen}. Attention au faux ami \all{aktiv} : « actif » se dit bien \all{aktiv}, mais « un acte » se dit \all{die Handlung} ou \all{die Aktion}, jamais \all{der Akt} dans ce contexte. Enfin, « s'engager » exige le réfléchi : \all{Wir haben uns für die Ärmsten engagiert.} — sans \all{sich}, la phrase est fautive.
\end{attention}

\courssub{Proposition de production et corrigé commenté}

\begin{exoresolu}[Article modèle pour „Brücke der Jugend“]
Rédige un article de 10 à 12 lignes sur une opération de nettoyage organisée par les délégués de Terminale A dans le quartier du lycée, en respectant toutes les contraintes de la consigne 2.
\tcblower
\textbf{„Sauberes Viertel, starke Gemeinschaft“}\\
Am Samstagmorgen haben die Schüler der Terminale A unseres Gymnasiums in Yaoundé eine große Putzaktion im Viertel organisiert. Über fünfzig Jugendliche nahmen daran teil, darunter auch den gewählten Klassensprechern. Zuerst wurden die Straßen um die Schule gefegt, dann sammelten die Gruppen den Plastikmüll ein. Es war anstrengend, aber die Stimmung war ausgezeichnet. Der Klassensprecher erklärte, dass diese Aktion allen Bewohnern ein Beispiel geben sollte. Zum Schluss wurden zwanzig Säcke Müll zur Müllabfuhr gebracht. Wir haben gelernt, dass Zivilcourage mit kleinen Gesten beginnen. Diese Aktion war wichtig für unsere Gemeinschaft, weil ein sauberes Viertel die Gesundheit aller schützt. Nächsten Monat wollen wir eine Hilfsaktion für bedürftige Familien starten.\\
\emph{Traduction : « Un quartier propre, une communauté forte » — Samedi matin, les élèves de Terminale A de notre lycée de Yaoundé ont organisé une grande opération de nettoyage dans le quartier. Plus de cinquante jeunes y ont participé, dont les délégués élus. D'abord, les rues autour de l'école ont été balayées, puis les groupes ont ramassé les déchets plastiques. C'était fatigant, mais l'ambiance était excellente. Le délégué a expliqué que cette action devait donner l'exemple à tous les habitants. Pour finir, vingt sacs de déchets ont été apportés à la collecte. Nous avons appris que le courage civique commence par de petits gestes. Cette action était importante pour notre communauté parce qu'un quartier propre protège la santé de tous. Le mois prochain, nous voulons lancer une opération d'aide aux familles nécessiteuses.}

\textbf{Commentaire des choix de langue :}
\begin{itemize}
\item \textbf{Ouverture par un complément} : \all{Am Samstagmorgen haben die Schüler…} — inversion sujet/verbe correcte après le complément initial (structure V2).
\item \textbf{Particule séparable} : \all{nahmen daran teil} (\all{teilnehmen an + Dat.}) — la particule \all{teil} va en fin de proposition principale.
\item \textbf{Préposition \all{darunter} + Datif} : \all{darunter auch den gewählten Klassensprechern} — \all{darunter} (parmi eux) gouverne le datif pluriel (\all{den Klassensprechern}).
\item \textbf{Voix passive au Präteritum} : \all{wurden die Straßen gefegt}, \all{wurden zwanzig Säcke gebracht} — structure \all{werden} (Präteritum) + Partizip II ; met l'action en avant, typique du style journalistique.
\item \textbf{Discours indirect} : \all{Der Klassensprecher erklärte, dass diese Aktion… ein Beispiel geben sollte.} — Konjunktiv II de renvoi (\all{sollte}) pour le verbe modal \all{sollen}, verbe en fin de subordonnée.
\item \textbf{Konjunktiv I dans le discours indirect} : \all{dass Zivilcourage mit kleinen Gesten beginnen.} — après \all{lernen, dass}, le Konjunktiv I (\all{beginnen}, forme = infinitif) marque le rapport de fait/opinion.
\item \textbf{Subordonnée de cause} : \all{weil ein sauberes Viertel die Gesundheit aller schützt} — verbe conjugué en dernière position.
\item \textbf{Lexique du chapitre réinvesti} : \all{die Gemeinschaft, die Zivilcourage, die Hilfsaktion, bedürftig, gewählt, die Putzaktion, der Klassensprecher}.
\end{itemize}
\end{exoresolu}

\courssub{Bilan de l'activité}

Cette correspondance citoyenne t'a permis de mobiliser simultanément les quatre savoir-faire du chapitre : rédiger un texte codifié (article ou lettre ouverte), traduire de l'allemand vers le français, réinvestir la grammaire révisée (ordre des mots, cas, passif, discours indirect) et relire une production avec une grille objective. Tu as aussi découvert la dimension concrète de la médiation linguistique : traduire, c'est d'abord comprendre le sens global avant de chercher les mots. Pour t'évaluer, pose-toi trois questions : mon texte allemand respecte-t-il la place du verbe dans chaque phrase ? Ai-je employé au moins cinq mots du lexique de la citoyenneté ? Ma version française est-elle naturelle et fidèle au sens ? Si tu réponds oui aux trois, tu es prêt pour l'épreuve de production et de traduction du baccalauréat — et ton texte a peut-être déjà pris le chemin de Hambourg.

\newpage

\coursec{Méthodes et exercices résolus}

\courssub{Révision grammaticale ciblée pour la production et la traduction}

\begin{aretenir}[Structures clés pour écrire et traduire sur la citoyenneté]
\begin{itemize}
\item \textbf{Ordre des mots (TeKaMoLo)} : Sujet -- Verbe (pos. 2) -- Temps -- Cause -- Manière -- Lieu -- Infinitif/Participe (fin). En subordonnée (\all{weil, dass, wenn, obwohl}) : verbe conjugué \textbf{à la fin}.
\item \textbf{Cas prépositionnels et verbaux} : \all{helfen} + \textbf{Datif} ; \all{teilnehmen an} + \textbf{Datif} ; \all{sich interessieren für} + \textbf{Accusatif} ; \all{kämpfen für} + \textbf{Accusatif} ; \all{fordern, dass} + subordonnée ; \all{bitten} + \textbf{Akkusativ} (personne) + \textbf{zu}-infinitif.
\item \textbf{Konjunktiv II (irréel / politesse)} : \all{würde} + infinitif (tous verbes) OU forme pleine (\all{hätte, wäre, ginge, käme, wüsste, wollte, dürfte, müsste, könnte, sollte}). Si-clause : \all{Wenn ... würde ..., würde ...}.
\item \textbf{Perfekt (passé oral)} : Auxiliaire \all{haben} (verbes transitifs, réflexifs, modaux) ou \all{sein} (mouvement, changement d'état, \all{sein, bleiben, werden}). Participe : verbes séparables \all{an-ge-fangen}, inséparables \all{teilgenommen} (pas de \all{ge-}), verbes en \all{-ieren} \all{organisiert}.
\item \textbf{Impératif} : \all{du} : radical (+ \all{-e} optionnel) ; \all{ihr} : forme \all{ihr} sans pronom ; \all{Sie} : forme \all{Sie} avec pronom après verbe. Séparables : préfixe à la fin (\all{Fangt an!}).
\item \textbf{Nominalisation} : \all{das Wählen, das Engagement, das Helfen} (neutre, majuscule). Prépositions : \all{beim} (\all{bei dem}), \all{zum} (\all{zu dem}), \all{für das} (\all{fürs} à l'oral).
\item \textbf{Discours indirect (Konjunktiv I)} : \all{Er sagte, er \textbf{sei} stolz} (3e pers. sg. \all{sei}, pas \all{ist}). Au passé : \all{er \textbf{habe} gemacht}.
\end{itemize}
\end{aretenir}

\courssub{Méthode 1 : Rédiger un discours (\all{eine Rede}) sur la citoyenneté}

\begin{methode}[Réussir un discours en cinq étapes]
\begin{enumerate}
\item \textbf{Analyser la consigne} : qui parle ? à qui ? à quelle occasion ? (fête nationale, élection, cérémonie scolaire). Cela détermine le registre : \all{Liebe Mitbürgerinnen und Mitbürger}, \all{Sehr geehrte Damen und Herren}, \all{Liebe Schülerinnen und Schüler}.
\item \textbf{Construire le plan en trois parties} : accroche (question rhétorique ou constat), développement (2 ou 3 arguments avec exemples concrets), appel à l'action (\all{Handeln wir gemeinsam!}).
\item \textbf{Utiliser les connecteurs de l'argumentation} : \all{zuerst} (d'abord), \all{außerdem} (de plus), \all{jedoch} (cependant), \all{deshalb} (c'est pourquoi), \all{zusammenfassend} (pour résumer).
\item \textbf{Soigner les structures fortes} : impératif (\all{Engagiert euch!}), \all{wir}-formes inclusives, phrases avec \all{es ist wichtig, dass ...} suivi du subordonné (verbe à la fin).
\item \textbf{Relire} : majuscule à tous les noms, verbe en position 2 dans les principales, verbe à la fin dans les subordonnées, accord des cas après prépositions (\all{für die Demokratie}, accusatif).
\end{enumerate}
\end{methode}

\begin{exoresolu}[Discours : la participation des jeunes]
\textbf{Énoncé (type Bac)} : Vous êtes le président du club d'allemand de votre lycée à Yaoundé. Rédigez en allemand un discours de 10 à 12 lignes pour la journée de la citoyenneté, afin d'encourager les élèves à s'engager dans la vie de leur quartier.
\tcblower
\textbf{Raisonnement} : registre = pairs, donc \all{Liebe Mitschülerinnen und Mitschüler}. Plan : accroche par une question, deux arguments (le quartier a besoin de nous ; l'engagement forme le citoyen), appel à l'action à l'impératif.

\textbf{Solution rédigée} :

\all{Liebe Mitschülerinnen und Mitschüler,}

\all{Was bedeutet es, ein guter Bürger zu sein? Zuerst müssen wir verstehen, dass die \cle{Bürgerschaft} nicht nur ein Wort ist, sondern eine Aufgabe. Unser Viertel braucht uns: Wir können den Müll sammeln, den jüngeren Kindern bei den Hausaufgaben helfen und an den Sitzungen des Quartierskomitees teilnehmen.}

\all{Außerdem lernen wir durch das Engagement Verantwortung. Wer heute für die Gemeinschaft arbeitet, wird morgen ein verantwortungsbewusster Wähler sein. Deshalb ist es wichtig, dass jeder von uns etwas tut.}

\all{Zusammenfassend sage ich: Wartet nicht auf die anderen! Engagiert euch jetzt, denn die Zukunft unseres Landes beginnt in unserem Viertel. Vielen Dank für eure Aufmerksamkeit.}

\textbf{Justifications grammaticales} : \all{ein guter Bürger} (adjectif épithète après article indéfini, nominatif masculin : \all{-er}) ; \all{dass die \cle{Bürgerschaft} ... eine Aufgabe ist} (subordonnée avec \all{dass}, verbe \all{ist} à la fin) ; \all{den jüngeren Kindern} (datif pluriel après \all{helfen} : \all{jemandem helfen}) ; \all{Wer heute ... arbeitet, wird morgen ... sein} (verbe en position 2 après la subordonnée qui occupe la position 1) ; \all{Engagiert euch} (impératif pluriel du verbe réfléchi \all{sich engagieren}) ; \all{Quartierskomitees} (nom composé avec \cle{Fugen-s}).

\textbf{Conclusion} : le discours respecte le plan accroche–arguments–appel, emploie quatre connecteurs et des structures variées : c'est ce que le correcteur valorise.
\end{exoresolu}

\courssub{Méthode 2 : Rédiger une lettre ouverte (\all{einen offenen Brief})}

\begin{methode}[Réussir une lettre ouverte en cinq étapes]
\begin{enumerate}
\item \textbf{Identifier le destinataire} (maire, ministre, rédacteur en chef) et choisir la formule d'appel : \all{Sehr geehrter Herr Bürgermeister}, \all{Sehr geehrte Damen und Herren}.
\item \textbf{Respecter la mise en page allemande} : lieu et date en haut à droite (\all{Yaoundé, den 15. Mai 2025}), objet possible (\all{Betreff: ...}), virgule après l'appel, première ligne du corps en minuscule.
\item \textbf{Structurer} : rappel du problème, faits et exemples, revendication claire (\all{Wir fordern, dass ...} + verbe à la fin), proposition de solution.
\item \textbf{Employer le registre formel} : \all{Sie} avec majuscule (et \all{Ihnen}, \all{Ihre}), tournures polies (\all{Wir bitten Sie, ...}), pas d'abréviations.
\item \textbf{Conclure} : formule de politesse \all{Mit freundlichen Grüßen} (sans virgule après !), puis signature.
\end{enumerate}
\end{methode}

\begin{exoresolu}[Lettre ouverte au maire]
\textbf{Énoncé (type Bac)} : Les jeunes de votre quartier à Douala n'ont aucun terrain de sport. Rédigez en allemand une lettre ouverte au maire (8 à 10 lignes) pour demander la construction d'un terrain.
\tcblower
\textbf{Raisonnement} : destinataire = le maire, donc \all{Sehr geehrter Herr Bürgermeister}. L'adjectif \all{geehrter} porte la désinence forte \all{-er} (nominatif masculin sans article : vocatif/appellation). Revendication avec \all{wir fordern, dass ...}.

\textbf{Solution rédigée} :

\all{Douala, den 15. Mai 2025}

\all{Sehr geehrter Herr Bürgermeister,}

\all{wir, die Jugendlichen des Quartiers Bonabéri, schreiben Ihnen, weil wir ein großes Problem haben: Es gibt bei uns keinen Sportplatz. Viele junge Leute spielen deshalb auf der Straße, und das ist gefährlich.}

\all{Sport ist aber wichtig für die Gesundheit und für die Erziehung der Bürger. Ein Sportplatz würde auch die Jugendkriminalität reduzieren, weil die Jugendlichen dann eine sinnvolle Beschäftigung hätten.}

\all{Wir fordern, dass die Stadt einen Sportplatz in unserem Quartier baut. Wir sind bereit, bei der Arbeit zu helfen.}

\all{Mit freundlichen Grüßen}

\all{Die Jugend von Bonabéri}

\textbf{Justifications grammaticales} : \all{keinen Sportplatz} (accusatif après \all{es gibt}) ; \all{weil wir ein großes Problem haben} (verbe \all{haben} à la fin après \all{weil}) ; \all{würde ... reduzieren} et \all{hätten} (Konjunktiv II pour l'hypothèse : \all{würde} + infinitif, \all{hätten} forme pleine de \all{haben}) ; \all{bei der Arbeit} (datif féminin après \all{bei}) ; \all{dass die Stadt ... baut} (verbe à la fin) ; \all{in unserem Quartier} (datif après \all{in} + lieu).

\textbf{Conclusion} : la lettre suit la structure problème–arguments–revendication–politesse ; le ton reste respectueux mais ferme, comme l'exige le genre de la lettre ouverte.
\end{exoresolu}

\courssub{Méthode 3 : Rédiger une interview et un article}

\begin{methode}[Réussir interview et article en cinq étapes]
\begin{enumerate}
\item \textbf{Interview} : préparer 4 à 6 questions variées — questions ouvertes avec \all{W-Fragen} (\all{Was, Warum, Wie, Seit wann}) et questions fermées avec inversion (\all{Glauben Sie, dass ...?}). Indiquer les intervenants : \all{Interviewer:} / \all{Bürgermeister:}.
\item \textbf{Alterner les types de questions} et faire des réponses de 2 à 3 phrases, avec des exemples concrets.
\item \textbf{Article} : titre court, \all{Einleitung} (qui ? quoi ? où ? quand ?), corps avec faits et citations au discours indirect (Konjunktiv I : \all{Der Bürgermeister sagte, er \textbf{sei} stolz auf die Jugend.}), conclusion brève.
\item \textbf{Employer les verbes de parole} : \all{sagen, erklären, meinen, betonen, hinzufügen} (ajouter).
\item \textbf{Relire} : inversion après \all{W-Fragen} et verbes de déclaration, majuscules des noms, cohérence des temps.
\end{enumerate}
\end{methode}

\begin{exoresolu}[Interview avec un jeune volontaire]
\textbf{Énoncé (type Bac)} : Vous êtes journaliste pour le journal du lycée. Rédigez en allemand une interview (4 questions minimum) d'un élève qui fait du bénévolat dans son quartier.
\tcblower
\textbf{Raisonnement} : questions ouvertes variées (\all{Seit wann}, \all{Warum}, \all{Was}, \all{Wie}) ; réponses au Perfekt ou au Präsens ; registre \all{du} ou \all{Sie} — ici entre jeunes, on choisit \all{du} de façon cohérente.

\textbf{Solution rédigée} :

\all{Interviewer: Seit wann engagierst du dich in deinem Viertel?}

\all{Paul: Seit zwei Jahren. Ich habe mit einer Müllsammelaktion angefangen.}

\all{Interviewer: Warum hast du dich entschieden, freiwillig zu arbeiten?}

\all{Paul: Weil ich etwas Nützliches für meine Gemeinschaft tun wollte. Ein guter Bürger wartet nicht; er handelt.}

\all{Interviewer: Was machst du genau?}

\all{Paul: Ich helfe den kleinen Kindern bei den Hausaufgaben und ich organisiere Fußballspiele am Wochenende.}

\all{Interviewer: Was würdest du den anderen Jugendlichen empfehlen?}

\all{Paul: Fangt an, auch wenn es klein ist! Jede kleine Aktion verändert das Viertel.}

\textbf{Justifications grammaticales} : \all{Seit wann engagierst du dich} (inversion après le mot interrogatif, verbe réfléchi) ; \all{Seit zwei Jahren} + Präsens (action qui continue : en français « depuis deux ans ») ; \all{Ich habe ... angefangen} (Perfekt du verbe à particule séparable : participe \all{angefangen} à la fin) ; \all{den kleinen Kindern} (datif pluriel après \all{helfen}) ; \all{Was würdest du ... empfehlen} (Konjunktiv II de politesse) ; \all{Fangt an} (impératif pluriel, préfixe séparable renvoyé à la fin) ; \all{auch wenn es klein ist} (concession, verbe à la fin).

\textbf{Conclusion} : l'interview alterne questions ouvertes et variées, les réponses donnent des détails concrets, et la grammaire (inversion, particules séparables, datif) est maîtrisée.
\end{exoresolu}

\courssub{Méthode 4 : Traduire du français vers l'allemand (thème)}

\begin{methode}[Réussir le thème en cinq étapes]
\begin{enumerate}
\item \textbf{Lire toute la phrase française} et repérer le verbe principal, son temps et son sujet avant d'écrire un seul mot.
\item \textbf{Traduire la structure, pas les mots} : construire d'abord le squelette allemand (sujet + verbe en position 2), puis placer les compléments selon la règle TeKaMoLo (temps, cause, manière, lieu).
\item \textbf{Choisir le bon temps} : passé composé français de conversation = \all{Perfekt} (\all{ich habe gemacht}) ; récit écrit = \all{Präteritum} ; futur = \all{werden} + infinitif.
\item \textbf{Vérifier chaque nom} : genre (article !), cas (fonction dans la phrase), majuscule.
\item \textbf{Relire en chassant les calques} : « être d'accord » = \all{zustimmen} ou \all{einverstanden sein} (jamais \all{*sein akkord}) ; « s'intéresser à » = \all{sich interessieren für} + accusatif ; « participer » = \all{teilnehmen an} + datif (pas \all{*partizipieren}).
\end{enumerate}
\end{methode}

\begin{exoresolu}[Thème : phrases sur la citoyenneté]
\textbf{Énoncé (type Bac)} : Traduisez en allemand : 1. « Les citoyens doivent respecter les lois de leur pays. » 2. « Hier, nous avons participé à une réunion au quartier. » 3. « Si les jeunes votaient plus, la politique changerait. »
\tcblower
\textbf{Phrase 1} : \all{Die Bürger müssen die Gesetze ihres Landes respektieren.}
Raisonnement : « doivent » = modal \all{müssen} en position 2, infinitif \all{respektieren} à la fin ; « les lois » = accusatif pluriel (\all{die Gesetze}) ; « de leur pays » = génitif : \all{ihres Landes} (génitif neutre de \all{das Land}, avec \all{-es} au nom).

\textbf{Phrase 2} : \all{Gestern haben wir an einer Sitzung im Viertel teilgenommen.}
Raisonnement : passé composé oral = Perfekt : auxiliaire \all{haben} en position 2, participe \all{teilgenommen} à la fin (verbe à particule inséparable \all{teilnehmen} : pas de \all{ge-} !) ; « participer à » = \all{an etwas teilnehmen} + datif : \all{an einer Sitzung} ; \all{im Viertel} = \all{in dem Viertel} (datif, lieu).

\textbf{Phrase 3} : \all{Wenn die jungen Leute mehr wählen würden, würde sich die Politik ändern.}
Raisonnement : hypothèse irréelle = Konjunktiv II : \all{würden wählen} dans la subordonnée (verbe à la fin), \all{würde sich ändern} dans la principale avec le verbe en position 2 (la subordonnée occupe la position 1) ; « changer » ici = changement en soi, donc réfléchi \all{sich ändern}.

\textbf{Conclusion} : chaque phrase a été construite sur le squelette allemand avant d'ajouter les compléments : aucune n'est un calque mot à mot du français.
\end{exoresolu}

\courssub{Méthode 5 : Traduire de l'allemand vers le français (version)}

\begin{methode}[Réussir la version en cinq étapes]
\begin{enumerate}
\item \textbf{Lire le texte entier} une première fois sans traduire pour saisir le sens global et le thème.
\item \textbf{Découper chaque phrase} : repérer le verbe conjugué (position 2), le sujet, puis les éléments entre le verbe et la fin de phrase (participe, infinitif, préfixe séparable).
\item \textbf{Traduire les temps avec exactitude} : \all{Präteritum} = passé simple ou imparfait selon le contexte ; \all{Perfekt} = passé composé ; Konjunktiv II = conditionnel français.
\item \textbf{Rendre un français naturel} : ne pas garder l'ordre allemand ; reformuler (\all{Es ist wichtig, dass ...} = « il est important que » + subjonctif).
\item \textbf{Relire la traduction seule} : elle doit se comprendre sans le texte allemand, sans mot laissé de côté.
\end{enumerate}
\end{methode}

\begin{exoresolu}[Version : un court texte sur le vote]
\textbf{Énoncé (type Bac)} : Traduisez en français : \all{Viele junge Menschen gehen nicht wählen, weil sie denken, dass die Politik ihre Probleme nicht versteht. Dabei ist das Wählen ein wichtiges Recht, für das viele Menschen früher gekämpft haben. Wer nicht wählt, überlässt die Entscheidung den anderen.}
\tcblower
\textbf{Analyse phrase par phrase} :

Phrase 1 : verbe \all{gehen} (position 2), infinitif \all{wählen} à la fin (\all{wählen gehen} = „aller voter“) ; subordonnée \all{weil sie denken} (verbe à la fin) ; subordonnée imbriquée \all{dass die Politik ihre Probleme nicht versteht}.

Phrase 2 : \all{Dabei} = „pourtant / or“ ; \all{für das viele Menschen früher gekämpft haben} = relative avec préposition : « pour lequel beaucoup de gens ont lutté autrefois » (Perfekt = passé composé). \all{das Wählen} = nominalisation (neutre).

Phrase 3 : \all{Wer nicht wählt} = proposition indéfinie : « celui qui ne vote pas » ; \all{überlässt} (verbe à particule inséparable, de \all{überlassen} + datif) = « laisse » ; \all{den anderen} = datif pluriel : « aux autres ».

\textbf{Traduction rédigée} : « Beaucoup de jeunes ne vont pas voter, parce qu'ils pensent que la politique ne comprend pas leurs problèmes. Pourtant, voter est un droit important pour lequel beaucoup de gens ont lutté autrefois. Celui qui ne vote pas laisse la décision aux autres. »

\textbf{Conclusion} : la version exige de repérer d'abord les verbes et les constructions (relative avec préposition, \all{wer}-proposition), puis de reformuler en français idiomatique.
\end{exoresolu}

\begin{attention}[Les 5 erreurs qui coûtent des points au Bac]
\begin{enumerate}
\item \textbf{La place du verbe} : après \all{weil}, \all{dass}, \all{wenn}, \all{obwohl}, le verbe conjugué va \textbf{à la fin}. Écrire \all{*weil sie denken die Politik...} est une faute éliminatoire. Et quand la subordonnée ouvre la phrase, le verbe de la principale suit immédiatement : \all{Wer nicht wählt, überlässt ...} (jamais \all{*Wer nicht wählt, er überlässt...}).
\item \textbf{Les majuscules des noms} : tout nom allemand prend une majuscule : \all{das Recht, die Entscheidung, der Bürger}. Une minuscule sur un nom est sanctionnée à chaque occurrence. Attention aussi à \all{Sie} / \all{Ihnen} (vouvoiement) toujours en majuscule.
\item \textbf{Les cas après verbes et prépositions} : \all{helfen} + datif (\all{ich helfe den Kindern}, pas \all{*die Kinder}) ; \all{teilnehmen an} + datif ; \all{sich interessieren für} + accusatif ; \all{für das Recht kämpfen} (accusatif). Apprendre chaque verbe avec sa construction.
\item \textbf{Les faux amis et calques du français} : \all{aktual} n'existe pas (adj. \all{aktuell} = actuel) ; « participer » = \all{teilnehmen}, pas \all{*partizipieren} au lycée ; « une loi » = \all{das Gesetz}, pas \all{*die Loi} ; « être d'accord » = \all{einverstanden sein} / \all{zustimmen}, jamais de calque avec \all{akkord}.
\item \textbf{Le Perfekt des verbes à particule} : séparables : \all{ich habe angefangen} (le \all{-ge-} s'insère : \all{an-ge-fangen}) ; inséparables : \all{ich habe teilgenommen} (pas de \all{ge-}). Mélanger les deux (\all{*geteilnahmen}, \all{*anfangt}) coûte cher en thème comme en rédaction.
\end{enumerate}
\end{attention}

\newpage

\coursec{Exercices}

\courssub{Je vérifie mes connaissances}

\begin{exercice}[QCM : vocabulaire de la citoyenneté]
Entoure la bonne réponse.
\begin{enumerate}
\item \all{die Staatsbürgerschaft} signifie :
\begin{enumerate}
\item la nationalité / la citoyenneté \item la mairie \item le passeport
\end{enumerate}
\item \all{wählen} signifie :
\begin{enumerate}
\item voter / élire \item choisir un métier \item se plaindre
\end{enumerate}
\item Le contraire de \all{die Pflicht} (le devoir) est :
\begin{enumerate}
\item \all{das Recht} \item \all{die Freiwilligkeit} \item \all{die Wahl}
\end{enumerate}
\item \all{sich engagieren} se traduit par :
\begin{enumerate}
\item se marier \item s'engager, s'investir \item se fâcher
\end{enumerate}
\item \all{der Bürgermeister} est :
\begin{enumerate}
\item le citoyen \item le maire \item le député
\end{enumerate}
\end{enumerate}
\end{exercice}

\begin{exercice}[Vrai ou faux ? Justifie ta réponse]
Dis si les affirmations suivantes sont vraies ou fausses et justifie en une phrase.
\begin{enumerate}
\item En allemand, le verbe conjugué se trouve toujours en deuxième position dans la phrase principale.
\item Dans une subordonnée introduite par \all{weil}, le verbe conjugué reste en deuxième position.
\item Le passif se forme avec l'auxiliaire \all{haben} + Partizip II.
\item \all{damit} introduit une subordonnée de but et renvoie le verbe conjugué à la fin.
\item Tous les noms allemands prennent une majuscule, même au milieu de la phrase.
\end{enumerate}
\end{exercice}

\begin{exercice}[Lexique : trouve le mot allemand]
Donne le mot allemand (avec l'article et le pluriel pour les noms) correspondant à :
\begin{enumerate}
\item le vote, l'élection
\item le bénévolat, le volontariat
\item la responsabilité
\item la société
\item les droits de l'homme
\item la commune, la municipalité
\end{enumerate}
\end{exercice}

\begin{exercice}[Conjugaison : à compléter]
Complète les phrases avec la forme correcte du verbe entre parenthèses au temps indiqué.
\begin{enumerate}
\item Die Jugendlichen \all{\ldots} (sich engagieren, Präsens) in einem sozialen Projekt.
\item Gestern \all{\ldots} (haben, Präteritum) wir eine Sitzung im Gemeindehaus.
\item Der Bürgermeister \all{\ldots} (versprechen, Perfekt) mehr Mittel für die Schule.
\item Wenn ich mehr Zeit hätte, \all{\ldots} (helfen, Konjunktiv II) ich den alten Menschen.
\item Nächstes Jahr \all{\ldots} (wählen, Futur I) viele junge Leute zum ersten Mal.
\end{enumerate}
\end{exercice}

\courssub{Je m'entraîne}

\begin{exercice}[Thème : phrases à traduire]
Traduis les phrases suivantes en allemand.
\begin{enumerate}
\item Les jeunes devraient s'engager davantage parce que l'avenir du pays dépend d'eux.
\item Chaque citoyen a des droits, mais aussi des devoirs.
\item Si j'étais maire, je construirais une bibliothèque pour les élèves.
\item Nous avons planté des arbres afin de protéger l'environnement.
\item Le vote est un droit important dans une démocratie.
\end{enumerate}
\end{exercice}

\begin{exercice}[Transformation : actif et passif]
Réécris les phrases suivantes au passif (Präsens), puis ajoute à deux d'entre elles une subordonnée introduite par \all{weil} ou \all{damit}.
\begin{enumerate}
\item Die Schüler reinigen den Schulhof jeden Freitag.
\item Die Gemeinde baut einen neuen Markt.
\item Die Jugendlichen organisieren eine Baumpflanzaktion.
\item Der Verein hilft den Kindern ohne Eltern.
\end{enumerate}
\end{exercice}

\begin{exercice}[Texte à trous : les devoirs du citoyen]
Complète le texte suivant avec les mots de la liste : \all{wählen, Gesetze, Steuern, Müll, Verantwortung, Mitmenschen}.

\begin{textebox}[Ein guter Bürger]
Ein guter Bürger respektiert die \all{\ldots} seines Landes und zahlt regelmäßig seine \all{\ldots}. Er geht zu den Wahlen, denn \all{\ldots} ist ein wichtiges Recht. Er wirft keinen \all{\ldots} auf die Straße und hilft seinen \all{\ldots}, wenn sie in Not sind. Kurz gesagt: Er übernimmt \all{\ldots} für seine Gemeinde.
\end{textebox}
\end{exercice}

\begin{exercice}[Appariements : genres de textes]
Associe chaque extrait (1--4) au genre de texte qui lui correspond (a--d) : a) discours, b) lettre ouverte, c) interview, d) article de presse.
\begin{enumerate}
\item \all{„Sehr geehrter Herr Bürgermeister, ich schreibe Ihnen, weil die Situation in unserem Viertel unerträglich geworden ist \ldots“}
\item \all{Journalistin: Warum engagieren Sie sich für dieses Projekt? -- Schüler: Weil ich glaube, dass jeder etwas beitragen kann.}
\item \all{Liebe Mitbürgerinnen und Mitbürger! Heute stehen wir vor einer wichtigen Entscheidung für die Zukunft unserer Stadt.}
\item \all{Yaoundé. Seit einer Woche sammeln Jugendliche jeden Samstag Abfälle im Stadtzentrum. Die Aktion findet immer mehr Anhänger.}
\end{enumerate}
\end{exercice}

\courssub{Je me prépare au Bac}

\begin{exercice}[Compréhension de texte et questions de langue]
Lis le texte suivant, puis réponds aux questions.

\begin{textebox}[Jugendliche engagieren sich in Douala]
In Douala haben Schüler einer weiterführenden Schule ein bemerkenswertes Projekt gestartet: Jeden Samstagmorgen treffen sich etwa dreißig Jugendliche, um die Straßen ihres Viertels zu reinigen und die Bewohner über den Umweltschutz zu informieren. „Wir wollen nicht nur warten, bis der Staat alles macht“, erklärt Amina, die Leiterin der Gruppe. „Citoyenneté bedeutet für uns, selbst die Verantwortung zu übernehmen.“ Die Aktion, die vor sechs Monaten begann, wird inzwischen von der Gemeinde unterstützt, die Handschuhe und Säcke zur Verfügung stellt. Viele Nachbarn sind stolz auf die Jugendlichen und einige Erwachsene haben sogar angefangen, mitzumachen. Die Schüler hoffen, dass ihr Beispiel auch andere Schulen in Kamerun motivieren wird.
\end{textebox}

\textbf{A. Compréhension} — Réponds en français par des phrases complètes.
\begin{enumerate}
\item Que font les élèves chaque samedi et pourquoi ?
\item Que signifie la citoyenneté pour Amina ? Cite le texte en allemand.
\item Comment la commune soutient-elle le projet ?
\item Quel est l'espoir des élèves ?
\end{enumerate}

\textbf{B. Langue}
\begin{enumerate}
\item Trouve dans le texte : un verbe réfléchi, un verbe au passif, une subordonnée avec \all{dass}.
\item Conjugue \all{sich treffen} au Präsens (ich, du, wir) et au Perfekt (sie, pluriel).
\item Mets la phrase suivante au passif : \all{Die Gemeinde unterstützt die Aktion.}
\end{enumerate}
\end{exercice}

\begin{exercice}[Version (type Bac)]
Traduis en français le texte suivant.

\begin{textebox}[Wählen ab 16?]
In Deutschland diskutiert man seit Jahren darüber, ob Jugendliche schon mit 16 Jahren wählen dürfen sollen. Bei den Europawahlen ist das inzwischen möglich. Die Befürworter sagen: Junge Menschen sind gut informiert und Entscheidungen von heute betreffen vor allem ihre Zukunft, zum Beispiel beim Klimaschutz oder bei der Bildung. Die Gegner meinen dagegen, dass viele Sechzehnjährige noch nicht genug Erfahrung haben und sich zu wenig für Politik interessieren. Eine Schülerin aus Hamburg erklärt: „Wir müssen die Folgen der Politik am längsten tragen. Deshalb sollten wir auch mitentscheiden dürfen.“ Ob das Wahlalter in Deutschland bald gesenkt wird, ist noch offen.
\end{textebox}
\end{exercice}

\begin{exercice}[Thème et production écrite (type Bac)]
\textbf{A. Thème.} Traduis en allemand les cinq phrases suivantes.
\begin{enumerate}
\item Les jeunes devraient s'engager davantage parce que l'avenir du pays dépend d'eux.
\item Si tous les citoyens respectaient les lois, la vie serait plus facile.
\item Hier, nous avons aidé les vieux habitants de notre quartier.
\item Il faut que chaque élève comprenne l'importance du vote.
\item Le maire a promis qu'il construirait un centre pour les jeunes.
\end{enumerate}

\textbf{B. Production écrite.} Choisis \textbf{un} des deux sujets suivants.

\textbf{Sujet 1 : Lettre ouverte.} Rédige en allemand une lettre ouverte (12 à 15 lignes) au maire de ta commune pour lui proposer un projet citoyen (propreté du quartier, aide aux élèves défavorisés, journée de l'arbre). Respecte la forme de la lettre : lieu et date, formule d'appel, introduction, présentation du projet, arguments, formule de politesse.

\textbf{Sujet 2 : Interview.} Rédige en allemand le dialogue (au moins 8 répliques) entre un journaliste allemand et un élève camerounais engagé dans un projet social. Le journaliste pose des questions sur le projet, les motivations, les difficultés et les projets d'avenir de l'élève.
\end{exercice}

\newpage

\coursec{Corrigés des exercices}

\courssub{Je vérifie mes connaissances}

\begin{corrige}[Exercice 1 — QCM : vocabulaire de la citoyenneté]
\begin{enumerate}
\item \textbf{1.} \all{die Staatsbürgerschaft} = la nationalité / la citoyenneté. (Composé de \all{der Staat} + \all{der Bürger} + suffixe \all{-schaft}.)
\item \textbf{1.} \all{wählen} = voter / élire. Attention : \all{wählen} signifie aussi « composer » un numéro de téléphone (\all{eine Nummer wählen}).
\item \textbf{1.} \all{das Recht} (le droit) s'oppose juridiquement à \all{die Pflicht} (le devoir). \all{die Freiwilligkeit} désigne le caractère volontaire d'une action, ce n'est pas un contraire.
\item \textbf{2.} \all{sich engagieren} = s'engager, s'investir (bénévolement). Faux ami : « se marier » se dit \all{heiraten}.
\item \textbf{2.} \all{der Bürgermeister} = le maire. Le citoyen = \all{der Bürger / der Staatsbürger} ; le député = \all{der Abgeordnete}.
\end{enumerate}
\end{corrige}

\begin{corrige}[Exercice 2 — Vrai ou faux]
\begin{enumerate}
\item \textbf{Vrai.} C'est la règle de la place du verbe (\cle{V2}) : dans la phrase principale allemande, le verbe conjugué occupe toujours la deuxième position, précédé d'un seul élément (sujet, complément de temps, complément d'objet\ldots). Ex. : \all{Heute \textbf{gehen} wir wählen.}
\item \textbf{Faux.} Dans une subordonnée introduite par \all{weil}, le verbe conjugué est rejeté \textbf{en dernière position}. Ex. : \all{Ich bleibe zu Hause, weil es regnet.} Ne jamais écrire \all{weil es regnet} avec le verbe en deuxième position.
\item \textbf{Faux.} Le passif d'action (\cle{Vorgangspassiv}) se forme avec \all{werden} + Partizip II : \all{Der Brief wird geschrieben.} L'auxiliaire \all{haben} + Partizip II sert au Perfekt actif. (Le passif d'état, \cle{Zustandspassiv}, utilise \all{sein} + Partizip II : \all{Die Tür ist geöffnet.})
\item \textbf{Vrai.} \all{damit} (afin que) introduit une subordonnée de but et renvoie le verbe conjugué à la fin : \all{Ich lerne viel, damit ich die Prüfung bestehe.}
\item \textbf{Vrai.} Tous les noms allemands prennent la majuscule, quelle que soit leur place dans la phrase : \all{der Bürger, die Wahl, das Recht}.
\end{enumerate}
\end{corrige}

\begin{corrige}[Exercice 3 — Lexique : trouve le mot allemand]
\begin{enumerate}
\item le vote, l'élection : \all{die Wahl, die Wahlen}
\item le bénévolat, le volontariat : \all{das Ehrenamt, die Ehrenämter} (ou \all{die Freiwilligenarbeit}, sans pluriel usuel)
\item la responsabilité : \all{die Verantwortung} (pluriel rare : \all{die Verantwortungen})
\item la société : \all{die Gesellschaft, die Gesellschaften}
\item les droits de l'homme : \all{die Menschenrechte} (n'existe qu'au pluriel)
\item la commune, la municipalité : \all{die Gemeinde, die Gemeinden}
\end{enumerate}
\textbf{Point de vigilance :} apprends toujours chaque nom avec son article et son pluriel ; le genre est imprévisible en allemand.
\end{corrige}

\begin{corrige}[Exercice 4 — Conjugaison]
\begin{enumerate}
\item \all{Die Jugendlichen \textbf{engagieren sich} in einem sozialen Projekt.} (Verbe pronominal : le pronom réfléchi \all{sich} suit immédiatement le verbe conjugué.)
\item \all{Gestern \textbf{hatten} wir eine Sitzung im Gemeindehaus.} (Präteritum de \all{haben} : ich hatte, du hattest, er/sie/es hatte, wir hatten, ihr hattet, sie/Sie hatten.)
\item \all{Der Bürgermeister \textbf{hat versprochen} mehr Mittel für die Schule.} (Perfekt : auxiliaire \all{haben} + Partizip II \all{versprochen} ; verbe fort mais participe régulier sans changement de radical : versprechen -- versprach -- hat versprochen.)
\item \all{Wenn ich mehr Zeit hätte, \textbf{würde ich helfen} den alten Menschen.} Forme correcte et plus élégante : \all{\ldots würde ich den alten Menschen helfen.} (Konjunktiv II analytique : \all{würde} + infinitif en fin de phrase ; la forme simple \all{hülfe} est soutenue et rare.) Attention : \all{helfen} + Datif : \all{den alten Menschen}.
\item \all{Nächstes Jahr \textbf{werden} viele junge Leute zum ersten Mal \textbf{wählen}.} (Futur I : \all{werden} conjugué en position 2 + infinitif \all{wählen} en fin de phrase.)
\end{enumerate}
\end{corrige}

\courssub{Je m'entraîne}

\begin{corrige}[Exercice 5 — Thème : phrases à traduire]
\begin{enumerate}
\item \all{Die Jugendlichen sollten sich mehr engagieren, weil die Zukunft des Landes von ihnen abhängt.}\\
Points de vigilance : \all{sich engagieren} (pronom réfléchi après le verbe) ; \all{abhängen von} + Datif (\all{von ihnen}) ; dans la subordonnée \all{weil}, le verbe \all{abhängt} va à la fin ; génitif \all{des Landes}.
\item \all{Jeder Bürger hat Rechte, aber auch Pflichten.}\\
\all{jeder} + nominatif singulier masculin (\all{jeder Bürger}) ; \all{Rechte} et \all{Pflichten} à l'accusatif pluriel (sans article, donc terminaison \all{-e} : \all{Rechte}, \all{Pflichten}).
\item \all{Wenn ich Bürgermeister wäre, würde ich eine Bibliothek für die Schüler bauen.}\\
Hypothèse irréelle : Konjunktiv II \all{wäre} (et non \all{würde sein}, lourd) ; pas d'article devant la fonction après \all{sein} : \all{Bürgermeister} ; \all{würde} + infinitif \all{bauen} en fin de phrase principale.
\item \all{Wir haben Bäume gepflanzt, um die Umwelt zu schützen.}\\
Perfekt : \all{haben} + \all{gepflanzt} (verbe faible) ; but exprimé par \all{um \ldots\ zu} + infinitif (même sujet dans les deux propositions).
\item \all{Die Wahl ist ein wichtiges Recht in einer Demokratie.}\\
\all{die Wahl} est féminin ; \all{das Recht} est neutre, d'où \all{ein wichtiges Recht} (accusatif neutre après \all{sein} : en fait nominatif attribut) ; \all{in einer Demokratie} : datif féminin après \all{in} (localisation).
\end{enumerate}
\end{corrige}

\begin{corrige}[Exercice 6 — Transformation : actif $\rightarrow$ passif]
Rappel de méthode : le complément d'objet direct (accusatif) de la phrase active devient le sujet (nominatif) de la phrase passive ; le verbe se met sous la forme \all{werden} + Partizip II ; l'agent est introduit par \all{von} + Datif.
\begin{enumerate}
\item \all{Der Schulhof wird jeden Freitag von den Schülern gereinigt.}\\
Ajout possible : \all{Der Schulhof wird jeden Freitag von den Schülern gereinigt, damit die Schule sauber bleibt.}
\item \all{Ein neuer Markt wird von der Gemeinde gebaut.}\\
Ajout possible : \all{Ein neuer Markt wird von der Gemeinde gebaut, weil der alte Markt zu klein geworden ist.}
\item \all{Eine Baumpflanzaktion wird von den Jugendlichen organisiert.}
\item \all{Den Kindern ohne Eltern wird vom Verein geholfen.}\\
\textbf{Cas particulier :} \all{helfen} se construit avec le \textbf{datif} ; il n'y a donc pas d'accusatif à transformer en sujet. Le datif est conservé (\all{den Kindern}) et le passif est impersonnel : \all{Den Kindern wird geholfen} ou \all{Es wird den Kindern geholfen.}
\end{enumerate}
\end{corrige}

\begin{corrige}[Exercice 7 — Texte à trous]
\begin{textebox}[Ein guter Bürger — Lösung]
Ein guter Bürger respektiert die \all{Gesetze} seines Landes und zahlt regelmäßig seine \all{Steuern}. Er geht zu den Wahlen, denn \all{Wählen} ist ein wichtiges Recht. Er wirft keinen \all{Müll} auf die Straße und hilft seinen \all{Mitmenschen}, wenn sie in Not sind. Kurz gesagt: Er übernimmt \all{Verantwortung} für seine Gemeinde.
\end{textebox}
Justifications : \all{die Gesetze} (pluriel de \all{das Gesetz}, accord avec \all{die}) ; \all{seine Steuern} (pluriel de \all{die Steuer}) ; \all{das Wählen} (infinitif substantivé, neutre, majuscule : « voter » comme action) ; \all{keinen Müll} (accusatif masculin après \all{kein-}) ; \all{seinen Mitmenschen} (datif pluriel après \all{helfen}, d'où \all{-n}) ; \all{Verantwortung übernehmen} (expression figée : assumer la responsabilité).
\end{corrige}

\begin{corrige}[Exercice 8 — Appariements : genres de textes]
\begin{enumerate}
\item \textbf{b) Lettre ouverte} : formule d'appel formelle \all{Sehr geehrter Herr Bürgermeister}, verbe typique \all{ich schreibe Ihnen}, destinataire nommé.
\item \textbf{c) Interview} : alternance question/réponse entre deux locuteurs identifiés (\all{Journalistin} / \all{Schüler}), tiret de dialogue.
\item \textbf{a) Discours} : apostrophe au public \all{Liebe Mitbürgerinnen und Mitbürger!}, ton solennel, \all{wir} fédérateur, enjeu collectif.
\item \textbf{d) Article de presse} : ligne de lieu (\all{Yaoundé.}), faits rapportés à la troisième personne, style informatif et objectif.
\end{enumerate}
\end{corrige}

\courssub{Je me prépare au Bac}

\begin{corrige}[Exercice 9 — Compréhension de texte et questions de langue]
\textbf{A. Compréhension}
\begin{enumerate}
\item Chaque samedi matin, une trentaine d'élèves se retrouvent pour nettoyer les rues de leur quartier et informer les habitants sur la protection de l'environnement. Ils le font parce qu'ils ne veulent pas seulement attendre que l'État fasse tout.
\item Pour Amina, le sens civique (\all{Bürgersinn}) signifie assumer soi-même la responsabilité. Citation : \all{„Bürgersinn bedeutet für uns, selbst die Verantwortung zu übernehmen.“}
\item La commune soutient le projet en fournissant des gants et des sacs : \all{die \ldots\ Handschuhe und Säcke zur Verfügung stellt}.
\item Les élèves espèrent que leur exemple motivera aussi d'autres écoles du Cameroun.
\end{enumerate}
\textbf{B. Langue}
\begin{enumerate}
\item Verbe pronominal : \all{sich treffen} (\all{treffen sich etwa dreißig Jugendliche}) ; verbe au passif : \all{wird \ldots\ von der Gemeinde unterstützt} ; subordonnée en \all{dass} : \all{dass ihr Beispiel auch andere Schulen in Kamerun motivieren wird}.
\item \all{sich treffen} au Präsens : \all{ich treffe mich, du triffst dich, wir treffen uns} (verbe fort à changement de vocalisme e $\rightarrow$ i). Au Perfekt (3\textsuperscript{e} pers. plur.) : \all{sie haben sich getroffen} (auxiliaire \all{haben} : verbe sans idée de déplacement).
\item \all{Die Aktion wird von der Gemeinde unterstützt.} (L'accusatif \all{die Aktion} devient sujet nominatif ; \all{wird} + Partizip II \all{unterstützt} ; agent introduit par \all{von} + datif.)
\end{enumerate}
\textbf{Barème indicatif (sur 20) :} A : 4 $\times$ 2 = 8 pts ; B1 : 3 $\times$ 2 = 6 pts ; B2 : 4 pts ; B3 : 2 pts.
\end{corrige}

\begin{corrige}[Exercice 10 — Version (type Bac)]
\textbf{Traduction proposée :}

En Allemagne, on discute depuis des années de la question de savoir si les jeunes devraient déjà avoir le droit de voter à seize ans. Aux élections européennes, c'est désormais possible. Les partisans disent : les jeunes sont bien informés et les décisions d'aujourd'hui concernent avant tout leur avenir, par exemple en matière de protection du climat ou d'éducation. Les opposants estiment au contraire que beaucoup de jeunes de seize ans n'ont pas encore assez d'expérience et s'intéressent trop peu à la politique. Une lycéenne de Hambourg explique : « Nous devons supporter les conséquences de la politique le plus longtemps. C'est pourquoi nous devrions aussi pouvoir participer aux décisions. » On ne sait pas encore si l'âge du vote sera bientôt abaissé en Allemagne.

\textbf{Points de vigilance :}
\begin{itemize}
\item \all{die Befürworter} = les partisans ; \all{die Gegner} = les opposants (ne pas traduire par « les ennemis »).
\item \all{betreffen} = concerner (faux ami : ce n'est pas « trahir » ni « affecter » au sens émotionnel).
\item \all{die Folgen tragen} = supporter les conséquences (expression figée).
\item \all{mitentscheiden dürfen} = pouvoir participer à la décision, co-décider.
\item \all{das Wahlalter senken} = abaisser l'âge du vote ; \all{ist noch offen} = reste une question ouverte.
\end{itemize}
\textbf{Barème indicatif (sur 20) :} 10 segments de sens $\times$ 2 pts ; tolérance pour les maladresses de style si le sens est respecté ; erreur de sens : $-1$ pt ; contresens : $-2$ pts.
\end{corrige}

\begin{corrige}[Exercice 11 A — Thème (type Bac)]
\begin{enumerate}
\item \all{Die Jugendlichen sollten sich mehr engagieren, weil die Zukunft des Landes von ihnen abhängt.}\\
(\all{sollen} au Konjunktiv II \all{sollten} pour « devraient » ; verbe en fin de subordonnée \all{weil}.)
\item \all{Wenn alle Bürger die Gesetze respektieren würden, wäre das Leben einfacher.}\\
(Konjunktiv II dans les deux propositions : \all{würden respektieren} + \all{wäre} ; forme simple \all{respektierten} également correcte.)
\item \all{Gestern haben wir den alten Bewohnern unseres Viertels geholfen.}\\
(\all{helfen} + datif pluriel : \all{den alten Bewohnern} ; génitif : \all{unseres Viertels} ; Perfekt : \all{haben \ldots\ geholfen}.)
\item \all{Jeder Schüler muss die Bedeutung der Wahl verstehen.} ou, avec subordonnée : \all{Es ist wichtig, dass jeder Schüler die Bedeutung der Wahl versteht.}\\
(« Il faut que » + subjonctif français se rend le plus naturellement par \all{müssen} + infinitif ou \all{es ist wichtig, dass} + indicatif.)
\item \all{Der Bürgermeister hat versprochen, dass er ein Zentrum für die Jugend bauen wird.}\\
(Discours rapporté : Futur I \all{bauen wird} ; le Konjunktiv I \all{bauen werde} est la forme soutenue du discours indirect : \all{\ldots, dass er ein Zentrum für die Jugend bauen werde.})
\end{enumerate}
\textbf{Barème indicatif (sur 10) :} 5 phrases $\times$ 2 pts (1 pt pour le sens et le lexique, 1 pt pour la morphosyntaxe : place du verbe, cas, temps).
\end{corrige}

\begin{corrige}[Exercice 11 B — Production écrite : propositions de correction]
\textbf{Sujet 1 : Lettre ouverte (modèle)}
\begin{textebox}[Offener Brief an den Bürgermeister]
Yaoundé, den 15. Oktober 2024

Sehr geehrter Herr Bürgermeister,

ich schreibe Ihnen als Schülerin eines Gymnasiums in Yaoundé, um Ihnen ein Projekt vorzuschlagen, das unser Viertel sauberer und grüner machen soll: eine monatliche „Aktion Saubere Straße“.

Jeden ersten Samstag im Monat möchten sich Schüler, Eltern und Anwohner treffen, um Müll zu sammeln, Bäume zu pflanzen und die Grünflächen zu pflegen. Die Schule stellt Handschuhe und Müllsäcke zur Verfügung; wir bitten die Gemeinde nur um die Abholung des gesammelten Mülls und um Gießkannen für die jungen Bäume.

Dieses Projekt stärkt den Bürgersinn der Jugend, verschönert das Stadtbild und kostet die Gemeinde wenig. Wenn alle mitmachen, wird unser Viertel ein Vorbild für ganz Yaoundé sein.

Ich hoffe auf Ihre positive Antwort und stehe für ein Gespräch gern zur Verfügung.

Mit freundlichen Grüßen

Aminata Ngo Bell, Klasse Tle A
\end{textebox}
\textbf{Points de vigilance (lettre) :} lieu et date avec \all{den} + date (\all{Yaoundé, den 15. Oktober}) ; formule d'appel suivie d'une virgule et d'une minuscule à la ligne suivante ; \all{Ihnen} avec majuscule (vouvoiement) ; formule finale \all{Mit freundlichen Grüßen} sans virgule ; infinitifs en \all{um \ldots\ zu} pour exprimer le but.

\textbf{Sujet 2 : Interview (modèle)}
\begin{textebox}[Interview mit einem engagierten Schüler]
\textbf{Journalist:} Guten Tag, Paul! Sie engagieren sich in Douala in einem sozialen Projekt. Was machen Sie genau?

\textbf{Paul:} Guten Tag! Mit einigen Freunden gebe ich Nachhilfe für Grundschulkinder, deren Eltern sich keinen Privatlehrer leisten können.

\textbf{Journalist:} Wie sind Sie auf diese Idee gekommen?

\textbf{Paul:} Ich habe gesehen, dass viele Kinder nach der Schule auf der Straße spielen, statt ihre Hausaufgaben zu machen. Da wollte ich helfen.

\textbf{Journalist:} Gab es am Anfang Schwierigkeiten?

\textbf{Paul:} Ja, natürlich. Uns fehlten Räume und Bücher. Jetzt dürfen wir zum Glück einen Raum in der Gemeindebibliothek benutzen, aber wir suchen noch mehr Freiwillige.

\textbf{Journalist:} Was motiviert Sie, trotzdem weiterzumachen?

\textbf{Paul:} Wenn ein Kind sagt: „Jetzt verstehe ich Mathematik!“, dann ist das der schönste Lohn. Für mich ist das echter Bürgersinn.

\textbf{Journalist:} Haben Sie Pläne für die Zukunft?

\textbf{Paul:} Ja, wir möchten ein „Lerncafé“ eröffnen, auch für ältere Jugendliche. Und später möchte ich Sozialarbeiter werden.

\textbf{Journalist:} Vielen Dank für das Gespräch und viel Erfolg!

\textbf{Paul:} Ich danke Ihnen!
\end{textebox}
\textbf{Points de vigilance (interview) :} alternance claire des locuteurs en gras ; questions variées (avec \all{W-Fragen} : \all{Was, Wie, Warum}) ; inversion après les questions ; registre de politesse (\all{Sie}) ; vocabulaire du chapitre réinvesti (\all{sich engagieren, die Nachhilfe, der Freiwillige, der Bürgersinn}).

\textbf{Barème indicatif (production, sur 10) :} respect du type de texte et de la consigne : 2 pts ; richesse et exactitude du lexique : 2 pts ; correction grammaticale (place du verbe, cas, temps) : 3 pts ; cohérence et articulation du texte : 2 pts ; orthographe et majuscules : 1 pt.
\end{corrige}

\newpage

\coursec{Fiche bilan}

\begin{popfigure}
\begin{center}
\begin{tikzpicture}[mindmap, grow cyclic, every node/.style={concept}, root concept/.append style={concept color=popblue, font=\large\bfseries, text=white}, level 1/.append style={level distance=3.4cm, sibling angle=60}, level 1 concept/.append style={font=\small\bfseries, text=white}]
\node[root concept]{AL30 Citoyenneté : produire et traduire}
  child[concept color=poporange]{ node{Lexique : die Staatsbürgerschaft, wählen, die Pflicht} }
  child[concept color=popgreen]{ node{Grammaire : verbe en 2\up{e} position, cas, Konjunktiv II} }
  child[concept color=poppurple]{ node{Méthode : commentaire de texte en 5 étapes} }
  child[concept color=poppink]{ node{Production : Rede, offener Brief, Interview, Artikel} }
  child[concept color=popgold]{ node{Thème : français $\rightarrow$ allemand} }
  child[concept color=popdark]{ node{Version : allemand $\rightarrow$ français} };
\end{tikzpicture}
\end{center}
\legende{Carte mentale de la leçon AL30 : les six compétences à mobiliser pour produire et traduire des textes sur la citoyenneté.}
\end{popfigure}

\begin{bilan}
\textbf{1. Lexique clé de la citoyenneté (à connaître par cœur)}
\begin{itemize}
\item \all{die Staatsbürgerschaft, -en} : la nationalité, la citoyenneté ; \all{der Bürger, - / die Bürgerin, -nen} : le citoyen / la citoyenne
\item \all{das Recht, -e} : le droit ; \all{die Pflicht, -en} : le devoir ; \all{die Freiheit, -en} : la liberté ; \all{die Verantwortung, -en} : la responsabilité
\item \all{wählen (wählte, hat gewählt)} : élire, voter ; \all{die Wahl, -en} : l'élection ; \all{wählen gehen} : aller voter ; \all{der Wähler, - / die Wählerin, -nen} : l'électeur / l'électrice
\item \all{die Demokratie, -n} : la démocratie ; \all{die Menschenrechte (Pl.)} : les droits de l'homme ; \all{die Gleichheit} : l'égalité ; \all{die Gerechtigkeit} : la justice
\item \all{sich engagieren (für + Akk.)} : s'engager (pour) ; \all{der/die Freiwillige, -n} : le/la bénévole (déclinaison comme un adjectif) ; \all{die Jugend} : la jeunesse
\item \all{die Meinungsfreiheit} : la liberté d'expression ; \all{die Presse, -n} : la presse ; \all{der Artikel, -} : l'article
\item \all{respektieren} : respecter ; \all{tolerant sein} : être tolérant ; \all{der Frieden} : la paix ; \all{die Solidarität} : la solidarité
\item \all{der Staat, -en} : l'État ; \all{die Regierung, -en} : le gouvernement ; \all{der Bürgermeister, - / die Bürgermeisterin, -nen} : le maire
\item \all{abstimmen (über + Akk.)} : voter (sur une question) ; \all{teilnehmen (an + Dat.)} : participer (à) ; \all{die Verfassung, -en} : la constitution
\end{itemize}

\textbf{2. Grammaire : les règles à connaître par cœur}
\begin{center}
\begin{tabularx}{\linewidth}{|l|X|X|}
\hline
\rowcolor{popblueL} \textbf{Point} & \textbf{Règle} & \textbf{Exemple} \\
\hline
Place du verbe & Le verbe conjugué est toujours en 2\up{e} position dans la phrase déclarative. & \all{Heute gehen viele Jugendliche wählen.} \\
\hline
Inversion & Si un complément ouvre la phrase, le sujet suit immédiatement le verbe. & \all{In Kamerun engagieren sich viele junge Bürger.} \\
\hline
Subordonnée & Verbe conjugué en dernière position après \all{dass, weil, wenn, obwohl, damit}. & \all{Ich glaube, dass die Jugend wählen muss.} \\
\hline
Verbes à particule & La particule séparable se détache et va en fin de proposition principale. & \all{Der Bürgermeister nimmt an der Diskussion teil.} \\
\hline
Cas & Akkusativ après \all{für, durch, gegen, ohne, um} ; Dativ après \all{mit, von, zu, bei, nach, aus}. & \all{für die Demokratie ; mit dem Bürgermeister} \\
\hline
Konjunktiv II & Hypothèse, souhait, conseil : \all{würde + Infinitiv}, \all{hätte, wäre, könnte, sollte}. & \all{Wenn ich Bürgermeister wäre, würde ich mehr Schulen bauen.} \\
\hline
Passif & \all{werden + Partizip II} ; au Perfekt : \all{ist ... worden}. & \all{Der Präsident wird gewählt. / Er ist gewählt worden.} \\
\hline
Discours indirect & Konjunktiv I (ou Konjunktiv II de remplacement si la forme est identique au présent). & \all{Er sagte, er habe gewählt. / Sie sagten, sie hätten gewählt.} \\
\hline
\end{tabularx}
\end{center}

\textbf{3. Méthodes express}
\begin{itemize}
\item \textbf{Commentaire de texte} : 1) lire deux fois et souligner les mots-clés ; 2) identifier le type de texte, la source et le thème ; 3) dégager la thèse de l'auteur ; 4) analyser les arguments et les exemples ; 5) donner son avis argumenté (\all{Meiner Meinung nach ...}, \all{Ich bin der Ansicht, dass ...}).
\item \textbf{Thème (français $\rightarrow$ allemand)} : repérer le temps, les cas et l'ordre des mots ; traduire les idées, pas les mots ; vérifier la majuscule des noms et le verbe en 2\up{e} position.
\item \textbf{Version (allemand $\rightarrow$ français)} : identifier sujet, verbe, compléments ; découper les subordonnées ; rendre en français naturel, jamais mot à mot.
\item \textbf{Production} : adapter le registre au genre demandé : \all{die Rede} (formule d'appel : \all{Liebe Mitbürgerinnen und Mitbürger!}), \all{der offene Brief} (destinataire, lieu et date, salutations), \all{das Interview} (questions avec inversion : \all{Engagieren Sie sich ...?}), \all{der Artikel} (titre, chapeau, paragraphes).
\end{itemize}
\end{bilan}

\begin{attention}[Erreurs fréquentes des francophones]
\begin{itemize}
\item Ne dites pas \all{Ich gehe wählen heute} : le complément de temps ne se place jamais après l'infinitif. Correct : \all{Ich gehe heute wählen.}
\item \all{wählen} signifie « élire / voter », pas « choisir » au sens large : pour « choisir un métier », employez \all{wählen} uniquement s'il s'agit d'un choix entre options, sinon \all{sich entscheiden für + Akk.}
\item Attention au verbe à particule : \all{teilnehmen} se sépare : \all{Viele Jugendliche nehmen an der Wahl teil.} (et non \all{... teilnehmen an der Wahl}).
\item Faux ami : \all{die Wahl} = l'élection, mais \all{die Wahl treffen} = faire un choix ; « le choix » n'est pas toujours \all{die Wahl} dans le sens politique.
\item N'oubliez pas la majuscule à tous les noms : \all{das Recht, die Pflicht, der Bürger}.
\item Après \all{für}, toujours l'accusatif : \all{für die Demokratie}, jamais \all{für der Demokratie}.
\end{itemize}
\end{attention}

\begin{exoresolu}[Thème guidé : une phrase à traduire]
Énoncé : traduisez en allemand : « Si j'étais citoyen engagé, je participerais à toutes les élections. »
\tcblower
Solution détaillée :
\begin{enumerate}
\item La phrase exprime une hypothèse irréelle : il faut le Konjunktiv II.
\item « Si j'étais » : \all{Wenn ich ... wäre} ; « citoyen engagé » : \all{ein engagierter Bürger} (adjectif épithète décliné : après \all{ein}, nominatif masculin \all{engagierter}).
\item « je participerais » : \all{würde ich teilnehmen} ; attention, \all{teilnehmen an + Dativ} : \all{an allen Wahlen}.
\item Dans la subordonnée en \all{wenn}, le verbe conjugué va à la fin ; dans la principale qui suit, le verbe ouvre la proposition (inversion).
\end{enumerate}
Traduction correcte : \all{Wenn ich ein engagierter Bürger wäre, würde ich an allen Wahlen teilnehmen.}
\end{exoresolu}

\begin{aretenir}[Checklist avant le Bac]
\begin{itemize}
\item[$\square$] Je connais le lexique de la citoyenneté avec articles et pluriels (\all{die Wahl, -en}).
\item[$\square$] Je place toujours le verbe conjugué en 2\up{e} position dans la phrase principale.
\item[$\square$] Je mets le verbe conjugué à la fin des subordonnées introduites par \all{dass, weil, wenn, obwohl}.
\item[$\square$] Je sépare correctement les verbes à particule : \all{Er nimmt an der Wahl teil.}
\item[$\square$] Je mets une majuscule à tous les noms allemands.
\item[$\square$] Je choisis le bon cas après les prépositions (\all{für + Akk.}, \all{mit + Dat.}).
\item[$\square$] J'emploie le Konjunktiv II pour exprimer un souhait ou une hypothèse : \all{Ich würde wählen gehen.}
\item[$\square$] Je sais structurer un discours, une lettre ouverte, une interview et un article.
\item[$\square$] En thème, je traduis les idées et je vérifie l'ordre des mots.
\item[$\square$] En version, je rends un français naturel et fidèle au sens.
\item[$\square$] Je relis ma copie : orthographe, accords, ponctuation, guillemets allemands „ ... “.
\item[$\square$] Je gère mon temps : lecture, brouillon, rédaction, relecture.
\end{itemize}
\end{aretenir}

\findecours

\end{document}
